% ============================================================
% Notas de Biomate — LAYOUT MAESTO
% Índice basado en notes_drafts_biomate_1.pdf (29 págs., transcrito
% en /tmp/biomath_ocr/draft_layout_transcription.md).
% Contenido disperso en "Notas Biomate 1.pdf" (español) y
% "Notes_ Non Linear Dynamics.pdf" (inglés) mapeado por capítulo
% en los comentarios de cada archivo de chapters/.
% Estilo: Variante A (Tong minimal), heredada de demo-a-tong.tex.
% ============================================================
\documentclass[11pt,a4paper]{article}
\usepackage[utf8]{inputenc}
\usepackage[spanish,es-noquoting]{babel}
\usepackage{amsmath,amssymb,amsthm}
\usepackage{graphicx}
\graphicspath{{figures/}}
\usepackage{tikz}
\usetikzlibrary{arrows.meta,decorations.markings,calc,shapes,positioning}
\usepackage{pgfplots}
\pgfplotsset{compat=newest}

\usepackage{minted}
\usemintedstyle{friendly_grayscale}
\usepackage[backend=biber, style=numeric, sorting=none]{biblatex}
\addbibresource{refs.bib}
\usepackage[hidelinks]{hyperref}
\usepackage[hang,flushmargin]{footmisc}
\usepackage{cleveref}

\newtheorem{theorem}{Teorema}[section]
\newtheorem{definition}{Definición}[section]
\newtheorem{example}{Ejemplo}[section]
\newcommand{\aside}[1]{\footnote{#1}}

\title{Notas de Matemáticas Biológicas\\[4pt]
       \large Sistemas Dinámicos No Lineales}
\author{Ulises Rayón}
\date{}

\begin{document}
\maketitle
\tableofcontents

% ---- Parte I — Preliminares (draft págs. 1--7) ----
% ============================================================
% FUENTES: draft p.1--2 (transcripción ## Page 1--2).
% ============================================================
\section{Sobre la evaluación y el temario}

\subsection{Reglas del juego y esquema de evaluación}

Bienvenidos al curso de Matemáticas Biológicas. Antes de adentrarnos formalmente en las ecuaciones diferenciales y en la dinámica no lineal, dejemos claras las reglas de evaluación y la dinámica de trabajo:

\begin{itemize}
    \item \textbf{Mini Proyectos:} Tareas-proyecto extensas de desarrollo teórico y numérico. Usualmente se asignan con un plazo de 15 días para su entrega y resolución cuidadosa.
    \item \textbf{Mini Labs:} Prácticas computacionales opcionales, pero sumamente útiles para perderle el miedo a la programación y a la simulación numérica de sistemas dinámicos.
    \item \textbf{Proyecto Final:} Se desarrollará durante el último cuarto del semestre (aproximadamente en el último mes) y consiste en la modelación y análisis cualitativo completo de un fenómeno biológico de interés.
\end{itemize}

\paragraph{Política sobre citas, copiado y plagio.}
Sé bien cómo funcionan estas dinámicas y sé que van a consultar y discutir material entre ustedes. Los problemas del curso han sido diseñados específicamente para estas notas, por lo que la gran mayoría no tiene soluciones directas flotando en internet. Si colaboran o toman una idea de un compañero o de un libro, no me molesta en absoluto, \emph{siempre y cuando incluyan la referencia explícita y den el crédito correspondiente}. Lo que no tolero es el plagio deshonesto (presentar ideas ajenas como propias sin citar).

\aside{Una advertencia fundamental: lo que bajo ninguna circunstancia se aceptará es que entreguen fotos o capturas de pantalla del trabajo de otra persona. Al menos, usar las propias manos para transcribir, escribir en \LaTeX{} o teclear el código garantiza que procesen e incorporen los conceptos.}

\subsection{Filosofía del curso y precedentes necesarios}

Los cursos de modelación biológica y sistemas dinámicos son demandantes porque requieren conectar herramientas de diversas ramas de las matemáticas:
\begin{enumerate}
    \item Ecuaciones diferenciales ordinarias (EDOs).
    \item Álgebra lineal (valores propios, transformaciones lineales).
    \item Cálculo diferencial e integral (esencialmente saber graficar curvas, raíces e intersecciones).
    \item Pensamiento geométrico y cualitativo.
    \item Cierto sentido común o ``conocimiento biológico''\aside{Sea lo que sea que eso signifique en la práctica transdisciplinaria.} para interpretar variables y parámetros con significado biológico concreto.
\end{enumerate}

Si han cursado materias de otras disciplinas o tienen interés en la transdisciplina, este marco les resultará natural y divertido. La modelación matemática consiste en aprender a jugar con las herramientas analíticas para develar la estructura oculta de los fenómenos naturales.

\paragraph{Un enfoque geométrico, no axiomático-formalista.}
Este curso no seguirá la estructura clásica de \emph{Definición $\to$ Teorema $\to$ Demostración $\to$ Corolario}. En lugar de perdernos en tecnicismos analíticos que a menudo ocultan la física del problema, nos apoyaremos en intuición cualitativa, argumentos geométricos, retratos de fase y muchos diagramas.

\aside{No se dejen engañar por la aparente sencillez visual de los primeros temas: si descuidan los detalles y las sutilezas geométricas, el curso puede deparar sorpresas mayúsculas al avanzar hacia bifurcaciones y caos.}

Para construir estos cimientos, comenzaremos repasando la evolución histórica que llevó a la creación de la dinámica no lineal moderna.


% ============================================================
% FUENTES: draft p.3--7.
% ============================================================
\section{Un brevísimo resumen histórico de los sistemas dinámicos no lineales}

La historia de los sistemas dinámicos es, en gran medida, la historia del enfrentamiento entre la búsqueda de fórmulas analíticas exactas y la inevitable complejidad del mundo natural.

\subsection{Del problema de los dos cuerpos a Poincaré (1666--fines de 1800)}

En 1666, un joven Isaac Newton sentó las bases del cálculo diferencial e integral y formuló las leyes del movimiento y de la gravitación universal. Con estas herramientas resolvió de manera exacta el \textbf{problema de los dos cuerpos}: dos masas puntuales que interactúan bajo la atracción gravitacional mutua describen órbitas cónicas (elipses, parábolas o hipérbolas), justificando matemáticamente las leyes empíricas de Kepler.

El éxito fue tan rotundo que inauguró la era del mecanicismo determinista. Sin embargo, cuando Newton intentó extender el análisis al \textbf{problema de tres cuerpos} ($N=3$), en particular para calcular con precisión la órbita de la Luna bajo la influencia combinada de la Tierra y el Sol, la dificultad resultó infranqueable. En una célebre carta a Edmond Halley, Newton confesó con amargura:
\begin{quote}
    \emph{``Ningún problema me ha hecho doler tanto la cabeza como el problema del Sol-Tierra-Luna.''}
\end{quote}

Durante los siguientes dos siglos, los matemáticos más brillantes de Europa —entre ellos Leonhard Euler, Joseph-Louis Lagrange y Carl Friedrich Gauss— intentaron infructuosamente encontrar soluciones analíticas cerradas en términos de funciones elementales o series convergentes para el problema de $N$ cuerpos.

A finales del siglo XIX, con motivo del concurso en honor al sexagésimo cumpleaños del rey Óscar II de Suecia (1889), el matemático francés Henri Poincaré abordó el problema restringido de tres cuerpos. Poincaré demostró rigurosamente que, en general, no existen integrales primeras algebraicas adicionales ni series uniformemente convergentes que resuelvan el problema.

\paragraph{El nacimiento del espacio de fases y el caos.}
En lugar de buscar fórmulas algebraicas exactas que no existen, Poincaré propuso un cambio de paradigma radical: el \textbf{análisis cualitativo y geométrico}. Introdujo la noción de \emph{espacio de fases} (o plano de fase en 2D), transformando el estudio de una ecuación diferencial en el estudio de las trayectorias de un flujo continuo y de sus intersecciones con superficies transversales (secciones de Poincaré).

En su análisis de las órbitas homoclínicas y heteroclínicas, Poincaré descubrió la existencia de trayectorias con un entrelazamiento infinitamente complejo (el hoy llamado entramado homoclínico). Había descubierto el \textbf{caos determinista}:
\begin{definition}[Caos determinista]
    Comportamiento aperiódico a largo plazo que ocurre dentro de un sistema dinámico determinista no lineal, el cual exhibe una sensibilidad extrema a las condiciones iniciales (pequeñas variaciones iniciales divergen exponencialmente en el tiempo).
\end{definition}

\aside{Una anécdota imperdible: Henri Poincaré era pésimo dibujando; se cuenta que en los exámenes de admisión a la École Polytechnique obtuvo un rotundo cero en la prueba de dibujo mecánico. Paradójicamente, su intuición era profundamente geométrica y visual. Como en sus memorias publicadas casi no incluía figuras sino extensas descripciones geométricas abstractas, y como la atención de la física a principios del siglo XX se volcó hacia la mecánica cuántica y la teoría de la relatividad, sus revolucionarias ideas sobre dinámica permanecieron en relativa sombra durante varias décadas.}

\subsection{Osciladores no lineales, ingeniería y el clima (1920s--1963)}

Entre las décadas de 1920 y 1950, la dinámica no lineal resurgió impulsada por las necesidades tecnológicas de la física y la ingeniería aplicada:
\begin{itemize}
    \item El desarrollo de la radiocomunicación y los tubos de vacío (tríodos y diodos) condujo al estudio de osciladores auto-sostenidos no lineales, como el célebre modelo de Balthasar van der Pol (1926).
    \item Los sistemas de radar durante la Segunda Guerra Mundial y el posterior diseño de láseres requirieron comprender la sincronización y los ciclos límite.
    \item El control de frecuencia en telecomunicaciones dio origen a los lazos de enganche de fase (\emph{Phase Locked Loops}, PLL), cuyo funcionamiento depende críticamente del bloqueo de fase no lineal entre un oscilador local y una señal entrante (véase la \cref{fig:pll_block_diagram}).
\end{itemize}

\begin{figure}[htbp]
\centering
\begin{tikzpicture}[>=Stealth, auto, node distance=2.2cm, >=Latex]
    \tikzstyle{block} = [draw, rectangle, minimum height=1.1cm, minimum width=2.2cm, fill=black!5, font=\small]
    \tikzstyle{input} = [coordinate]
    \tikzstyle{output} = [coordinate]
    
    \node [input] (in) {};
    \node [block, right of=in, node distance=2.0cm] (pd) {$\Phi$ \begin{tabular}{c}Detector\\de fase\end{tabular}};
    \node [block, right of=pd, node distance=3.2cm] (lf) {\begin{tabular}{c}Filtro\\de lazo\end{tabular}};
    \node [block, right of=lf, node distance=3.2cm] (vco) {\begin{tabular}{c}VCO\\Oscilador\end{tabular}};
    \node [output, right of=vco, node distance=2.5cm] (out) {};
    \node [coordinate, below of=lf, node distance=1.6cm] (fb) {};

    \draw [->, thick] (in) -- node[above] {$v_i(t)$} (pd);
    \draw [->, thick] (pd) -- node[above] {$e(t)$} (lf);
    \draw [->, thick] (lf) -- node[above] {$v_c(t)$} (vco);
    \draw [->, thick] (vco) -- node[above] (outlabel) {$v_o(t)$} (out);
    \draw [->, thick] ($(vco.east)!0.5!(out)$) |- (fb) -| (pd.south);
\end{tikzpicture}
\caption{Diagrama de bloques de un lazo de enganche de fase (\emph{Phase Locked Loop}, PLL). La señal de entrada $v_i(t)$ se compara con la salida $v_o(t)$ generada por el oscilador controlado por voltaje (VCO); el error pasa por un filtro de lazo que corrige dinámicamente la frecuencia y la fase.}
\label{fig:pll_block_diagram}
\end{figure}

Hacia 1950, la invención de las computadoras digitales posibilitó simular numéricamente sistemas de ecuaciones no integrables. En 1963, el matemático y meteorólogo Edward Lorenz, trabajando en el MIT en modelos de convección atmosférica basados en las ecuaciones de Navier--Stokes y Rayleigh--Bénard, redujo el sistema a tres ecuaciones diferenciales no lineales acopladas. Al simularlas, descubrió trayectorias que nunca se repetían pero que permanecían confinadas en un objeto geométrico fractal: el primer atractor extraño documentado.

Lorenz publicó su trabajo seminal \emph{``Deterministic Nonperiodic Flow''} en 1963 en el \emph{Journal of the Atmospheric Sciences} \cite{lorenz1963}. Inicialmente, el artículo pasó inadvertido para la comunidad matemática porque estaba publicado en una revista especializada en meteorología, y los físicos lo consideraron en su momento un modelo de juguete demasiado simplificado.

\subsection{Caos determinista, biología y universalidad (1963--1990s)}

A partir de los años sesenta y setenta, la revolución de los sistemas dinámicos se extendió a múltiples disciplinas:
\begin{itemize}
    \item \textbf{Topología y matemáticas puras:} Stephen Smale introdujo la transformación de herradura (\emph{Smale horseshoe}), demostrando la estructura geométrica robusta del caos. En mecánica hamiltoniana, el teorema KAM (Kolmogorov--Arnold--Moser) explicó la persistencia y destrucción paulatina de toros invariantes en sistemas casi-integrables.
    \item \textbf{Biología de poblaciones y mapas discretos:} En 1976, Robert M. May publicó en \emph{Nature} su célebre artículo \emph{``Simple Mathematical Models with Very Complicated Dynamics''} \cite{may1976}. May demostró que la ecuación logística discreta $x_{n+1} = r x_n(1 - x_n)$, a pesar de su sencillez, exhibe una cascada infinita de duplicación de periodo y caos. May hizo un enérgico llamado pedagógico a no restringir la enseñanza a modelos lineales: \emph{``si permitimos a los sistemas ser no lineales, es donde surge la verdadera riqueza dinámica''}.
    \item \textbf{Geometría fractal:} Benoit Mandelbrot acuñó el término \emph{fractal} y reveló que las geometrías de la naturaleza (costas, nubes, bronquios, redes vasculares) son inherentemente no euclidianas y auto-similares.
    \item \textbf{Osciladores biológicos:} Arthur T. Winfree investigó la sincronización biológica en ritmos circadianos, agregación de células cardíacas y ondas espaciotemporales \cite{winfree2001}.
    \item \textbf{Turbulencia y dinámica de fluidos:} David Ruelle y Floris Takens (1971) propusieron que la turbulencia en fluidos no surge de una suma infinita de frecuencias independientes (como sugería Landau), sino del colapso del flujo hacia un atractor extraño de baja dimensión \cite{ruelle1971}.
    \item \textbf{Universalidad cuantitativa:} Hacia 1978, el físico Mitchell Feigenbaum \cite{feigenbaum1978} descubrió que la tasa geométrica con la que se suceden las bifurcaciones de duplicación de periodo converge a una constante universal:
    \begin{equation}
        \delta = \lim_{k \to \infty} \frac{r_k - r_{k-1}}{r_{k+1} - r_k} \approx 4.6692016\dots
    \end{equation}
    aplicando técnicas de grupo de renormalización tomadas de las transiciones de fase de la física estadística.
    \item \textbf{Difusión cultural y aplicaciones de ingeniería:} En 1987, James Gleick publicó el superventas \emph{Chaos: Making a New Science} \cite{gleick1987}, llevando estos conceptos al público general (popularizados luego en la cultura masiva con la novela y película \emph{Jurassic Park} en 1993 y la famosa metáfora del ``efecto mariposa''). En la década de 1990, el caos se aplicó al cifrado de comunicaciones seguras, sincronización de láseres y control de arritmias cardíacas.
\end{itemize}

\aside{Un ejemplo fascinante de la mecánica celeste moderna: en 1985, los científicos de la NASA emplearon la dinámica no lineal de $N$ cuerpos y trayectorias caóticas alrededor de los puntos de Lagrange para redirigir la sonda espacial ISEE-3 (renombrada ICE, \emph{International Comet Explorer}) y realizar el primer sobrevuelo exitoso de un cometa (Giacobini--Zinner) consumiendo un mínimo de combustible.}

\subsection{Formulación matemática general}

Formalmente, un sistema dinámico determinista continuo en tiempo se describe mediante un sistema de ecuaciones diferenciales ordinarias autónomas de primer orden:
\begin{equation}
    \dot{\vec{x}} \equiv \frac{d\vec{x}}{dt} = \vec{F}(\vec{x}), \quad \vec{x} \in \mathbb{R}^n, \quad \vec{F}: \mathbb{R}^n \to \mathbb{R}^n
    \label{eq:sistema_general}
\end{equation}
donde el espacio euclidiano $\mathbb{R}^n$ (o una variedad diferenciable $M$) se denomina el \textbf{espacio de fases} (o espacio de estados), y $\vec{x}(t) = (x_1(t), x_2(t), \dots, x_n(t))^\top$ describe el estado instantáneo del sistema.

En coordenadas escalares, el sistema \eqref{eq:sistema_general} toma la forma explícita:
\begin{equation}
\begin{aligned}
    \dot{x}_1 &= f_1(x_1, x_2, \dots, x_n) \\
    \dot{x}_2 &= f_2(x_1, x_2, \dots, x_n) \\
    &\;\;\vdots \\
    \dot{x}_n &= f_n(x_1, x_2, \dots, x_n)
\end{aligned}
\label{eq:sistema_escalar}
\end{equation}

\begin{definition}[Linealidad versus No linealidad]
    Decimos que el sistema \eqref{eq:sistema_general} es \textbf{lineal} si el campo vectorial $\vec{F}(\vec{x})$ satisface el principio de superposición, es decir, $\vec{F}(\alpha\vec{x} + \beta\vec{y}) = \alpha\vec{F}(\vec{x}) + \beta\vec{F}(\vec{y})$, lo cual equivale a que todas las funciones $f_i$ dependan únicamente de combinaciones lineales homogéneas de primer orden:
    \begin{equation}
        \dot{\vec{x}} = A \vec{x}, \quad A \in \mathbb{R}^{n \times n}
    \end{equation}
    En caso contrario, si en las funciones $f_i$ aparecen potencias ($x_i^2, x_i^3$), productos cruzados entre variables ($x_i x_j$) o funciones trascendentes ($\sin x_i, e^{x_i}, \tanh x_i$), el sistema se clasifica como \textbf{no lineal}.
\end{definition}

En los sistemas no lineales, el principio de superposición se quiebra: el todo no es simplemente la suma de las partes. Por ello, el análisis requiere estudiar la geometría del campo vectorial en el espacio de fases, como veremos en el siguiente capítulo.



% ---- Parte II — Flujos en una dimensión (draft págs. 7--21) ----
% ============================================================
% FUENTES: draft p.7--8; "Notes_ Non Linear Dynamics.pdf" (inglés),
% sección "2.1 A Geometric way of thinking" (campo vectorial, x'=sin x).
% ============================================================
\section{Espacio de fases y pensamiento geométrico}

La premisa central de la dinámica no lineal es sustituir la búsqueda obsesiva de fórmulas analíticas por una visión geométrica global. Como enfatizaba Henri Poincaré, el objetivo no es integrar mecánicamente una ecuación diferencial para obtener una función complicada e inescrutable, sino entender el comportamiento cualitativo de todas las soluciones posibles a partir de la geometría del campo vectorial.

\subsection{El péndulo: oscilador lineal vs. péndulo no lineal}

Para ilustrar este cambio de perspectiva, comparemos el oscilador armónico lineal clásico con el péndulo simple no lineal.

\paragraph{El oscilador armónico lineal.}
Una masa $m$ sujeta a un resorte lineal con constante de restitución $k$ (sin fricción) satisface:
\begin{equation}
    m \ddot{x} + k x = 0 \implies \ddot{x} + \omega^2 x = 0, \quad \omega = \sqrt{k/m}
\end{equation}
Al definir la velocidad $v = \dot{x}$, este sistema de segundo orden se desacopla en el sistema lineal en el plano $(x, v)$:
\begin{align}
    \dot{x} &= v \label{eq:armonico_x}\\
    \dot{v} &= -\omega^2 x \label{eq:armonico_v}
\end{align}
Las trayectorias en el plano $(x, v)$ son elipses concéntricas cerradas centradas en el origen, correspondientes a la conservación de la energía mecánica $E = \frac{1}{2}mv^2 + \frac{1}{2}kx^2 = \text{cte}$.

\paragraph{El péndulo no lineal.}
Si consideramos un péndulo rígido de longitud $L$ y masa $m$ bajo la gravedad $g$, el balance de momento angular sin la aproximación de oscilaciones pequeñas ($\sin x \approx x$) conduce a la ecuación no lineal:
\begin{equation}
    \ddot{x} + \frac{g}{L} \sin x = 0
\end{equation}
donde $x(t)$ representa el ángulo de desviación respecto a la vertical inferior. Escogiendo unidades tales que $g/L = 1$ y definiendo la velocidad angular $v = \dot{x}$, obtenemos el sistema bidimensional autónomo no lineal:
\begin{align}
    \dot{x} &= v \label{eq:pendulo_x}\\
    \dot{v} &= -\sin x \label{eq:pendulo_v}
\end{align}

\subsection{Pensar la ecuación diferencial como campo vectorial (Poincaré)}

La ecuación \eqref{eq:pendulo_x}--\eqref{eq:pendulo_v} asigna a cada punto del plano de estado $(x, v)$ un vector de velocidad de fase $(\dot{x}, \dot{v}) = (v, -\sin x)$. Una solución $(x(t), v(t))$ no es más que una curva parametrizada por el tiempo $t$ que fluye siguiendo tangencialmente las flechas del campo vectorial en cada instante.

\begin{definition}[Retrato de fases (\emph{Phase Portrait})]
    Es la representación geométrica global de las trayectorias de un sistema dinámico en su espacio de estados, mostrando todos los comportamientos cualitativos posibles (puntos de equilibrio, órbitas periódicas, separatrices y cuencas de atracción) sin necesidad de resolver analíticamente las ecuaciones.
\end{definition}

\aside{Intentar resolver analíticamente ecuaciones como $\dot{x} = \sin x$ mediante separación de variables conduce a expresiones implícitas engorrosas:
\[
    \int \csc x \, dx = \int dt \implies -\ln|\csc x + \cot x| = t + C
\]
Esta fórmula es prácticamente inútil para responder preguntas dinámicas elementales (¿hacia dónde tiende el sistema cuando $t \to \infty$ si $x(0) = \pi/4$?). En cambio, observar el signo de $\sin x$ en el eje $x$ responde la pregunta en un segundo.}

\subsection{Libraciones, rotaciones y la separatrix}

Al trazar el retrato de fases del péndulo no lineal en el plano $(x, v)$ (véase la \cref{fig:retratos_fase_poincare}), descubrimos una estructura cualitativa mucho más rica que la del oscilador armónico:

\begin{enumerate}
    \item \textbf{Libraciones (órbitas cerradas):} Alrededor del equilibrio estable inferior $(x, v) = (0, 0)$ y de sus réplicas periódicas $(2k\pi, 0)$, existen curvas cerradas anidadas. Corresponden a oscilaciones periódicas confinadas donde el péndulo oscila hacia adelante y hacia atrás sin llegar a la posición vertical superior.
    \item \textbf{Puntos de ensilladura (\emph{saddles}):} En $(x, v) = (\pm \pi, 0)$, el péndulo se encuentra en equilibrio inestable invertido (la posición vertical superior).
    \item \textbf{Separatriz (\emph{separatrix}):} Es la trayectoria límite (órbita homoclínica) de energía crítica $E_c$ que une los puntos de ensilladura inestables, dada analíticamente por la curva $v(x) = \pm 2\cos(x/2)$. La separatriz divide el espacio de fases en dos regímenes dinámicamente incompatibles.
    \item \textbf{Rotaciones (curvas ondulantes abiertas):} Por encima y por debajo de la separatriz ($|v| > 2$ cuando $x=0$), la energía cinética es suficiente para superar la barrera de potencial gravitacional: el péndulo nunca se detiene y gira indefinidamente en sentido antihorario ($v > 0$) o horario ($v < 0$).
\end{enumerate}

\begin{figure}[htbp]
\centering
\begin{tikzpicture}[>=Stealth, scale=0.85]
    % --- Panel Izquierdo: Oscilador armónico ---
    \begin{scope}[shift={(-4.5,0)}]
        \draw[->, thick, black!60] (-3,0) -- (3,0) node[right, font=\small] {$x$};
        \draw[->, thick, black!60] (0,-2.5) -- (0,2.5) node[above, font=\small] {$v$};
        \node[above, font=\bfseries\small] at (0, 2.7) {Oscilador armónico (lineal)};
        
        \foreach \r in {0.7, 1.4, 2.1} {
            \draw[thick, black!80] (0,0) ellipse ({\r*1.1} and {\r*0.85});
            \draw[->, thick, black!80] ({\r*1.1*cos(45)}, {\r*0.85*sin(45)}) -- ++(0.01,-0.02);
            \draw[->, thick, black!80] ({-\r*1.1*cos(45)}, {-\r*0.85*sin(45)}) -- ++(-0.01,0.02);
        }
        \filldraw[black] (0,0) circle (2pt) node[below right, font=\footnotesize] {Centro $(0,0)$};
        \node[font=\footnotesize, align=center] at (0, -3.2) {Órbitas elípticas concéntricas\\(todas de igual periodo)};
    \end{scope}

    % --- Panel Derecho: Péndulo no lineal ---
    \begin{scope}[shift={(4.5,0)}]
        \draw[->, thick, black!60] (-3.8,0) -- (3.8,0) node[right, font=\small] {$x$};
        \draw[->, thick, black!60] (0,-2.5) -- (0,2.5) node[above, font=\small] {$v$};
        \node[above, font=\bfseries\small] at (0, 2.7) {Péndulo simple (no lineal)};
        
        % Centros y sillas
        \filldraw[black] (0,0) circle (2pt);
        \draw[black, fill=white, thick] (-3.1416, 0) circle (2pt) node[below, font=\footnotesize] {$-\pi$};
        \draw[black, fill=white, thick] (3.1416, 0) circle (2pt) node[below, font=\footnotesize] {$\pi$};
        \node[below, font=\footnotesize] at (0,0) {$0$};

        % Libraciones (órbitas cerradas)
        \draw[thick, black!80] (0,0) ellipse (1.0 and 0.6);
        \draw[thick, black!80] (0,0) ellipse (2.0 and 1.2);
        \draw[->, thick, black!80] (0, 1.2) -- ++(0.1,0);
        \draw[->, thick, black!80] (0, -1.2) -- ++(-0.1,0);

        % Separatriz
        \draw[very thick, black!90, domain=-3.1415:3.1415, samples=60] 
            plot (\x, {1.9*cos(\x/2 r)});
        \draw[very thick, black!90, domain=-3.1415:3.1415, samples=60] 
            plot (\x, {-1.9*cos(\x/2 r)});
        
        % Rotaciones superiores e inferiores
        \draw[thick, dashed, black!70, domain=-3.6:3.6, samples=60] 
            plot (\x, {2.2 + 0.2*cos(\x r)});
        \draw[->, thick, black!70] (0, 2.4) -- ++(0.2, 0);
        \draw[thick, dashed, black!70, domain=-3.6:3.6, samples=60] 
            plot (\x, {-2.2 - 0.2*cos(\x r)});
        \draw[->, thick, black!70] (0, -2.4) -- ++(-0.2, 0);

        % Etiquetas descriptivas
        \node[font=\tiny, fill=white, inner sep=1pt] at (0, 0.7) {Libraciones};
        \node[font=\tiny, fill=white, inner sep=1pt] at (1.8, 1.6) {Separatriz};
        \node[font=\tiny, fill=white, inner sep=1pt] at (0, 2.2) {Rotación ($v>0$)};
        \node[font=\footnotesize, align=center] at (0, -3.2) {Libraciones, rotaciones\\y separatriz homoclínica};
    \end{scope}
\end{tikzpicture}
\caption{Comparación geométrica en el espacio de fases $(x, v)$. Izquierda: oscilador armónico lineal con órbitas elípticas globales. Derecha: péndulo no lineal con libraciones cerradas alrededor de $(0,0)$, separatriz crítica que conecta los equilibrios inestables $(\pm\pi, 0)$ y rotaciones abiertas para energías superiores.}
\label{fig:retratos_fase_poincare}
\end{figure}


% ============================================================
% FUENTES: draft p.9--11; inglés §2 intro (n=1, sistemas de primer orden,
% restricción a sistemas autónomos f(x,t)->f(x)).
% ============================================================
\section{Flujos en una dimensión}
\label{sec:flujos1d}

En esta sección exploramos la clase más sencilla de sistemas dinámicos continuos: los flujos en una línea (espacio de fases unidimensional, $n=1$). Veremos cómo la representación geométrica supera las limitaciones del cálculo analítico tradicional.

\subsection{El ejemplo lineal $\dot{x} = ax$}

Consideremos el problema de valor inicial para el sistema lineal homogéneo de primer orden:
\begin{equation}
    \dot{x} = ax, \quad x(0) = x_0, \quad a \in \mathbb{R}
    \label{eq:flujo_lineal}
\end{equation}

\paragraph{Solución analítica.}
Separando variables de manera rigurosa:
\begin{align}
    \frac{dx}{dt} = ax &\implies \frac{1}{x} \frac{dx}{dt} = a \nonumber \\
    &\implies \int \frac{1}{x} \frac{dx}{dt} \, dt = \int a \, dt \nonumber \\
    &\implies \ln|x| = at + C \nonumber \\
    &\implies x(t) = x_0 e^{at}
\end{align}

\paragraph{Interpretación geométrica en la línea de fase.}
En lugar de integrar, graficamos el campo vectorial $f(x) = ax$ frente a $x$ (véase la \cref{fig:flujo_lineal_ax}).
\begin{itemize}
    \item Si $a > 0$: $f(x) > 0$ cuando $x > 0$, lo que significa que la velocidad es positiva ($\dot{x} > 0$) y el estado se desplaza hacia la derecha. Para $x < 0$, $f(x) < 0 \implies \dot{x} < 0$, desplazándose hacia la izquierda. El origen $x^* = 0$ es un punto de velocidad nula; cualquier pequeña perturbación aleja al sistema de $0$. Por ello, $x^* = 0$ es un \textbf{punto fijo inestable} o \textbf{repulsor}.
    \item Si $a < 0$: el signo del campo se invierte; las trayectorias convergen exponencialmente hacia $x^* = 0$, convirtiéndolo en un \textbf{punto fijo estable} o \textbf{atractor}.
\end{itemize}

\begin{figure}[htbp]
\centering
\begin{tikzpicture}[>=Stealth, scale=0.9]
    \draw[->, thick, black!60] (-4,0) -- (4,0) node[right, font=\small] {$x$};
    \draw[->, thick, black!60] (0,-2.5) -- (0,2.5) node[above, font=\small] {$\dot{x} = f(x)$};
    
    % Recta f(x) = a x con a > 0
    \draw[very thick, black!80] (-3.2,-2.0) -- (3.2,2.0) node[right, font=\small] {$f(x) = ax \ (a>0)$};
    
    % Flechas sobre el eje x
    \draw[->, very thick, black!90] (0.5, 0) -- (1.5, 0);
    \draw[->, very thick, black!90] (2.0, 0) -- (3.0, 0);
    \draw[->, very thick, black!90] (-0.5, 0) -- (-1.5, 0);
    \draw[->, very thick, black!90] (-2.0, 0) -- (-3.0, 0);
    
    % Punto fijo inestable en el origen
    \draw[black, fill=white, thick] (0,0) circle (2.5pt) node[below right=2pt, font=\footnotesize] {$x^*=0$ (inestable)};
    
    \node[font=\footnotesize, black!70, align=center] at (2.2, -1.2) {Velocidad positiva ($\dot{x} > 0$)\\\emph{flujo hacia la derecha}};
    \node[font=\footnotesize, black!70, align=center] at (-2.2, 1.2) {Velocidad negativa ($\dot{x} < 0$)\\\emph{flujo hacia la izquierda}};
\end{tikzpicture}
\caption{Línea de fase y campo vectorial para el sistema lineal $\dot{x} = ax$ con $a > 0$. El origen $x^* = 0$ es un punto de equilibrio inestable del cual divergen todas las trayectorias.}
\label{fig:flujo_lineal_ax}
\end{figure}

\subsection{El ejemplo no lineal $\dot{x} = \sin x$}

Consideremos ahora un sistema no lineal fundamental:
\begin{equation}
    \dot{x} = \sin x, \quad x(0) = x_0
    \label{eq:flujo_sinx}
\end{equation}

\paragraph{Integración analítica.}
Separando variables e integrando mediante la sustitución trigonométrica estándar:
\begin{align}
    \int \frac{1}{\sin x} \, dx = \int dt &\implies \int \csc x \, dx = t + C \nonumber \\
    &\implies -\ln|\csc x + \cot x| = t + C
\end{align}
Imponiendo la condición inicial en $t=0$, $x(0) = x_0$:
\begin{equation}
    t = \ln \left| \frac{\csc x_0 + \cot x_0}{\csc x(t) + \cot x(t)} \right|
    \label{eq:sol_implicita_sinx}
\end{equation}
La ecuación \eqref{eq:sol_implicita_sinx} es exacta, pero extraordinariamente opaca: despejar explícitamente $x(t)$ requiere manipulaciones algebraicas engorrosas y no ofrece intuición directa sobre la física global del sistema.

\paragraph{Análisis cualitativo y línea de fase.}
Grafiquemos $f(x) = \sin x$ (véase la \cref{fig:flujo_sin_x}). Los puntos de equilibrio $\dot{x} = 0$ ocurren en los múltiplos enteros de $\pi$:
\begin{equation}
    x_k^* = k\pi, \quad k \in \mathbb{Z}
\end{equation}

\begin{itemize}
    \item En el intervalo $(0, \pi)$, $\sin x > 0 \implies \dot{x} > 0$: el flujo se mueve hacia la derecha.
    \item En el intervalo $(\pi, 2\pi)$, $\sin x < 0 \implies \dot{x} < 0$: el flujo se mueve hacia la izquierda.
    \item En consecuencia, $x^* = 0, \pm 2\pi, \dots$ son \textbf{repulsores inestables}, mientras que $x^* = \pm \pi, \pm 3\pi, \dots$ son \textbf{atractores estables}.
\end{itemize}

\begin{figure}[htbp]
\centering
\begin{tikzpicture}[>=Stealth, scale=0.95]
    \draw[->, thick, black!60] (-5.5,0) -- (5.5,0) node[right, font=\small] {$x$};
    \draw[->, thick, black!60] (0,-1.8) -- (0,1.8) node[above, font=\small] {$\dot{x} = \sin x$};
    
    % Curva sin(x)
    \draw[very thick, black!85, domain=-5.2:5.2, samples=100] plot (\x, {1.2*sin(\x r)});
    
    % Puntos fijos
    \draw[black, fill=white, thick] (-3.1416, 0) circle (2.5pt); % -pi estable
    \filldraw[black] (-3.1416, 0) circle (2.5pt) node[below=3pt, font=\footnotesize] {$-\pi$};
    
    \draw[black, fill=white, thick] (0,0) circle (2.5pt) node[below=3pt, font=\footnotesize] {$0$};
    
    \filldraw[black] (3.1416, 0) circle (2.5pt) node[below=3pt, font=\footnotesize] {$\pi$};
    
    % Flechas de fase
    \draw[->, very thick, black!90] (-2.2, 0) -- (-2.8, 0);
    \draw[->, very thick, black!90] (-4.2, 0) -- (-3.5, 0);
    \draw[->, very thick, black!90] (0.8, 0) -- (1.8, 0);
    \draw[->, very thick, black!90] (2.2, 0) -- (2.8, 0);
    \draw[->, very thick, black!90] (4.5, 0) -- (3.6, 0);
    
    \node[font=\footnotesize, black!70] at (1.5, 1.4) {$\dot{x} > 0$ (velocidad positiva)};
    \node[font=\footnotesize, black!70] at (4.5, -1.2) {$\dot{x} < 0$ (velocidad negativa)};
\end{tikzpicture}
\caption{Campo vectorial y línea de fase para $\dot{x} = \sin x$. Los círculos negros sólidos representan puntos fijos estables (atractores en $\pm \pi$), y los círculos blancos abiertos representan equilibrios inestables (repulsores en $0, \pm 2\pi$).}
\label{fig:flujo_sin_x}
\end{figure}

Si una partícula inicia en $x_0 = \pi/4$, como $\sin(\pi/4) > 0$, el flujo la empuja irrevocablemente hacia la derecha con velocidad máxima en $x = \pi/2$ (donde la derivada $\ddot{x} = \cos x \cdot \dot{x} = 0$, punto de inflexión), desacelerando suavemente hasta alcanzar asintóticamente $x^* = \pi$ cuando $t \to \infty$ (véase la \cref{fig:trayectoria_sin_x}).

\begin{figure}[htbp]
\centering
\begin{tikzpicture}[>=Stealth, scale=0.9]
    \draw[->, thick, black!60] (0,0) -- (6,0) node[right, font=\small] {$t$};
    \draw[->, thick, black!60] (0,0) -- (0,3.8) node[above, font=\small] {$x(t)$};
    
    % Asíntota en pi
    \draw[dashed, black!50, thick] (0, 3.1416) -- (5.5, 3.1416) node[right, font=\footnotesize] {$x = \pi$ (asíntota estable)};
    
    % Trayectoria sigmoide
    \draw[very thick, black!85, domain=0:5.2, samples=60] 
        plot (\x, {2*rad(atan(exp(\x - 1.2))) + 0.3});
    
    % Puntos clave
    \filldraw[black] (0, 0.785) circle (2pt) node[left, font=\footnotesize] {$x_0 = \pi/4$};
    \draw[dotted, thick, black!60] (1.2, 0) node[below, font=\footnotesize] {$t_{\text{inflexión}}$} -- (1.2, 1.57) -- (0, 1.57) node[left, font=\footnotesize] {$\pi/2$};
    \filldraw[black] (1.2, 1.57) circle (2pt);
    
    \node[font=\footnotesize, black!75, align=left] at (4.0, 1.2) {Convergencia asintótica:\\$\lim_{t \to \infty} x(t) = \pi$};
\end{tikzpicture}
\caption{Evolución temporal de la solución $x(t)$ para $\dot{x} = \sin x$ con condición inicial $x(0) = \pi/4$. La trayectoria presenta una forma sigmoidal con punto de inflexión en $x = \pi/2$ y se nivela asintóticamente en el atractor $x^* = \pi$.}
\label{fig:trayectoria_sin_x}
\end{figure}

\subsection{Definición formal: sistemas de primer orden, autónomos}

Formalicemos el marco teórico de los flujos unidimensionales:

\begin{definition}[Flujo autónomo unidimensional]
    Un sistema dinámico unidimensional continuo y autónomo se rige por la ecuación diferencial:
    \begin{equation}
        \dot{x} = f(x), \quad x \in \mathbb{R}
        \label{eq:flujo_1d_def}
    \end{equation}
    donde la función $f: \mathbb{R} \to \mathbb{R}$ es de clase $\mathcal{C}^1$ (continuamente diferenciable), garantizando por el teorema de Picard--Lindelöf la existencia y unicidad local de las trayectorias.
\end{definition}

\paragraph{Dos precisiones conceptuales fundamentales:}
\begin{enumerate}
    \item \textbf{El significado de ``sistema'':} Empleamos el término en el sentido moderno de la teoría de sistemas dinámicos (un estado determinista que evoluciona bajo una regla de velocidad en su espacio de estados), no en el sentido escolar de un conjunto acoplado de múltiples ecuaciones simultáneas. Una única ecuación diferencial $\dot{x} = f(x)$ es un sistema dinámico unidimensional.
    \item \textbf{Restricción a ecuaciones autónomas:} A lo largo de esta parte del curso asumiremos siempre que $f$ no depende explícitamente del tiempo ($f(x)$, no $f(x, t)$). Una ecuación no autónoma $\dot{x} = f(x, t)$ requiere especificar dos cantidades $(x, t)$ para determinar la velocidad instantánea, por lo que matemáticamente equivale a un sistema bidimensional en el espacio aumentado $(x, t)$.
\end{enumerate}

\subsection{Ejercicios}

\paragraph{Ejercicio 1 (mecánico; línea de fase).} Para $\dot{x} = x^2 - 1$, halla los puntos fijos, determina el signo de $f(x)$ en cada intervalo que definen y dibuja la línea de fase con sus flechas. Clasifica cada equilibrio como estable o inestable a partir únicamente de la dirección del flujo.
\aside{Pista: $f(x)=(x-1)(x+1)$; estudia el signo en $(-\infty,-1)$, $(-1,1)$, $(1,\infty)$.}
% SOL: FP x=-1 (estable, f pasa de + a -) y x=1 (inestable, f pasa de - a +).

\paragraph{Ejercicio 2 (análisis; destino sin integrar).} Sea $\dot{x} = \sin x$ con $x(0) = \pi/4$. Sin usar la solución implícita \eqref{eq:sol_implicita_sinx}, determina $\lim_{t\to\infty} x(t)$ y el valor de $x$ (dentro de la trayectoria) donde la velocidad $|\dot x|$ alcanza su máximo. Justifícalo con el signo de $\sin x$ y con la condición de inflexión $\ddot x=0$.
\aside{Pista: en $(0,\pi)$, $\sin x>0$; el flujo sube hacia el atractor $\pi$. La velocidad es máxima donde $\ddot x=\cos x\,\dot x=0$, es decir en $x=\pi/2$.}
% SOL: lim = pi (atractor estable). La trayectoria parte de pi/4, cruza x=pi/2 (velocidad maxima, sin=1) y se nivela en pi. Coincide con la Fig. de la trayectoria sigmoide.

\paragraph{Ejercicio 3 (interpretación; autonomía).} Explica por qué $\dot{x} = x + \sin t$ \emph{no} es un flujo unidimensional genuino en el sentido de la \cref{eq:flujo_1d_def}, y describe el espacio de fases mínimo en el que sí puede representarse como sistema autónomo.
\aside{Pista: la velocidad depende de $t$; añade la ecuación trivial $\dot\tau=1$ para absorber el tiempo como variable de estado.}
% SOL: depende de t explicitamente -> no autonomo. Se autonomiza en el plano (x,tau) con xdot=x+sin(tau), taudot=1: es un sistema 2D.


% ============================================================
% FUENTES: draft p.12--13; inglés §2.2 "Fixed Points & Stability".
% ============================================================
\section{Puntos fijos y estabilidad}

En los sistemas continuos, los estados donde el sistema cesa de cambiar son los elementos vertebradores de toda la dinámica.

\subsection{Clasificación cualitativa en la línea de fase}

Consideremos el flujo unidimensional general $\dot{x} = f(x)$.

\begin{definition}[Punto fijo o estado de equilibrio]
    Un punto $x^* \in \mathbb{R}$ es un \textbf{punto fijo} (o punto de equilibrio) del sistema $\dot{x} = f(x)$ si:
    \begin{equation}
        f(x^*) = 0
    \end{equation}
    Si la condición inicial coincide con el punto fijo, $x(0) = x^*$, la solución constante $x(t) \equiv x^*$ es válida para todo $t \ge 0$.
\end{definition}

El comportamiento en la vecindad de $x^*$ depende exclusivamente del signo de $f(x)$ a ambos lados:

\begin{enumerate}
    \item \textbf{Punto fijo estable (atractor / sumidero / \emph{sink}):} Si $f(x) > 0$ a la izquierda de $x^*$ y $f(x) < 0$ a la derecha, el flujo empuja a las partículas hacia $x^*$ desde ambas direcciones ($\to \bullet \leftarrow$). Geométricamente, la gráfica de $f(x)$ cruza el eje $x$ con pendiente negativa ($f'(x^*) < 0$).
    \item \textbf{Punto fijo inestable (repulsor / fuente / \emph{source}):} Si $f(x) < 0$ a la izquierda y $f(x) > 0$ a la derecha, el flujo expulsa a las partículas lejos de $x^*$ ($\leftarrow \circ \to$). Geométricamente, $f(x)$ cruza el eje $x$ con pendiente positiva ($f'(x^*) > 0$).
    \item \textbf{Punto fijo semi-estable (nodo semi-estable / cuello de botella):} Ocurre cuando la gráfica de $f(x)$ toca tangencialmente el eje $x$ sin cruzarlo (por ejemplo, $f(x) = (x-x^*)^2$ o $-(x-x^*)^2$). El flujo atrae por un flanco pero repele por el opuesto ($\to \bullet \to$ o $\leftarrow \bullet \leftarrow$).
\end{enumerate}

La \cref{fig:clasificacion_puntos_fijos} resume los cuatro casos prototípicos en la línea de fase.

\begin{figure}[htbp]
\centering
\begin{tikzpicture}[>=Stealth, scale=0.82]
    % Panel 1: Parábola arriba
    \begin{scope}[shift={(-6.5, 3.2)}]
        \draw[->, thick, black!50] (-2.2,0) -- (2.2,0) node[right, font=\tiny] {$x$};
        \draw[->, thick, black!50] (0,-1.2) -- (0,1.5) node[above, font=\tiny] {$f(x)$};
        \draw[thick, black!85, domain=-1.8:1.8] plot (\x, {\x*\x - 0.8});
        \filldraw[black] (-0.894, 0) circle (2pt) node[above left, font=\tiny] {$a$};
        \draw[black, fill=white, thick] (0.894, 0) circle (2pt) node[above right, font=\tiny] {$b$};
        \draw[->, thick, black!90] (-1.8,0) -- (-1.2,0);
        \draw[->, thick, black!90] (0,0) -- (-0.5,0);
        \draw[->, thick, black!90] (1.2,0) -- (1.7,0);
        \node[font=\footnotesize, align=center] at (0,-1.6) {\textbf{(a)} $a$ estable, $b$ inestable};
    \end{scope}

    % Panel 2: Tangencia (semi-estable)
    \begin{scope}[shift={(1.5, 3.2)}]
        \draw[->, thick, black!50] (-3.0,0) -- (3.0,0) node[right, font=\tiny] {$x$};
        \draw[->, thick, black!50] (0,-1.2) -- (0,1.5) node[above, font=\tiny] {$f(x)$};
        \draw[thick, black!85, domain=-2.3:2.3, samples=60] plot (\x, {-0.3*(\x+1.6)*(\x)*(\x)*(\x-1.6)});
        \filldraw[black] (-1.6, 0) circle (2pt) node[below, font=\tiny] {$a$};
        \filldraw[black!60] (0, 0) circle (2pt) node[below right, font=\tiny] {$b$};
        \draw[black, fill=white, thick] (1.6, 0) circle (2pt) node[below, font=\tiny] {$c$};
        \draw[->, thick, black!90] (-2.2,0) -- (-1.8,0);
        \draw[->, thick, black!90] (-0.8,0) -- (-1.2,0);
        \draw[->, thick, black!90] (0.5,0) -- (1.0,0);
        \draw[->, thick, black!90] (2.2,0) -- (2.6,0);
        \node[font=\footnotesize, align=center] at (0,-1.6) {\textbf{(b)} $a$ estable, $b$ semi-estable, $c$ inestable};
    \end{scope}

    % Panel 3: Lineal decreciente
    \begin{scope}[shift={(-6.5, -2.5)}]
        \draw[->, thick, black!50] (-2.2,0) -- (2.2,0) node[right, font=\tiny] {$x$};
        \draw[->, thick, black!50] (0,-1.2) -- (0,1.5) node[above, font=\tiny] {$f(x)$};
        \draw[thick, black!85] (-1.5,1.2) -- (1.5,-1.2);
        \filldraw[black] (0, 0) circle (2pt) node[above right, font=\tiny] {$a$};
        \draw[->, thick, black!90] (-1.5,0) -- (-0.6,0);
        \draw[->, thick, black!90] (1.5,0) -- (0.6,0);
        \node[font=\footnotesize, align=center] at (0,-1.6) {\textbf{(c)} Lineal estable ($f'<0$)};
    \end{scope}

    % Panel 4: Lineal creciente
    \begin{scope}[shift={(1.5, -2.5)}]
        \draw[->, thick, black!50] (-2.2,0) -- (2.2,0) node[right, font=\tiny] {$x$};
        \draw[->, thick, black!50] (0,-1.2) -- (0,1.5) node[above, font=\tiny] {$f(x)$};
        \draw[thick, black!85] (-1.5,-1.2) -- (1.5,1.2);
        \draw[black, fill=white, thick] (0, 0) circle (2pt) node[below right, font=\tiny] {$a$};
        \draw[->, thick, black!90] (-0.6,0) -- (-1.5,0);
        \draw[->, thick, black!90] (0.6,0) -- (1.5,0);
        \node[font=\footnotesize, align=center] at (0,-1.6) {\textbf{(d)} Lineal inestable ($f'>0$)};
    \end{scope}
\end{tikzpicture}
\caption{Clasificación geométrica de puntos fijos en la línea de fase según el cruce de $f(x)$ con el eje horizontal: atractores estables ($\bullet$), repulsores inestables ($\circ$) y nodos semi-estables.}
\label{fig:clasificacion_puntos_fijos}
\end{figure}

\subsection{Definiciones formales (Lyapunov, asintótica)}

Para dar rigor matemático a los conceptos intuitivos de atracción y repulsión, recurrimos a las definiciones formales introducidas por Aleksandr Lyapunov:

\begin{definition}[Estabilidad en el sentido de Lyapunov]
    Un punto de equilibrio $x^*$ es \textbf{estable en el sentido de Lyapunov} si para cada $\varepsilon > 0$, existe un $\delta = \delta(\varepsilon) > 0$ tal que si:
    \begin{equation}
        \|x(0) - x^*\| < \delta \implies \|x(t) - x^*\| < \varepsilon, \quad \forall t \ge 0
    \end{equation}
    En palabras cotidianas: si una solución comienza ``suficientemente cerca'' de $x^*$ (dentro de la bola $\delta$), permanecerá acotada y arbitrariamente cerca (dentro de la bola $\varepsilon$) por todo el tiempo futuro.
\end{definition}

\begin{definition}[Estabilidad Asintótica]
    Un punto de equilibrio $x^*$ es \textbf{asintóticamente estable} si cumple simultáneamente dos condiciones:
    \begin{enumerate}
        \item Es estable en el sentido de Lyapunov.
        \item Es un atractor local: existe un radio $\delta_0 > 0$ tal que si $\|x(0) - x^*\| < \delta_0$, entonces:
        \begin{equation}
            \lim_{t \to \infty} \|x(t) - x^*\| = 0
        \end{equation}
    \end{enumerate}
    Es decir, las trayectorias no solo permanecen atrapadas en la vecindad de $x^*$, sino que convergen asintóticamente a él cuando el tiempo tiende a infinito\aside{Una pregunta que abordaremos en el análisis lineal: ¿a qué velocidad converge el sistema hacia el punto fijo? Como veremos, la velocidad de convergencia está gobernada por la escala de tiempo característica impuesta por $f'(x^*)$.}.
\end{definition}

\begin{figure}[htbp]
\centering
\begin{tikzpicture}[>=Stealth, scale=0.95]
    % Eje t
    \draw[->, thick, black!60] (0,0) -- (7.5,0) node[right, font=\small] {$t$};
    \draw[->, thick, black!60] (0,-2.5) -- (0,2.5) node[above, font=\small] {$x(t)$};
    
    % Punto fijo en 0
    \draw[thick, black!90] (0,0) -- (7.0,0) node[right, font=\footnotesize] {$x^*$};
    
    % Banda épsilon
    \draw[dashed, black!40, thick] (0,1.8) -- (7.0,1.8) node[right, font=\tiny] {$x^* + \varepsilon$};
    \draw[dashed, black!40, thick] (0,-1.8) -- (7.0,-1.8) node[right, font=\tiny] {$x^* - \varepsilon$};
    
    % Intervalo delta en t=0
    \draw[very thick, black!80] (0,-0.9) -- (0,0.9);
    \draw[black!80, fill=black!80] (0,0.9) circle (1.5pt) node[left, font=\tiny] {$x^* + \delta$};
    \draw[black!80, fill=black!80] (0,-0.9) circle (1.5pt) node[left, font=\tiny] {$x^* - \delta$};
    
    % Trayectorias convergentes
    \draw[thick, black!85, domain=0:6.5, samples=50] plot (\x, {0.8*exp(-0.7*\x)});
    \draw[thick, black!85, domain=0:6.5, samples=50] plot (\x, {-0.6*exp(-0.9*\x)});
    \filldraw[black] (0, 0.8) circle (2pt) node[above right, font=\tiny] {$x_0$};
    
    \node[font=\footnotesize, black!75, align=left] at (4.5, 1.2) {Estabilidad Asintótica:\\$\|x(t)-x^*\| < \varepsilon \ \forall t \ge 0$\\$\lim_{t\to\infty} x(t) = x^*$};
\end{tikzpicture}
\caption{Esquema conceptual de la estabilidad de Lyapunov y la estabilidad asintótica: la condición inicial dentro del intervalo $\delta$ garantiza permanencia dentro de $\varepsilon$ y decaimiento monótono hacia $x^*$.}
\label{fig:lyapunov_estabilidad}
\end{figure}


% ============================================================
% FUENTES: draft p.14--17.
% ============================================================
\section{Ejemplos y aplicaciones en una dimensión}

A continuación desarrollamos tres aplicaciones clásicas que demuestran la versatilidad del análisis de flujos unidimensionales en física, matemáticas y dinámica de poblaciones.

\subsection{Un circuito RC}

Consideremos un circuito eléctrico elemental compuesto por una resistencia $R$ y un capacitor $C$ conectados en serie con una fuente de voltaje continuo constante $V_0$ mediante un interruptor que se cierra en $t=0$ (véase la \cref{fig:circuito_rc}). Supongamos que el capacitor se encuentra descargado inicialmente, $Q(0) = 0$.

\begin{figure}[htbp]
\centering
\begin{tikzpicture}[>=Stealth, scale=0.85]
    % Esquema de circuito (panel izquierdo)
    \begin{scope}[shift={(-6.2, 0)}]
        % Fuente V0 (batería)
        \draw[thick] (0,0) -- (0,1.0);
        \draw[thick] (-0.3,1.0) -- (0.3,1.0); % placa corta (-)
        \draw[very thick] (-0.5,1.3) -- (0.5,1.3) node[left=8pt, font=\tiny] {$V_0$}; % placa larga (+)
        \draw[thick] (0,1.3) -- (0,2.5) -- (1.0,2.5);
        % Interruptor
        \draw[thick] (1.0,2.5) -- (1.5,2.9) node[above=1pt, font=\tiny] {$t=0$};
        \filldraw[black] (1.0,2.5) circle (1.5pt);
        \filldraw[black] (1.6,2.5) circle (1.5pt);
        % Resistencia R
        \draw[thick] (1.6,2.5) -- (2.0,2.5);
        \draw[thick, fill=black!5] (2.0,2.3) rectangle (3.2,2.7) node[pos=0.5, font=\tiny] {$R$};
        \draw[thick] (3.2,2.5) -- (4.4,2.5);
        % Capacitor C
        \draw[thick] (4.4,2.5) -- (4.4,1.4);
        \draw[very thick] (3.9,1.4) -- (4.9,1.4) node[right=4pt, font=\tiny] {$C$};
        \draw[very thick] (3.9,1.1) -- (4.9,1.1);
        \draw[thick] (4.4,1.1) -- (4.4,0) -- (0,0);
        % Flecha de corriente
        \draw[->, thick, black!75] (2.2, 3.0) -- node[above, font=\tiny] {$I(t) = \dot{Q}$} (3.0, 3.0);
        \node[font=\footnotesize, align=center] at (2.2, -0.6) {\textbf{(a)} Circuito RC en serie};
    \end{scope}


    % Línea de fase dot{Q} vs Q (panel central)
    \begin{scope}[shift={(0.2, 0)}]
        \draw[->, thick, black!60] (-0.3,0) -- (3.6,0) node[right, font=\tiny] {$Q$};
        \draw[->, thick, black!60] (0,-1.5) -- (0,2.5) node[above, font=\tiny] {$\dot{Q}$};
        \draw[thick, black!85] (-0.2, 2.2) -- (3.2, -1.2) node[below right, font=\tiny] {$\dot{Q} = \frac{V_0}{R} - \frac{Q}{RC}$};
        \filldraw[black] (2.0, 0) circle (2pt) node[below left=2pt, font=\tiny] {$Q^* = CV_0$};
        \draw[dotted, thick, black!50] (0, 2.0) node[left, font=\tiny] {$V_0/R$} -- (0,2.0);
        \draw[->, very thick, black!90] (0.5, 0) -- (1.5, 0);
        \draw[->, very thick, black!90] (3.2, 0) -- (2.5, 0);
        \node[font=\footnotesize, align=center] at (1.8, -0.6) {\textbf{(b)} Campo $\dot{Q}(Q)$};
    \end{scope}

    % Trayectoria Q(t) vs t (panel derecho)
    \begin{scope}[shift={(5.2, 0)}]
        \draw[->, thick, black!60] (0,0) -- (3.6,0) node[right, font=\tiny] {$t$};
        \draw[->, thick, black!60] (0,0) -- (0,2.5) node[above, font=\tiny] {$Q(t)$};
        \draw[dashed, black!50, thick] (0, 1.8) -- (3.2, 1.8) node[right, font=\tiny] {$Q^* = CV_0$};
        \draw[very thick, black!85, domain=0:3.2, samples=50] plot (\x, {1.8*(1 - exp(-\x))});
        \draw[thick, dashed, black!70, domain=0:3.2, samples=50] plot (\x, {1.8 + 0.6*exp(-\x)});
        \node[font=\footnotesize, align=center] at (1.8, -0.6) {\textbf{(c)} Solución $Q(t)$};
    \end{scope}
\end{tikzpicture}
\caption{Dinámica del circuito RC. (a) Diagrama esquemático con fuente continua $V_0$. (b) Campo vectorial lineal con pendiente negativa $-1/(RC)$ y atractor estable en $Q^* = CV_0$. (c) Convergencia monótona de la carga $Q(t)$ hacia la asíntota de equilibrio.}
\label{fig:circuito_rc}
\end{figure}

Por la ley de Kirchhoff de voltajes, la suma de caídas de potencial en la malla es:
\begin{equation}
    V_0 = V_R + V_C = R I + \frac{Q}{C} = R \dot{Q} + \frac{Q}{C}
\end{equation}
Despejando la derivada temporal de la carga:
\begin{equation}
    \dot{Q} = f(Q) = \frac{V_0}{R} - \frac{Q}{RC}
    \label{eq:circuito_rc_ode}
\end{equation}

El punto de equilibrio $\dot{Q} = 0$ corresponde a $Q^* = C V_0$. Como $f'(Q) = -1/(RC) < 0$, la velocidad disminuye linealmente a medida que el capacitor se carga, garantizando que $Q^* = CV_0$ es un punto fijo global y asintóticamente estable al que el sistema converge como $Q(t) = C V_0 (1 - e^{-t/(RC)})$.

\subsection{Raíces transcendentes: $\dot{x} = x - \cos x$}

Consideremos el flujo no lineal:
\begin{equation}
    \dot{x} = x - \cos x
    \label{eq:flujo_x_cosx}
\end{equation}

Determinar los puntos fijos analíticamente requiere resolver la ecuación trascendente $x^* = \cos x^*$, la cual no admite despeje explícito en términos de funciones algebraicas estándar. Sin embargo, el método geométrico resuelve el problema de inmediato: graficamos simultáneamente en el mismo plano las curvas elementales $y = x$ e $y = \cos x$ (véase la \cref{fig:raiz_transcendente_cos}).

\begin{figure}[htbp]
\centering
\begin{tikzpicture}[>=Stealth, scale=0.9]
    \draw[->, thick, black!60] (-3.2,0) -- (3.5,0) node[right, font=\small] {$x$};
    \draw[->, thick, black!60] (0,-1.8) -- (0,2.5) node[above, font=\small] {$y$};
    
    % Curva cos(x)
    \draw[very thick, black!85, domain=-3.0:3.2, samples=80] plot (\x, {cos(\x r)}) node[right, font=\footnotesize] {$y = \cos x$};
    % Recta y = x
    \draw[very thick, black!60, domain=-1.6:2.4] plot (\x, {\x}) node[above right, font=\footnotesize] {$y = x$};
    
    % Punto de intersección x* = 0.739
    \draw[dotted, thick, black!50] (0.739, 0) node[below=3pt, font=\footnotesize] {$x^* \approx 0.739$} -- (0.739, 0.739) -- (0, 0.739);
    \draw[black, fill=white, thick] (0.739, 0) circle (2.5pt);
    \filldraw[black] (0.739, 0.739) circle (2pt);
    
    % Flechas en la línea de fase sobre el eje x
    \draw[->, very thick, black!90] (0.3, 0) -- (-0.6, 0);
    \draw[->, very thick, black!90] (-1.2, 0) -- (-2.0, 0);
    \draw[->, very thick, black!90] (1.2, 0) -- (2.0, 0);
    \draw[->, very thick, black!90] (2.2, 0) -- (3.0, 0);
    
    \node[font=\footnotesize, black!75] at (2.2, -1.0) {$x > \cos x \implies \dot{x} > 0$ ($\rightarrow$)};
    \node[font=\footnotesize, black!75] at (-1.8, 1.2) {$x < \cos x \implies \dot{x} < 0$ ($\leftarrow$)};
\end{tikzpicture}
\caption{Determinación geométrica del punto fijo y estabilidad para $\dot{x} = x - \cos x$. La recta $y = x$ y la curva $y = \cos x$ se cortan en un único punto $x^* \approx 0.739085$. Como $x - \cos x > 0$ a la derecha y negativo a la izquierda, $x^*$ es un equilibrio inestable.}
\label{fig:raiz_transcendente_cos}
\end{figure}

El análisis de signos se deduce por simple inspección visual:
\begin{itemize}
    \item Para $x > x^*$, la recta $y=x$ se encuentra por encima del coseno, luego $x - \cos x > 0 \implies \dot{x} > 0$ (flujo hacia la derecha).
    \item Para $x < x^*$, la recta se encuentra por debajo del coseno, luego $x - \cos x < 0 \implies \dot{x} < 0$ (flujo hacia la izquierda).
\end{itemize}
Por lo tanto, concluimos rigurosamente que $x^*$ es un \textbf{punto fijo inestable}, habiendo obtenido la dinámica completa sin necesidad de aproximaciones numéricas previas.

\subsection{Dinámica de poblaciones: Malthus vs. logístico}

\paragraph{El modelo de Malthus (crecimiento exponencial).}
El modelo poblacional más simple postula que la tasa de nacimientos y muertes es directamente proporcional al tamaño actual de la población $N(t)$:
\begin{equation}
    \dot{N} = r N, \quad (r > 0)
    \label{eq:malthus}
\end{equation}
donde $r$ es la tasa intrínseca de crecimiento per cápita:
\begin{equation}
    \frac{\dot{N}}{N} = r = \text{constante}
\end{equation}
La solución analítica es la explosión exponencial $N(t) = N_0 e^{rt}$. En este modelo, el origen $N^* = 0$ es un punto fijo inestable: cualquier población inicial positiva $N_0 > 0$ crece sin límite, lo cual resulta biológicamente inviable a largo plazo debido al agotamiento de recursos.

\paragraph{El modelo logístico (Verhulst).}
En un entorno con recursos finitos (alimento, espacio, refugio), los individuos compiten entre sí (\emph{crowding}). Pierre-François Verhulst (1838) propuso que la tasa de crecimiento per cápita disminuye linealmente con la densidad de población:
\begin{equation}
    \frac{\dot{N}}{N} = r \left(1 - \frac{N}{K}\right)
\end{equation}
donde $K > 0$ es la \textbf{capacidad de carga} (\emph{carrying capacity}) impuesta por el ecosistema. Multiplicando por $N$, obtenemos la célebre ecuación logística:
\begin{equation}
    \dot{N} = f(N) = r N \left(1 - \frac{N}{K}\right)
    \label{eq:logistico}
\end{equation}

\begin{figure}[htbp]
\centering
\begin{tikzpicture}[>=Stealth, scale=0.85]
    % Panel a: Tasa per cápita
    \begin{scope}[shift={(-6.2, 0)}]
        \draw[->, thick, black!60] (0,0) -- (3.5,0) node[right, font=\tiny] {$N$};
        \draw[->, thick, black!60] (0,0) -- (0,2.5) node[above, font=\tiny] {$\dot{N}/N$};
        \draw[dashed, black!50, thick] (0, 1.8) node[left, font=\tiny] {$r$} -- (3.2, 1.8) node[right, font=\tiny] {Malthus};
        \draw[very thick, black!85] (0, 1.8) -- (2.5, 0) node[below=2pt, font=\tiny] {$K$};
        \node[font=\footnotesize, align=center] at (1.6, -0.7) {\textbf{(a)} Tasa per cápita $\dot{N}/N$};
    \end{scope}

    % Panel b: Parábola dot{N} vs N
    \begin{scope}[shift={(0.2, 0)}]
        \draw[->, thick, black!60] (-0.5,0) -- (3.6,0) node[right, font=\tiny] {$N$};
        \draw[->, thick, black!60] (0,-0.8) -- (0,2.5) node[above, font=\tiny] {$\dot{N}$};
        \draw[very thick, black!85, domain=0:2.6, samples=50] plot (\x, {-2.5*\x*(\x - 2.2)/1.21});
        \draw[black, fill=white, thick] (0,0) circle (2pt) node[below left=1pt, font=\tiny] {$0$};
        \filldraw[black] (2.2,0) circle (2pt) node[below=2pt, font=\tiny] {$K$};
        \draw[->, very thick, black!90] (0.5, 0) -- (1.4, 0);
        \draw[->, very thick, black!90] (3.0, 0) -- (2.5, 0);
        \node[font=\footnotesize, align=center] at (1.5, -0.7) {\textbf{(b)} Campo $\dot{N}(N)$};
    \end{scope}

    % Panel c: Trayectorias sigmoides
    \begin{scope}[shift={(5.5, 0)}]
        \draw[->, thick, black!60] (0,0) -- (3.5,0) node[right, font=\tiny] {$t$};
        \draw[->, thick, black!60] (0,0) -- (0,2.5) node[above, font=\tiny] {$N(t)$};
        \draw[dashed, black!50, thick] (0, 1.8) -- (3.2, 1.8) node[right, font=\tiny] {$K$};
        \draw[very thick, black!85, domain=0:3.2, samples=50] plot (\x, {1.8/(1 + 5*exp(-2.0*\x))});
        \draw[thick, dashed, black!70, domain=0:3.2, samples=50] plot (\x, {1.8/(1 - 0.4*exp(-2.0*\x))});
        \node[font=\footnotesize, align=center] at (1.6, -0.7) {\textbf{(c)} Trayectorias $N(t)$};
    \end{scope}
\end{tikzpicture}
\caption{Comparación entre el crecimiento malthusiano y el logístico. (a) La tasa per cápita decrece linealmente hasta anularse en $N=K$. (b) Parábola invertida con punto inestable en $N^*=0$ y atractor estable en la capacidad de carga $N^*=K$. (c) Curvas sigmoides que convergen asintóticamente a $K$.}
\label{fig:malthus_vs_logistico}
\end{figure}

La gráfica de $f(N)$ es una parábola invertida con raíces en $N_1^* = 0$ y $N_2^* = K$. Para cualquier población inicial biológicamente realista $N_0 > 0$, el flujo conduce monótonamente a la población hacia el atractor $N^* = K$.

\paragraph{Limitaciones biológicas del modelo logístico.}
Es crucial comprender que el modelo logístico es una metáfora simplificada de la ecología de poblaciones. Como señala Charles J. Krebs en su tratado clásico de ecología \cite{krebs1972}, el crecimiento sigmoide logístico reproduce fielmente experimentos de laboratorio con organismos unicelulares (bacterias y levaduras en cajas Petri de medio controlado). Sin embargo, en poblaciones naturales de organismos multicelulares (como moscas de la fruta \emph{Drosophila} o escarabajos de la harina \emph{Tribolium}), el modelo suele fallar al predecir oscilaciones persistentes debidas a dos factores que la EDO \eqref{eq:logistico} ignora:
\begin{enumerate}
    \item \textbf{Estructura de edad:} Los organismos atraviesan fases fisiológicas diferenciadas ($\text{huevos} \to \text{larvas} \to \text{pupas} \to \text{adultos}$), donde solo ciertas clases de edad consumen recursos y se reproducen.
    \item \textbf{Retardos temporales (\emph{time delays}):} El efecto del hacinamiento no es instantáneo. Una puesta masiva de huevos no satura inmediatamente los recursos el día de hoy, pero genera una catástrofe de hambruna semanas después cuando eclosionan larvas hambrientas.
\end{enumerate}


% ============================================================
% FUENTES: draft p.17--20.
% ============================================================
\section{Análisis de estabilidad lineal}

Aunque el trazado gráfico del campo vectorial permite clasificar cualitativamente los puntos fijos, en muchas situaciones (especialmente en dimensiones superiores o con parámetros libres) se requiere un método algebraico cuantitativo. Este método es el \textbf{análisis de estabilidad lineal}.

\subsection{Derivación: linealización alrededor de $x^{*}$}

Sea $\dot{x} = f(x)$ un flujo unidimensional y sea $x^*$ un punto fijo, de modo que $f(x^*) = 0$. Consideremos una perturbación infinitesimal $\eta(t)$ alrededor del equilibrio:
\begin{equation}
    x(t) = x^* + \eta(t), \quad |\eta(t)| \ll 1
    \label{eq:perturbacion_def}
\end{equation}

Para determinar si la perturbación crece o decae con el tiempo, derivamos $\eta(t)$ respecto a $t$:
\begin{equation}
    \dot{\eta} = \frac{d}{dt}(x(t) - x^*) = \dot{x} = f(x(t)) = f(x^* + \eta)
\end{equation}
Expandiendo $f(x^* + \eta)$ en serie de Taylor alrededor de $x^*$:
\begin{equation}
    f(x^* + \eta) = f(x^*) + f'(x^*)\eta + \frac{f''(x^*)}{2!}\eta^2 + \mathcal{O}(\eta^3)
\end{equation}
Dado que $f(x^*) = 0$ por definición de punto fijo, obtenemos:
\begin{equation}
    \dot{\eta} = f'(x^*)\eta + \mathcal{O}(\eta^2)
\end{equation}
Si $f'(x^*) \neq 0$ y nos mantenemos suficientemente cerca de $x^*$ tal que $|\eta| \ll 1$, los términos de orden superior $\mathcal{O}(\eta^2)$ son completamente despreciables frente al término lineal. Obtenemos así la \textbf{ecuación linealizada}:
\begin{equation}
    \dot{\eta} \approx f'(x^*)\eta
    \label{eq:linealizada_eta}
\end{equation}
cuya solución exacta es:
\begin{equation}
    \eta(t) = \eta(0) e^{f'(x^*) t}
    \label{eq:sol_linealizada}
\end{equation}

\begin{theorem}[Criterio de estabilidad lineal en 1D]
    Sea $x^*$ un punto fijo de $\dot{x} = f(x)$ con $f \in \mathcal{C}^1$:
    \begin{itemize}
        \item Si $f'(x^*) < 0$, las perturbaciones decaen exponencialmente a cero ($\eta(t) \to 0$); el punto fijo $x^*$ es \textbf{linealmente estable} (atractor asintótico).
        \item Si $f'(x^*) > 0$, las perturbaciones crecen exponencialmente ($\eta(t) \to \pm\infty$); el punto fijo $x^*$ es \textbf{linealmente inestable} (repulsor).
        \item Si $f'(x^*) = 0$, el término lineal se anula y la linealización es \textbf{inconclusa} (punto fijo degenerado o no hiperbólico). La estabilidad queda determinada por los términos no lineales de orden superior.
    \end{itemize}
\end{theorem}

\subsection{Escala de tiempo característica $\tau = 1/|f'(x^{*})|$}

El valor numérico de la derivada $f'(x^*)$ no solo determina el signo de la estabilidad, sino que cuantifica la tasa de respuesta temporal del sistema en las inmediaciones del equilibrio.

\begin{definition}[Escala de tiempo característica]
    Para un punto fijo hiperbólico ($f'(x^*) \neq 0$), definimos el \textbf{tiempo característico de relajación} $\tau$ como:
    \begin{equation}
        \tau = \frac{1}{|f'(x^*)|}
        \label{eq:tiempo_caracteristico}
    \end{equation}
    $\tau$ representa el tiempo necesario para que una perturbación inicial $\eta(0)$ se reduzca (o se amplifique) en un factor de $1/e \approx 36.8\%$.
\end{definition}

Un valor grande de $|f'(x^*)|$ implica una respuesta dinámica rápida ($\tau$ muy pequeño), mientras que $|f'(x^*)| \approx 0$ denota una relajación sumamente lenta hacia el equilibrio.

\subsection{Ejemplos: $\sin x$ y el modelo logístico}

\begin{example}[Campo trigonométrico $\dot{x} = \sin x$]
    Los puntos fijos satisfacen $\sin x_k^* = 0 \implies x_k^* = k\pi$, con $k \in \mathbb{Z}$. Evaluando la derivada $f'(x) = \cos x$:
    \begin{equation}
        f'(x_k^*) = \cos(k\pi) = (-1)^k = \begin{cases} 
            +1 > 0 & \text{si } k \text{ es par } (0, \pm 2\pi, \dots) \implies \text{\textbf{Inestable}} \\
            -1 < 0 & \text{si } k \text{ es impar } (\pm \pi, \pm 3\pi, \dots) \implies \text{\textbf{Estable}}
        \end{cases}
    \end{equation}
    En ambos casos, la escala de tiempo característica es $\tau = 1/| \pm 1 | = 1$.
\end{example}

\begin{example}[Modelo logístico $\dot{N} = rN(1 - N/K)$]
    Reescribiendo $f(N) = rN - \frac{rN^2}{K}$, la derivada respecto a $N$ es:
    \begin{equation}
        f'(N) = r - \frac{2rN}{K}
    \end{equation}
    Evaluando en los dos equilibrios $N_1^* = 0$ y $N_2^* = K$:
    \begin{itemize}
        \item En $N_1^* = 0$: $f'(0) = r > 0 \implies$ \textbf{Inestable}, con tiempo característico $\tau = 1/r$.
        \item En $N_2^* = K$: $f'(K) = r - \frac{2rK}{K} = -r < 0 \implies$ \textbf{Estable}, con idéntico tiempo característico $\tau = 1/r$.
    \end{itemize}
\end{example}

\begin{figure}[htbp]
\centering
\begin{tikzpicture}[>=Stealth, scale=0.85]
    % Panel izquierdo: sin(x)
    \begin{scope}[shift={(-4.5, 0)}]
        \draw[->, thick, black!60] (-0.5,0) -- (4.2,0) node[right, font=\tiny] {$x$};
        \draw[->, thick, black!60] (0,-1.5) -- (0,1.8) node[above, font=\tiny] {$f(x)$};
        \draw[very thick, black!85, domain=0:3.8, samples=50] plot (\x, {1.2*sin(\x r)});
        
        \draw[black, fill=white, thick] (0,0) circle (2pt) node[below left=2pt, font=\tiny] {$0$ (inestable)};
        \filldraw[black] (3.1416,0) circle (2pt) node[below=2pt, font=\tiny] {$\pi$ (estable)};
        
        % Tangentes
        \draw[dashed, black!60] (-0.6,-0.6) -- (0.6,0.6) node[above right, font=\tiny] {$f'(0)=+1$};
        \draw[dashed, black!60] (2.5,0.6) -- (3.7,-0.6) node[below right, font=\tiny] {$f'(\pi)=-1$};
        \node[font=\footnotesize, align=center] at (1.8, -1.2) {\textbf{(a)} $\dot{x} = \sin x$};
    \end{scope}

    % Panel derecho: Logístico
    \begin{scope}[shift={(4.0, 0)}]
        \draw[->, thick, black!60] (-0.5,0) -- (4.2,0) node[right, font=\tiny] {$N$};
        \draw[->, thick, black!60] (0,-1.5) -- (0,1.8) node[above, font=\tiny] {$f(N)$};
        \draw[very thick, black!85, domain=0:3.6, samples=50] plot (\x, {-1.5*\x*(\x - 3.0)/2.25});
        
        \draw[black, fill=white, thick] (0,0) circle (2pt) node[below left=2pt, font=\tiny] {$N^*=0$};
        \filldraw[black] (3.0,0) circle (2pt) node[below=2pt, font=\tiny] {$N^*=K$};
        
        % Tangentes
        \draw[dashed, black!60] (-0.5,-0.6) -- (0.8,0.96) node[above left, font=\tiny] {$f'(0)=r$};
        \draw[dashed, black!60] (2.2,0.96) -- (3.5,-0.6) node[below right, font=\tiny] {$f'(K)=-r$};
        \node[font=\footnotesize, align=center] at (1.8, -1.2) {\textbf{(b)} $\dot{N} = rN(1 - N/K)$};
    \end{scope}
\end{tikzpicture}
\caption{Linealización geométrica en $\dot{x} = \sin x$ y en el modelo logístico. Las líneas discontinuas representan las tangentes locales cuyas pendientes $f'(x^*)$ determinan la estabilidad.}
\label{fig:estabilidad_lineal_ejemplos}
\end{figure}

\subsection{Puntos fijos degenerados y no existencia de oscilaciones en 1D}

\paragraph{Puntos fijos de orden superior ($f'(x^*) = 0$).}
Cuando $f'(x^*) = 0$, la aproximación lineal desaparece y el destino de las trayectorias depende de la primera derivada no nula en la serie de Taylor:
\begin{enumerate}
    \item $\dot{x} = -x^3$: Aquí $f'(0) = 0$, pero $f(x) > 0$ para $x < 0$ y $f(x) < 0$ para $x > 0$. El origen es \textbf{asintóticamente estable}, aunque la convergencia no es exponencial sino algebraica y lenta: integrando $\int -x^{-3}dx = \int dt \implies x(t) = \frac{x_0}{\sqrt{1 + 2x_0^2 t}} \sim t^{-1/2}$.
    \item $\dot{x} = x^3$: $f'(0) = 0$, pero $f(x)$ expulsa las trayectorias. El origen es \textbf{inestable} (con escape a infinito en tiempo finito).
    \item $\dot{x} = x^2$: $f'(0) = 0$, pero $f(x) \ge 0$ para todo $x$. El origen es un \textbf{nodo semi-estable} (atrae desde la izquierda, repele hacia la derecha).
\end{enumerate}

\begin{figure}[htbp]
\centering
\begin{tikzpicture}[>=Stealth, scale=0.8]
    % Panel a: -x^3
    \begin{scope}[shift={(-5.5, 0)}]
        \draw[->, thick, black!60] (-1.8,0) -- (1.8,0) node[right, font=\tiny] {$x$};
        \draw[->, thick, black!60] (0,-1.5) -- (0,1.5) node[above, font=\tiny] {$\dot{x}$};
        \draw[very thick, black!85, domain=-1.2:1.2, samples=50] plot (\x, {-0.9*\x*\x*\x});
        \filldraw[black] (0,0) circle (2pt);
        \draw[->, thick, black!90] (-1.2,0) -- (-0.4,0);
        \draw[->, thick, black!90] (1.2,0) -- (0.4,0);
        \node[font=\footnotesize, align=center] at (0,-1.8) {\textbf{(a)} $\dot{x} = -x^3$\\Estable};
    \end{scope}

    % Panel b: x^3
    \begin{scope}[shift={(0, 0)}]
        \draw[->, thick, black!60] (-1.8,0) -- (1.8,0) node[right, font=\tiny] {$x$};
        \draw[->, thick, black!60] (0,-1.5) -- (0,1.5) node[above, font=\tiny] {$\dot{x}$};
        \draw[very thick, black!85, domain=-1.2:1.2, samples=50] plot (\x, {0.9*\x*\x*\x});
        \draw[black, fill=white, thick] (0,0) circle (2pt);
        \draw[->, thick, black!90] (-0.4,0) -- (-1.2,0);
        \draw[->, thick, black!90] (0.4,0) -- (1.2,0);
        \node[font=\footnotesize, align=center] at (0,-1.8) {\textbf{(b)} $\dot{x} = x^3$\\Inestable};
    \end{scope}

    % Panel c: x^2
    \begin{scope}[shift={(5.5, 0)}]
        \draw[->, thick, black!60] (-1.8,0) -- (1.8,0) node[right, font=\tiny] {$x$};
        \draw[->, thick, black!60] (0,-1.5) -- (0,1.5) node[above, font=\tiny] {$\dot{x}$};
        \draw[very thick, black!85, domain=-1.2:1.2, samples=50] plot (\x, {1.0*\x*\x});
        \filldraw[black!60] (0,0) circle (2pt);
        \draw[->, thick, black!90] (-1.2,0) -- (-0.4,0);
        \draw[->, thick, black!90] (0.4,0) -- (1.2,0);
        \node[font=\footnotesize, align=center] at (0,-1.8) {\textbf{(c)} $\dot{x} = x^2$\\Semi-estable};
    \end{scope}
\end{tikzpicture}
\caption{Comportamiento de sistemas no hiperbólicos con $f'(0) = 0$. (a) $\dot{x} = -x^3$ es estable. (b) $\dot{x} = x^3$ es inestable. (c) $\dot{x} = x^2$ es un nodo semi-estable.}
\label{fig:puntos_fijos_orden_superior}
\end{figure}

\paragraph{Imposibilidad topológica de oscilaciones en 1D.}
Un resultado cualitativo de primer orden es el siguiente:
\begin{theorem}[No existencia de oscilaciones en la recta real]
    En un flujo autónomo unidimensional $\dot{x} = f(x)$ con $f \in \mathcal{C}^1(\mathbb{R})$, las soluciones $x(t)$ son funciones estrictamente monótonas en el tiempo (o estacionarias). Por tanto, \textbf{no pueden existir oscilaciones periódicas ni oscilaciones amortiguadas con sobreimpulso (\emph{overshoot})}.
\end{theorem}

\begin{figure}[htbp]
\centering
\begin{tikzpicture}[>=Stealth, scale=0.85]
    \draw[->, thick, black!60] (0,0) -- (6,0) node[right, font=\small] {$t$};
    \draw[->, thick, black!60] (0,-1.8) -- (0,2.2) node[above, font=\small] {$x(t)$};
    
    % Onda oscilatoria amortiguada
    \draw[very thick, black!40, domain=0:5.5, samples=100] plot (\x, {1.8*exp(-0.6*\x)*cos(4*\x r)});
    
    % Gran cruz roja/negra de prohibido
    \draw[very thick, black!85] (1.5, -1.2) -- (4.5, 1.5);
    \draw[very thick, black!85] (1.5, 1.5) -- (4.5, -1.2);
    
    \node[font=\bfseries\small, align=center] at (3.0, -2.2) {¡IMPOSIBLE EN FLUJOS 1D!\\El flujo sobre $\mathbb{R}$ no puede invertir su marcha};
\end{tikzpicture}
\caption{Esquema que ilustra la imposibilidad de oscilaciones en flujos continuos sobre $\mathbb{R}$. Para invertir el sentido de movimiento ($\dot{x}$ cambiando de signo), la trayectoria debería cruzar un punto con $\dot{x} = 0$, donde quedaría detenida para siempre por unicidad de soluciones.}
\label{fig:imposibilidad_oscilaciones}
\end{figure}

La razón es geométrica: estamos confinados a una línea recta. Para que $x(t)$ oscile, la velocidad $\dot{x}$ tendría que cambiar de signo, pasando forzosamente por un punto donde $\dot{x} = 0$ (un punto fijo). Pero por el teorema de unicidad de Picard--Lindelöf, una trayectoria que llega a un punto fijo requiere tiempo infinito para alcanzarlo o queda atrapada en él; jamás puede cruzarlo ni invertir su sentido\aside{¿Podría haber oscilaciones periódicas si el espacio de fases fuera una circunferencia $S^1$ en vez de la recta $\mathbb{R}$? ¡Sí! En el círculo, una partícula puede moverse monótonamente en una sola dirección ($\dot{\theta} = \omega > 0$) y completar ciclos indefinidamente sin retroceder jamás. Pero en la recta $\mathbb{R}$, las soluciones solo pueden converger monótonamente hacia un atractor o escapar hacia $\pm\infty$.}.


% ============================================================
% FUENTES: draft p.21; "Notas Biomate 1.pdf" (español).
% ============================================================
\section{Potenciales y la analogía mecánica}

Una de las formas más intuitivas y poderosas de visualizar la dinámica de un flujo unidimensional es recurrir al concepto físico de \textbf{energía potencial}.

\subsection{Sistemas sobreamortiguados: la bola en miel}

Consideremos una partícula de masa $m$ acoplada a un resorte y sometida a una fuerza externa no lineal $F(x)$ dentro de un medio extremadamente viscoso con coeficiente de arrastre $b > 0$ (imaginemos una canica moviéndose dentro de un frasco de miel o glicerina espesa, como en la \cref{fig:masa_en_miel}).

\begin{figure}[htbp]
\centering
\begin{tikzpicture}[>=Stealth, scale=0.9]
    % Recipiente con miel
    \draw[thick, fill=black!8] (-3.5,-1.2) rectangle (3.5,1.2);
    \node[font=\bfseries\footnotesize, black!40] at (0, -0.8) {MEDIO ALTAMENTE VISCOSO (MIEL)};
    
    % Pared de soporte
    \draw[very thick] (-3.5,-1.5) -- (-3.5,1.5);
    \foreach \y in {-1.4, -1.0, -0.6, -0.2, 0.2, 0.6, 1.0, 1.4} {
        \draw[thick] (-3.5, \y) -- (-3.9, \y-0.2);
    }
    
    % Resorte
    \draw[thick, decoration={aspect=0.4, segment length=3mm, amplitude=3mm, coil}, decorate] (-3.5,0) -- (-0.8,0);
    
    % Masa m
    \draw[thick, fill=black!20] (-0.8,-0.5) rectangle (0.8,0.5) node[pos=0.5, font=\small\bfseries] {$m$};
    
    % Fuerzas
    \draw[->, very thick, black!90] (0.8,0) -- (2.2,0) node[right, font=\footnotesize] {$F(x)$};
    \draw[->, thick, black!70] (-0.8,-0.2) -- (-2.0,-0.2) node[below, font=\tiny] {Arrastre $-b\dot{x}$};
    
    % Coordenada x
    \draw[->, thick, black!60] (-3.5,-1.6) -- (3.5,-1.6) node[right, font=\tiny] {$x$};
    \draw[dotted, thick] (0, 0.5) -- (0, -1.6) node[below, font=\tiny] {$x=0$};
\end{tikzpicture}
\caption{Analogía mecánica del sistema sobreamortiguado (\emph{overdamped}): una masa $m$ inmersa en un fluido de altísima viscosidad donde la fuerza de amortiguamiento domina por completo sobre la inercia ($b\dot{x} \gg m\ddot{x}$).}
\label{fig:masa_en_miel}
\end{figure}

La segunda ley de Newton establece:
\begin{equation}
    m \ddot{x} + b \dot{x} = F(x)
\end{equation}
En el régimen de \textbf{sobreamortiguamiento límite} (\emph{overdamped limit} o número de Reynolds extremadamente bajo), la inercia $m\ddot{x}$ es insignificante comparada con las fuerzas viscosas ($m\ddot{x} \ll b\dot{x}$). Despreciando el término acelerativo, obtenemos:
\begin{equation}
    b \dot{x} \approx F(x) \implies \dot{x} = \frac{1}{b}F(x) \equiv f(x)
    \label{eq:sobreamortiguado_ode}
\end{equation}
En este régimen, la velocidad de la partícula es instantáneamente proporcional a la fuerza aplicada: la partícula carece de ``memoria cinética'' y se detiene en el momento exacto en que la fuerza total se anula.

\subsection{El potencial $V(x)$, con $f(x) = -\,dV/dx$}

Para cualquier flujo unidimensional $\dot{x} = f(x)$, siempre es posible introducir una función escalar diferenciable $V(x)$, denominada \textbf{potencial}, definida mediante la convención estándar de la física:
\begin{equation}
    f(x) = -\frac{dV}{dx} \iff V(x) = -\int f(x) \, dx
    \label{eq:def_potencial}
\end{equation}

\paragraph{Monotonía del potencial a lo largo de trayectorias.}
Calculemos la derivada temporal de $V(x(t))$ siguiendo la regla de la cadena:
\begin{equation}
    \dot{V} = \frac{dV}{dt} = \frac{dV}{dx} \frac{dx}{dt} = \left(-\,f(x)\right) \left(f(x)\right) = -\left(\frac{dV}{dx}\right)^2 = -\left(f(x)\right)^2 \le 0
    \label{eq:derivada_potencial}
\end{equation}

\begin{theorem}[Disipación del potencial]
    A lo largo de cualquier solución no trivial de $\dot{x} = f(x)$, la función potencial $V(x(t))$ es monótonamente decreciente ($\dot{V} < 0$ siempre que $\dot{x} \neq 0$). El estado del sistema fluye siempre hacia regiones de menor potencial, deteniéndose únicamente en los puntos críticos donde $\frac{dV}{dx} = 0$.
\end{theorem}

\begin{figure}[htbp]
\centering
\begin{tikzpicture}[>=Stealth, scale=0.9]
    \draw[->, thick, black!60] (-4.5,0) -- (4.5,0) node[right, font=\small] {$x$};
    \draw[->, thick, black!60] (0,-1.0) -- (0,3.5) node[above, font=\small] {$V(x)$};
    
    % Paisaje de potencial V(x)
    \draw[very thick, black!85, domain=-3.6:3.6, samples=80] 
        plot (\x, {0.15*(\x+2.2)*(\x+2.2)*(\x-1.8)*(\x-1.8) + 0.3*(\x) + 0.8});
    
    % Puntos críticos
    \filldraw[black] (-2.2, 0.4) circle (2.5pt) node[below=4pt, font=\footnotesize] {$x_1^*$ (mínimo estable)};
    \draw[black, fill=white, thick] (-0.2, 2.3) circle (2.5pt) node[above=4pt, font=\footnotesize] {$x_2^*$ (máximo inestable)};
    \filldraw[black] (1.8, 0.9) circle (2.5pt) node[below=4pt, font=\footnotesize] {$x_3^*$ (mínimo estable)};
    
    % Bola rodando
    \shade[ball color=black!60] (-1.2, 1.6) circle (5pt);
    \draw[->, very thick, black!90] (-1.0, 1.4) -- (-1.6, 0.8);
    \node[font=\tiny, black!80] at (-0.6, 1.6) {Bola rodando cuesta abajo};
\end{tikzpicture}
\caption{Paisaje de energía potencial $V(x)$. El estado del sistema se comporta como una partícula sin inercia que rueda cuesta abajo hacia los mínimos locales de $V(x)$ (puntos fijos estables) y se aleja de los máximos locales (puntos fijos inestables).}
\label{fig:paisaje_potencial}
\end{figure}

\paragraph{Identificación de estabilidad mediante el potencial:}
\begin{itemize}
    \item \textbf{Puntos fijos estables (atractores):} Corresponden a los \textbf{mínimos locales} de $V(x)$, donde $V'(x^*) = 0$ y $V''(x^*) > 0 \implies f'(x^*) = -V''(x^*) < 0$.
    \item \textbf{Puntos fijos inestables (repulsores):} Corresponden a los \textbf{máximos locales} de $V(x)$, donde $V'(x^*) = 0$ y $V''(x^*) < 0 \implies f'(x^*) = -V''(x^*) > 0$.
    \item \textbf{Nodos semi-estables:} Corresponden a \textbf{puntos de inflexión con tangente horizontal}, donde $V'(x^*) = 0$ y $V''(x^*) = 0$.
\end{itemize}

\begin{example}[El potencial de doble pozo para $\dot{x} = x - x^3$]
    Integrando el campo $f(x) = x - x^3$:
    \begin{equation}
        V(x) = -\int (x - x^3) \, dx = -\frac{1}{2}x^2 + \frac{1}{4}x^4
        \label{eq:potencial_doble_pozo}
    \end{equation}
    (donde fijamos la constante de integración $C=0$). Los puntos críticos satisfacen $V'(x) = -x + x^3 = 0 \implies x^* = 0, \pm 1$.
    \begin{itemize}
        \item En $x^* = 0$: $V''(0) = -1 < 0$ (máximo local $\implies$ origen inestable).
        \item En $x^* = \pm 1$: $V''(\pm 1) = -1 + 3(1) = 2 > 0$ (dos mínimos locales simétricos de igual profundidad $V = -1/4 \implies$ equilibrios estables biestables).
    \end{itemize}
\end{example}

\begin{figure}[htbp]
\centering
\begin{tikzpicture}[>=Stealth, scale=0.85]
    \draw[->, thick, black!60] (-3.0,0) -- (3.0,0) node[right, font=\small] {$x$};
    \draw[->, thick, black!60] (0,-1.5) -- (0,2.5) node[above, font=\small] {$V(x)$};
    
    % Doble pozo V(x) = -0.5 x^2 + 0.25 x^4
    \draw[very thick, black!85, domain=-2.2:2.2, samples=70] plot (\x, {-1.2*\x*\x + 0.6*\x*\x*\x*\x});
    
    % Mínimos y máximo
    \draw[black, fill=white, thick] (0,0) circle (2.5pt) node[above=3pt, font=\footnotesize] {$x^*=0$};
    \filldraw[black] (-1.0, -0.6) circle (2.5pt) node[below=3pt, font=\footnotesize] {$x^*=-1$};
    \filldraw[black] (1.0, -0.6) circle (2.5pt) node[below=3pt, font=\footnotesize] {$x^*=+1$};
    
    \draw[dashed, black!50] (-1.0,0) -- (-1.0,-0.6);
    \draw[dashed, black!50] (1.0,0) -- (1.0,-0.6);
    
    % Flechas hacia los pozos
    \draw[->, thick, black!80] (-1.8, 0.8) -- (-1.3, -0.2);
    \draw[->, thick, black!80] (-0.3, 0.0) -- (-0.7, -0.4);
    \draw[->, thick, black!80] (0.3, 0.0) -- (0.7, -0.4);
    \draw[->, thick, black!80] (1.8, 0.8) -- (1.3, -0.2);
\end{tikzpicture}
\caption{Potencial simétrico de doble pozo $V(x) = -\frac{1}{2}x^2 + \frac{1}{4}x^4$. La barrera central en $x^* = 0$ separa dos cuencas de atracción que convergen a los mínimos biestables $x^* = \pm 1$.}
\label{fig:potencial_doble_pozo}
\end{figure}

Esta visión por potenciales será la herramienta clave cuando analicemos cómo se deforma la geometría del sistema al variar un parámetro en las bifurcaciones.

\subsection{Ejercicios}

\paragraph{Ejercicio 1 (mecánico; construir el potencial).} Para $\dot{x} = x - x^2$, calcula $V(x)$ con \eqref{eq:def_potencial} (fija $C=0$), localiza sus puntos críticos y clasifícalos usando el signo de $V''(x^*)$.
\aside{Pista: $V(x)=-\tfrac12 x^2+\tfrac13 x^3$; los críticos coinciden con los ceros de $f$, y $V''=-f'$.}
% SOL: V=-x^2/2 + x^3/3. Criticos x=0,1. V''(0)=-1<0 max -> inestable; V''(1)=+1>0 min -> estable.

\paragraph{Ejercicio 2 (análisis; $V$ como función de Lyapunov).} Usando el Teorema de disipación \eqref{eq:derivada_potencial}, argumenta que en el doble pozo \eqref{eq:potencial_doble_pozo} una trayectoria con $x(0)=0.5$ converge a $x^*=+1$ y no a $-1$. Calcula el descenso total de potencial $V(0.5)-V(1)$ que ``libera'' esa trayectoria.
\aside{Pista: $\dot V<0$ salvo en equilibrios; con $x_0=0.5>0$ el estado está en la cuenca del mínimo derecho. $V(x)=-\tfrac12 x^2+\tfrac14 x^4$.}
% SOL: x0=0.5 en cuenca de +1 (a la derecha de la barrera x=0) -> converge a +1. V(0.5)=-0.125+0.0156=-0.109; V(1)=-0.25; descenso = V(0.5)-V(1)=0.141.

\paragraph{Ejercicio 3 (interpretación; interruptor celular).} Interpreta el doble pozo $\dot{x}=x-x^3$ como un interruptor genético biestable con dos estados de expresión $x^*=\pm 1$. ¿Qué representa la barrera en $x=0$ y su altura $V(0)-V(\pm1)=1/4$? Explica cómo el ruido molecular podría inducir una transición espontánea entre estados.
\aside{Pista: la barrera es el umbral que separa las dos cuencas; una fluctuación mayor que la barrera empuja al sistema al otro pozo.}
% SOL: x=0 es el umbral inestable; altura 1/4 = energia para conmutar. Ruido con energia > 1/4 puede saltar la barrera y voltear el interruptor (transicion estocastica entre estados).



% ---- Parte III — Bifurcaciones (draft págs. 22--29 + Notas Biomate 1) ----
% ============================================================
% FUENTES: draft p.22--24.
% ============================================================
\section{Bifurcación saddle-node y formas normales}
\label{sec:saddlenode}

En los capítulos previos analizamos sistemas dinámicos con parámetros constantes fijos. Sin embargo, en los problemas reales de la física y la biología, los parámetros del entorno (temperatura, concentración de nutrientes, tasa de bombeo óptico, voltajes aplicados) varían continuamente. El estudio de cómo cambia cualitativamente la estructura del espacio de fases al alterar estos parámetros constituye la \textbf{teoría de bifurcaciones}.

\subsection{Idea de bifurcación; el ejemplo del pandeo}

\begin{definition}[Bifurcación]
    Una \textbf{bifurcación} es un cambio cualitativo o topológico en la estructura de las soluciones de un sistema dinámico (creación, destrucción o intercambio de estabilidad de puntos fijos y órbitas) que ocurre al cruzar un valor crítico de uno o más parámetros de control.
\end{definition}

\paragraph{Analogía del pandeo mecánico (\emph{buckling}).}
Un ejemplo intuitivo es una columna o viga elástica vertical delgada sujeta por su base y cargada axialmente en su extremo superior por una fuerza vertical $F$ (véase la \cref{fig:pandeo_mecanico}).
\begin{itemize}
    \item Para cargas bajas ($F < F_c$), la posición vertical perfectamente recta es un estado de equilibrio estable: si se perturba lateralmente la columna, oscila y regresa al centro.
    \item Al superar una carga crítica de Euler $F_c$, la posición recta pierde estabilidad; la columna pandea bruscamente (\emph{buckles}) hacia la izquierda o hacia la derecha, adoptando una nueva configuración curvada estable.
\end{itemize}

\begin{figure}[htbp]
\centering
\begin{tikzpicture}[>=Stealth, scale=0.85]
    % Viga recta (panel izquierdo)
    \begin{scope}[shift={(-3.5, 0)}]
        \draw[very thick] (-1.0, 0) -- (1.0, 0);
        \foreach \x in {-0.8, -0.4, 0, 0.4, 0.8} {
            \draw[thick] (\x, 0) -- (\x-0.2, -0.3);
        }
        \draw[ultra thick, black!85] (0, 0) -- (0, 3.2);
        \draw[->, very thick, black!90] (0, 4.2) -- node[right, font=\tiny] {$F < F_c$} (0, 3.2);
        \node[font=\footnotesize, align=center] at (0, -0.8) {Estado recto estable\\($F < F_c$)};
    \end{scope}

    % Viga pandeada (panel derecho)
    \begin{scope}[shift={(3.5, 0)}]
        \draw[very thick] (-1.0, 0) -- (1.0, 0);
        \foreach \x in {-0.8, -0.4, 0, 0.4, 0.8} {
            \draw[thick] (\x, 0) -- (\x-0.2, -0.3);
        }
        \draw[ultra thick, black!85] (0, 0) to[out=90, in=-80] (1.0, 1.8) to[out=100, in=-90] (0.2, 3.1);
        \draw[->, very thick, black!90] (0.2, 4.2) -- node[right, font=\tiny] {$F > F_c$} (0.2, 3.1);
        \node[font=\footnotesize, align=center] at (0, -0.8) {Pandeo lateral espontáneo\\($F > F_c$)};
    \end{scope}
\end{tikzpicture}
\caption{Analogía física del pandeo mecánico de una viga elástica. Al cruzar el umbral de carga axial $F_c$, la configuración recta vertical pierde estabilidad y el sistema bifurca hacia estados flexionados estables.}
\label{fig:pandeo_mecanico}
\end{figure}

\subsection{El prototipo $\dot{x} = r + x^{2}$ y la bifurcación blue-sky}

La bifurcación más elemental en una dimensión es la \textbf{bifurcación saddle-node} (o nodo-ensilladura, bifurcación tangente, o de pliegue / \emph{fold}). Su forma normal canónica es:
\begin{equation}
    \dot{x} = r + x^2
    \label{eq:saddle_node_normal}
\end{equation}
donde $r \in \mathbb{R}$ es el parámetro de bifurcación.

\paragraph{Análisis de los regímenes dinámicos:}
\begin{enumerate}
    \item \textbf{Para $r < 0$:} Existen dos puntos fijos reales en $x^* = \pm\sqrt{-r}$. Evaluando la derivada $f'(x) = 2x$:
    \begin{itemize}
        \item En $x_1^* = -\sqrt{-r}$: $f'(x_1^*) = -2\sqrt{-r} < 0 \implies$ \textbf{Estable} (atractor).
        \item En $x_2^* = +\sqrt{-r}$: $f'(x_2^*) = +2\sqrt{-r} > 0 \implies$ \textbf{Inestable} (repulsor).
    \end{itemize}
    \item \textbf{Para $r = 0$:} Los dos puntos fijos colisionan en el origen $x^* = 0$, formando un único \textbf{nodo semi-estable} no hiperbólico ($f'(0) = 0$).
    \item \textbf{Para $r > 0$:} Como $r + x^2 > 0$ para todo $x$, no existen puntos fijos reales. El campo vectorial es estrictamente positivo y las trayectorias viajan libremente hacia la derecha ($\dot{x} > 0$).
\end{enumerate}

\begin{figure}[htbp]
\centering
\begin{tikzpicture}[>=Stealth, scale=0.82]
    % Panel 1: r < 0
    \begin{scope}[shift={(-6.2, 3.2)}]
        \draw[->, thick, black!60] (-2.2,0) -- (2.2,0) node[right, font=\tiny] {$x$};
        \draw[->, thick, black!60] (0,-1.5) -- (0,1.8) node[above, font=\tiny] {$\dot{x}$};
        \draw[thick, black!85, domain=-1.6:1.6] plot (\x, {\x*\x - 1.0});
        \filldraw[black] (-1.0, 0) circle (2pt) node[above left, font=\tiny] {$-\sqrt{-r}$};
        \draw[black, fill=white, thick] (1.0, 0) circle (2pt) node[above right, font=\tiny] {$+\sqrt{-r}$};
        \draw[->, thick, black!90] (-1.8,0) -- (-1.3,0);
        \draw[->, thick, black!90] (0,0) -- (-0.6,0);
        \draw[->, thick, black!90] (1.3,0) -- (1.8,0);
        \node[font=\footnotesize, align=center] at (0,-1.8) {\textbf{(a)} $r < 0$ (2 equilibrios)};
    \end{scope}

    % Panel 2: r = 0
    \begin{scope}[shift={(0, 3.2)}]
        \draw[->, thick, black!60] (-2.2,0) -- (2.2,0) node[right, font=\tiny] {$x$};
        \draw[->, thick, black!60] (0,-1.5) -- (0,1.8) node[above, font=\tiny] {$\dot{x}$};
        \draw[thick, black!85, domain=-1.4:1.4] plot (\x, {\x*\x});
        \filldraw[black!60] (0, 0) circle (2pt) node[below=2pt, font=\tiny] {$x^*=0$};
        \draw[->, thick, black!90] (-1.4,0) -- (-0.6,0);
        \draw[->, thick, black!90] (0.5,0) -- (1.3,0);
        \node[font=\footnotesize, align=center] at (0,-1.8) {\textbf{(b)} $r = 0$ (colisión)};
    \end{scope}

    % Panel 3: r > 0
    \begin{scope}[shift={(6.2, 3.2)}]
        \draw[->, thick, black!60] (-2.2,0) -- (2.2,0) node[right, font=\tiny] {$x$};
        \draw[->, thick, black!60] (0,-1.5) -- (0,1.8) node[above, font=\tiny] {$\dot{x}$};
        \draw[thick, black!85, domain=-1.2:1.2] plot (\x, {\x*\x + 0.6});
        \draw[->, thick, black!90] (-1.2,0) -- (-0.4,0);
        \draw[->, thick, black!90] (0.2,0) -- (1.0,0);
        \node[font=\footnotesize, align=center] at (0,-1.8) {\textbf{(c)} $r > 0$ (sin equilibrios)};
    \end{scope}

    % Panel 4: Diagrama de bifurcación (r vs x*)
    \begin{scope}[shift={(0, -2.8)}]
        \draw[->, thick, black!60] (-3.5,0) -- (2.5,0) node[right, font=\small] {$r$};
        \draw[->, thick, black!60] (0,-2.2) -- (0,2.2) node[above, font=\small] {$x^*$};
        
        % Parábola r = - (x*)^2
        \draw[very thick, black!85, domain=-2.0:0, samples=50] plot ({-\x*\x}, \x); % rama estable inferior
        \draw[very thick, dashed, black!75, domain=0:2.0, samples=50] plot ({-\x*\x}, \x); % rama inestable superior
        
        \filldraw[black] (0,0) circle (2.5pt) node[above right, font=\footnotesize] {$(r_c, x^*) = (0,0)$};
        \node[below, font=\footnotesize] at (-1.8, -1.3) {Rama estable ($-\sqrt{-r}$)};
        \node[above, font=\footnotesize] at (-1.8, 1.3) {Rama inestable ($+\sqrt{-r}$)};
        \node[font=\bfseries\small] at (0, -2.6) {\textbf{(d)} Diagrama de bifurcación en el plano $(r, x^*)$};
    \end{scope}
\end{tikzpicture}
\caption{Bifurcación saddle-node en $\dot{x} = r + x^2$. Arriba: evolución del campo vectorial para $r < 0$, $r = 0$ y $r > 0$. Abajo: diagrama de bifurcación donde una rama estable (línea sólida) y una inestable (línea a trazos) nacen/mueren en la parábola $r = -(x^*)^2$.}
\label{fig:bifurcacion_saddle_node}
\end{figure}

\paragraph{La bifurcación ``blue-sky''.}
Ralph Abraham y Christopher Shaw (1988) \cite{abraham1988} bautizaron poéticamente a este fenómeno como \emph{``blue-sky bifurcation''}: al variar el parámetro $r$, de la nada absoluta (\emph{``out of the clear blue sky''}) aparece espontáneamente un par de puntos fijos (uno estable y otro inestable), o bien colisionan y se aniquilan en el aire\aside{Una forma equivalente de la saddle-node es $\dot{x} = r - x^2$, donde los equilibrios existen para $r > 0$ en $x^* = \pm\sqrt{r}$ y desaparecen para $r < 0$.}.

\subsection{Formas normales: Taylor local en $(x^{*}, r_{c})$}

¿Por qué la familia cuadrática $\dot{x} = r \pm x^2$ es tan importante? Porque es la \textbf{forma normal} representativa de \emph{toda} bifurcación saddle-node genérica.

\paragraph{Derivación rigurosa mediante serie de Taylor en dos variables.}
Sea $\dot{x} = f(x, r)$ una familia continua uniparamétrica de sistemas dinámicos que experimenta una bifurcación saddle-node en el punto $(x^*, r_c)$. Expandamos $f(x, r)$ en serie de Taylor de dos variables alrededor de $(x^*, r_c)$:
\begin{align}
    f(x, r) &= f(x^*, r_c) + \partial_x f(x^*, r_c)(x - x^*) + \partial_r f(x^*, r_c)(r - r_c) \nonumber \\
    &\quad + \frac{1}{2} \partial_{xx}^2 f(x^*, r_c)(x - x^*)^2 + \partial_{xr}^2 f(x^*, r_c)(x - x^*)(r - r_c) \nonumber \\
    &\quad + \frac{1}{2} \partial_{rr}^2 f(x^*, r_c)(r - r_c)^2 + \mathcal{O}\left((x - x^*)^3, (r - r_c)^3\right)
    \label{eq:taylor_2var}
\end{align}

Las condiciones definitorias de una bifurcación saddle-node en $(x^*, r_c)$ son:
\begin{enumerate}
    \item $f(x^*, r_c) = 0$ (existencia de punto fijo en el valor crítico).
    \item $\partial_x f(x^*, r_c) = 0$ (el punto fijo es no hiperbólico; la tangente es horizontal).
    \item $a \equiv \partial_r f(x^*, r_c) \neq 0$ (el parámetro $r$ desplaza la curva verticalmente a través del eje).
    \item $b \equiv \frac{1}{2} \partial_{xx}^2 f(x^*, r_c) \neq 0$ (la gráfica tiene curvatura cuadrática no nula, comportándose localmente como un cuenco parabólico).
\end{enumerate}

Sustituyendo estas condiciones en \eqref{eq:taylor_2var} y reteniendo los términos dominantes de menor orden:
\begin{equation}
    \dot{x} \approx a(r - r_c) + b(x - x^*)^2
    \label{eq:forma_normal_taylor}
\end{equation}
Definiendo las variables centradas y reescaladas $X = \sqrt{|b|}(x - x^*)$ y $R = a(r - r_c)$, la ecuación \eqref{eq:forma_normal_taylor} se reduce formalmente a $\dot{X} = R \pm X^2$.

\subsection{Ejemplo: $\dot{x} = r - x - e^{-x}$, $r_{c} = 1$}

Consideremos el sistema no lineal:
\begin{equation}
    \dot{x} = f(x, r) = r - x - e^{-x}
    \label{eq:ejemplo_sn_exp}
\end{equation}

Para encontrar el valor crítico de bifurcación $r_c$ y el punto fijo $x^*$:
\begin{align}
    f(x^*, r_c) = 0 &\implies r_c - x^* - e^{-x^*} = 0 \implies r_c - x^* = e^{-x^*} \label{eq:sn_cond1}\\
    \partial_x f(x^*, r_c) = 0 &\implies -1 + e^{-x^*} = 0 \implies e^{-x^*} = 1 \implies x^* = 0 \label{eq:sn_cond2}
\end{align}
Sustituyendo $x^* = 0$ en la condición \eqref{eq:sn_cond1}:
\begin{equation}
    r_c - 0 = e^0 = 1 \implies r_c = 1
\end{equation}

\begin{figure}[htbp]
\centering
\begin{tikzpicture}[>=Stealth, scale=0.9]
    \draw[->, thick, black!60] (-2.5,0) -- (3.5,0) node[right, font=\small] {$x$};
    \draw[->, thick, black!60] (0,-0.5) -- (0,3.5) node[above, font=\small] {$y$};
    
    % Curva exponencial e^{-x}
    \draw[very thick, black!85, domain=-1.1:3.0, samples=60] plot (\x, {exp(-\x)}) node[right, font=\footnotesize] {$y = e^{-x}$};
    
    % Recta r = 1 (tangente crítica)
    \draw[very thick, black!60, domain=-1.5:2.5] plot (\x, {1.0 - \x}) node[above right, font=\tiny] {$y = 1 - x \ (r_c=1)$};
    \filldraw[black] (0, 1.0) circle (2.5pt) node[above right=2pt, font=\footnotesize] {Tangencia en $(0,1)$};
    
    % Recta r = 1.6 (dos intersecciones)
    \draw[thick, dashed, black!75, domain=-1.0:2.8] plot (\x, {1.6 - \x}) node[above right, font=\tiny] {$y = 1.6 - x \ (r > 1)$};
    \filldraw[black] (-0.45, 1.57) circle (2pt);
    \draw[black, fill=white, thick] (1.1, 0.33) circle (2pt);
\end{tikzpicture}
\caption{Determinación geométrica de la bifurcación saddle-node en $\dot{x} = r - x - e^{-x}$. Para $r = r_c = 1$, la recta $y = r-x$ es exactamente tangente a la exponencial $y = e^{-x}$ en $x^* = 0$.}
\label{fig:ejemplo_transcendente_saddlenode}
\end{figure}

Expandiendo $e^{-x}$ en serie de Maclaurin: $e^{-x} = 1 - x + \frac{x^2}{2} - \frac{x^3}{6} + \cdots$:
\begin{align}
    \dot{x} &= r - x - \left(1 - x + \frac{x^2}{2} + \mathcal{O}(x^3)\right) \nonumber \\
    &= (r - 1) - \frac{1}{2}x^2 + \mathcal{O}(x^3)
\end{align}
Identificando $R = r - 1$ y reescalando $X = x/\sqrt{2}$, recuperamos $\dot{X} \approx R - X^2$, demostrando que el sistema experimenta una bifurcación saddle-node estándar en $r_c = 1$.

\subsection{Ejercicios}

\paragraph{Ejercicio 1 (mecánico; saddle-node explícita).} Para $\dot{x} = r - x^2$, halla los puntos fijos en función de $r$, clasifícalos con $f'(x^*)$ e identifica el valor crítico $r_c$ donde nacen/mueren. Dibuja el diagrama de bifurcación en el plano $(r,x^*)$.
\aside{Pista: hay equilibrios reales sólo si $r\ge 0$, en $x^*=\pm\sqrt{r}$; $f'(x)=-2x$.}
% SOL: r>0: x*=+sqrt(r) estable (f'=-2sqrt r<0), x*=-sqrt(r) inestable. r=0: colision semiestable. r<0: sin equilibrios. rc=0; parabola r=(x*)^2.

\paragraph{Ejercicio 2 (análisis; forma normal por Taylor).} Considera $\dot{x} = r - \cosh x$. Aplicando las condiciones definitorias $f=0$ y $\partial_x f=0$, localiza $(x^*,r_c)$ y, con la serie $\cosh x = 1 + x^2/2 + \mathcal{O}(x^4)$, redúcelo a la forma normal $\dot{X}\approx R - X^2$. Concluye el tipo de bifurcación.
\aside{Pista: $\partial_x f=-\sinh x=0\Rightarrow x^*=0$; luego $f(0)=r-1=0$.}
% SOL: x*=0, rc=1. xdot = (r-1) - x^2/2 -> R=r-1, X=x/sqrt2 -> R - X^2. Saddle-node en rc=1.

\paragraph{Ejercicio 3 (interpretación; punto de no retorno ecológico).} Un lago tiene un equilibrio estable de agua clara y, al aumentar la carga de nutrientes $r$, ese equilibrio colisiona con uno inestable y desaparece por una saddle-node, saltando el sistema a un estado turbio. Explica por qué la recuperación exige bajar $r$ \emph{muy por debajo} del umbral de colapso (histéresis) y por qué la metáfora ``blue-sky'' captura lo abrupto del cambio.
\aside{Pista: tras el colapso ya no hay equilibrio claro cercano; hay que recrear la rama estable, lo que ocurre a un $r$ distinto del de su destrucción si coexisten otras ramas.}
% SOL: Al desaparecer el atractor claro, el sistema cae al turbio; restaurar el claro requiere reaparicion de la rama (a r mas bajo) -> histeresis. Blue-sky: el par de equilibrios aparece/desaparece 'de la nada' al cruzar rc.


% ============================================================
% FUENTES: draft p.25--26.
% ============================================================
\section{Bifurcación transcrítica}

A diferencia de la bifurcación saddle-node, donde los puntos de equilibrio se crean o destruyen en pares, existen numerosos sistemas en la física, la química y la biología donde un punto de equilibrio específico persiste para todos los valores de los parámetros (por ejemplo, el estado de extinción $N=0$ en dinámica de poblaciones o el estado de reposo electromagnético en un láser). En tales casos, las ramas de equilibrio no desaparecen, sino que colisionan e intercambian su estabilidad. Este mecanismo se conoce como \textbf{bifurcación transcrítica}.

\subsection{Forma normal $\dot{x} = rx - x^{2}$ e intercambio de estabilidad}

La forma normal prototípica de la bifurcación transcrítica es:
\begin{equation}
    \dot{x} = rx - x^2 = x(r - x)
    \label{eq:transcritica_forma_normal}
\end{equation}
donde $r \in \mathbb{R}$ es el parámetro de control.

\paragraph{Puntos de equilibrio y estabilidad:}
Los puntos fijos se obtienen de $x(r - x) = 0$, rindiendo siempre dos soluciones:
\begin{equation}
    x_1^* = 0, \quad x_2^* = r
\end{equation}
Calculando la derivada del campo vectorial $f'(x) = r - 2x$:
\begin{itemize}
    \item En la rama trivial $x_1^* = 0$: $f'(0) = r$.
    \begin{itemize}
        \item Si $r < 0$: $f'(0) < 0 \implies x_1^* = 0$ es \textbf{Estable}.
        \item Si $r > 0$: $f'(0) > 0 \implies x_1^* = 0$ es \textbf{Inestable}.
    \end{itemize}
    \item En la rama móvil $x_2^* = r$: $f'(r) = r - 2r = -r$.
    \begin{itemize}
        \item Si $r < 0$: $f'(r) = -r > 0 \implies x_2^* = r$ es \textbf{Inestable}.
        \item Si $r > 0$: $f'(r) = -r < 0 \implies x_2^* = r$ es \textbf{Estable}.
    \end{itemize}
    \item En el valor crítico $r = 0$: Las dos ramas se cruzan en el origen $x^* = 0$, que se convierte en un nodo semi-estable $\dot{x} = -x^2$.
\end{itemize}

\begin{figure}[htbp]
\centering
\begin{tikzpicture}[>=Stealth, scale=0.82]
    % Panel 1: r < 0
    \begin{scope}[shift={(-6.2, 3.2)}]
        \draw[->, thick, black!60] (-2.2,0) -- (2.2,0) node[right, font=\tiny] {$x$};
        \draw[->, thick, black!60] (0,-1.5) -- (0,1.8) node[above, font=\tiny] {$\dot{x}$};
        \draw[thick, black!85, domain=-1.6:0.8] plot (\x, {-1.0*\x - \x*\x});
        \draw[black, fill=white, thick] (-1.0, 0) circle (2pt) node[below=2pt, font=\tiny] {$x^*=r$};
        \filldraw[black] (0, 0) circle (2pt) node[above right=1pt, font=\tiny] {$x^*=0$};
        \draw[->, thick, black!90] (-1.6,0) -- (-1.2,0);
        \draw[->, thick, black!90] (-0.3,0) -- (-0.7,0);
        \draw[->, thick, black!90] (0.8,0) -- (0.3,0);
        \node[font=\footnotesize, align=center] at (0,-1.8) {\textbf{(a)} $r < 0$ ($x=0$ estable)};
    \end{scope}

    % Panel 2: r = 0
    \begin{scope}[shift={(0, 3.2)}]
        \draw[->, thick, black!60] (-2.2,0) -- (2.2,0) node[right, font=\tiny] {$x$};
        \draw[->, thick, black!60] (0,-1.5) -- (0,1.8) node[above, font=\tiny] {$\dot{x}$};
        \draw[thick, black!85, domain=-1.2:1.2] plot (\x, {-\x*\x});
        \filldraw[black!60] (0, 0) circle (2pt) node[above=2pt, font=\tiny] {$x^*=0$};
        \draw[->, thick, black!90] (-1.2,0) -- (-0.4,0);
        \draw[->, thick, black!90] (1.2,0) -- (0.4,0);
        \node[font=\footnotesize, align=center] at (0,-1.8) {\textbf{(b)} $r = 0$ (semi-estable)};
    \end{scope}

    % Panel 3: r > 0
    \begin{scope}[shift={(6.2, 3.2)}]
        \draw[->, thick, black!60] (-2.2,0) -- (2.2,0) node[right, font=\tiny] {$x$};
        \draw[->, thick, black!60] (0,-1.5) -- (0,1.8) node[above, font=\tiny] {$\dot{x}$};
        \draw[thick, black!85, domain=-0.8:1.6] plot (\x, {1.0*\x - \x*\x});
        \draw[black, fill=white, thick] (0, 0) circle (2pt) node[below left=1pt, font=\tiny] {$x^*=0$};
        \filldraw[black] (1.0, 0) circle (2pt) node[above=2pt, font=\tiny] {$x^*=r$};
        \draw[->, thick, black!90] (-0.8,0) -- (-0.3,0);
        \draw[->, thick, black!90] (0.3,0) -- (0.7,0);
        \draw[->, thick, black!90] (1.6,0) -- (1.2,0);
        \node[font=\footnotesize, align=center] at (0,-1.8) {\textbf{(c)} $r > 0$ ($x=r$ estable)};
    \end{scope}

    % Panel 4: Diagrama de bifurcación (r vs x*)
    \begin{scope}[shift={(0, -2.8)}]
        \draw[->, thick, black!60] (-3.0,0) -- (3.0,0) node[right, font=\small] {$r$};
        \draw[->, thick, black!60] (0,-2.2) -- (0,2.2) node[above, font=\small] {$x^*$};
        
        % Rama trivial x* = 0
        \draw[very thick, black!85] (-2.5,0) -- (0,0); % estable para r < 0
        \draw[very thick, dashed, black!75] (0,0) -- (2.5,0); % inestable para r > 0
        
        % Rama móvil x* = r
        \draw[very thick, dashed, black!75] (-2.0,-2.0) -- (0,0); % inestable para r < 0
        \draw[very thick, black!85] (0,0) -- (2.0,2.0); % estable para r > 0
        
        \filldraw[black] (0,0) circle (2.5pt) node[below right=2pt, font=\footnotesize] {Bifurcación $(0,0)$};
        \node[above, font=\footnotesize] at (-1.5, 0.1) {Estable ($x^*=0$)};
        \node[below, font=\footnotesize] at (1.5, -0.1) {Inestable ($x^*=0$)};
        \node[above, font=\footnotesize] at (1.2, 1.3) {Estable ($x^*=r$)};
        \node[font=\bfseries\small] at (0, -2.6) {\textbf{(d)} Intercambio de estabilidad (\emph{exchange of stabilities})};
    \end{scope}
\end{tikzpicture}
\caption{Bifurcación transcrítica en $\dot{x} = rx - x^2$. Las ramas de equilibrio $x^*=0$ y $x^*=r$ coexisten para todo $r$ e intercambian su estabilidad al colisionar en el punto crítico $r=0$.}
\label{fig:bifurcacion_transcritica}
\end{figure}

\subsection{Ejemplo no polinomial y condición $ab = 1$}

Consideremos el siguiente flujo unidimensional no polinomial:
\begin{equation}
    \dot{x} = x(1 - x^2) - a(1 - e^{-bx})
    \label{eq:ejemplo_transcritica_nopol}
\end{equation}
con parámetros positivos $a > 0$ y $b > 0$.

Notemos en primer lugar que $x^* = 0$ es un punto de equilibrio para cualquier par $(a, b)$, ya que $0(1 - 0) - a(1 - e^0) = 0$.

\paragraph{Aproximación en serie de Taylor.}
Para analizar la dinámica cerca del origen ($|x| \ll 1$), expandimos la exponencial $e^{-bx}$ en serie de Maclaurin:
\begin{equation}
    e^{-bx} = 1 - bx + \frac{1}{2}b^2 x^2 - \frac{1}{6}b^3 x^3 + \mathcal{O}(x^4)
\end{equation}
Sustituyendo en la ecuación diferencial \eqref{eq:ejemplo_transcritica_nopol}:
\begin{align}
    \dot{x} &= x - x^3 - a\left[1 - \left(1 - bx + \frac{1}{2}b^2 x^2 - \frac{1}{6}b^3 x^3 + \mathcal{O}(x^4)\right)\right] \nonumber \\
    &= x - x^3 - a\left(bx - \frac{1}{2}b^2 x^2 + \mathcal{O}(x^3)\right) \nonumber \\
    &= (1 - ab)x + \frac{1}{2}ab^2 x^2 + \mathcal{O}(x^3)
    \label{eq:expansion_transcritica_ab}
\end{align}

\begin{figure}[htbp]
\centering
\begin{tikzpicture}[>=Stealth, scale=0.9]
    \draw[->, thick, black!60] (0,0) -- (4.5,0) node[right, font=\small] {$a$};
    \draw[->, thick, black!60] (0,0) -- (0,3.8) node[above, font=\small] {$b$};
    
    % Hipérbola ab = 1 => b = 1/a
    \draw[very thick, black!85, domain=0.3:4.0, samples=60] plot (\x, {1.0/\x}) node[above right, font=\footnotesize] {$ab = 1$ (curva de bifurcación)};
    
    \node[font=\footnotesize, black!75, align=center] at (2.8, 2.2) {Régimen $ab > 1$:\\Origen inestable ($f'(0) < 0$),\\nace rama no trivial estable $x^* > 0$};
    \node[font=\footnotesize, black!75, align=center] at (1.0, 0.4) {Régimen $ab < 1$:\\Origen $x^*=0$ estable};
    
    \filldraw[black] (1.0, 1.0) circle (2pt) node[below left, font=\tiny] {$(1,1)$};
\end{tikzpicture}
\caption{Curva de bifurcación transcrítica en el espacio de parámetros $(a, b)$. La hipérbola $ab = 1$ separa las regiones de estabilidad del estado trivial.}
\label{fig:curva_bifurcacion_ab}
\end{figure}

\paragraph{Condición de bifurcación y rama no nula.}
Comparando con la forma normal $\dot{x} = R x + c x^2$, la bifurcación transcrítica ocurre cuando el coeficiente lineal se anula:
\begin{equation}
    1 - ab = 0 \iff \boxed{ab = 1}
\end{equation}
Esta relación define una hipérbola crítica en el primer cuadrante del plano de parámetros $(a, b)$ (véase la \cref{fig:curva_bifurcacion_ab}).

Para encontrar la fórmula aproximada del punto fijo que bifurca desde $x=0$ en las inmediaciones de la curva ($ab \approx 1$):
\begin{equation}
    (1 - ab)x^* + \frac{ab^2}{2}(x^*)^2 \approx 0 \implies x^* \left[(1 - ab) + \frac{ab^2}{2}x^*\right] \approx 0
\end{equation}
Excluyendo la solución trivial $x^* = 0$, obtenemos la rama bifurcada:
\begin{equation}
    x^* \approx \frac{2(ab - 1)}{ab^2}
    \label{eq:rama_bifurcada_ab}
\end{equation}
Esta expresión es válida localmente para $|x^*| \ll 1$ y $|ab - 1| \ll 1$.


% ============================================================
% FUENTES: draft p.27--29 (Haken 1983).
% ============================================================
\section{El umbral del láser (modelo de Haken)}

Una de las aplicaciones físicas más elegantes de la bifurcación transcrítica es la explicación del funcionamiento del láser (\emph{Light Amplification by Stimulated Emission of Radiation}) formulada por Hermann Haken en el marco de la sinergética \cite{haken1983}.

\subsection{Ganancia, pérdida de cavidad y agotamiento de átomos excitados}

Un dispositivo láser básico consta de un medio activo (por ejemplo, un cristal de Nd:YAG o un tubo de gas) confinado dentro de una cavidad resonante entre dos espejos paralelos: un espejo totalmente reflectante en un extremo y un espejo parcialmente transmisor en el otro, a través del cual emerge el haz de luz (véase la \cref{fig:esquema_laser}). Una fuente de bombeo externa (una lámpara de destello o una descarga eléctrica alimentada por un banco de capacitores) excita los átomos del medio activo llevándolos a estados de mayor energía.

\begin{figure}[htbp]
\centering
\begin{tikzpicture}[>=Stealth, scale=0.85]
    % Espejo izquierdo (100% reflectante)
    \draw[very thick, fill=black!25] (-4.2, -1.2) rectangle (-3.9, 1.2);
    \node[above, font=\tiny, align=center] at (-4.05, 1.2) {Espejo 100\%\\reflectante};
    
    % Espejo derecho (parcialmente transmisor)
    \draw[very thick, fill=black!10] (3.9, -1.2) rectangle (4.2, 1.2);
    \node[above, font=\tiny, align=center] at (4.05, 1.2) {Espejo 95\%\\transmisor};
    
    % Medio activo (Cristal Nd:YAG)
    \draw[thick, fill=black!6] (-3.2, -0.9) rectangle (3.2, 0.9);
    \node[font=\footnotesize\bfseries] at (0, 0) {Medio Activo (Cristal Nd:YAG)};
    
    % Lámpara de bombeo óptico
    \draw[thick, fill=black!15] (-3.2, 1.5) rectangle (3.2, 2.0);
    \node[font=\footnotesize] at (0, 1.75) {Lámpara de bombeo óptico (\emph{Flashlamp})};
    \foreach \x in {-2.5, -1.5, -0.5, 0.5, 1.5, 2.5} {
        \draw[->, thick, black!70] (\x, 1.5) -- (\x, 0.9);
    }
    
    % Circuito de bombeo
    \draw[thick] (-3.2, 1.75) -- (-4.5, 1.75) -- (-4.5, 2.7) -- (-0.3, 2.7);
    \draw[very thick] (-0.3, 2.4) -- (-0.3, 3.0) node[above=2pt, font=\tiny] {$C$};
    \draw[very thick] (0.3, 2.4) -- (0.3, 3.0);
    \draw[thick] (0.3, 2.7) -- (4.5, 2.7) -- (4.5, 1.75) -- (3.2, 1.75);
    \node[above=8pt, font=\tiny] at (0, 3.0) {Descarga de pulso de bombeo};

    
    % Haz de salida
    \draw[->, ultra thick, black!90] (4.2, 0) -- (6.2, 0) node[right, font=\small\bfseries] {Haz Láser};
    \draw[thick, dashed, black!50] (4.2, 0.4) -- (6.0, 0.4);
    \draw[thick, dashed, black!50] (4.2, -0.4) -- (6.0, -0.4);
\end{tikzpicture}
\caption{Esquema físico de una cavidad láser de estado sólido. Un cristal activo entre dos espejos es excitado por una lámpara de destello; superado el umbral crítico de bombeo, la emisión estimulada genera un haz coherente macroscópico.}
\label{fig:esquema_laser}
\end{figure}

\paragraph{Emisión espontánea vs. estimulada y autoorganización.}
Cuando un átomo excitado decae a su estado fundamental, emite un fotón. Si el bombeo es débil, los átomos emiten de manera independiente y desordenada en direcciones y fases aleatorias (emisión espontánea, equivalente a una lámpara ordinaria). No obstante, cuando la densidad de fotones y de átomos excitados es suficientemente alta, un fotón que pasa cerca de un átomo excitado estimula la emisión de un segundo fotón idéntico en energía, dirección y fase (emisión estimulada de Einstein). Este proceso cooperativo conduce a la \textbf{sincronización y autoorganización} de los dipolos atómicos, emergiendo un campo electromagnético coherente macroscópico.

\subsection{La ecuación de cerrado $\dot{n} = (GN_{0} - k)n - \alpha G n^{2}$}

Siguiendo el modelo simplificado de Haken \cite{haken1983}, describimos el estado del sistema mediante dos variables macroscópicas:
\begin{itemize}
    \item $n(t) \ge 0$: número de fotones en el modo láser de la cavidad.
    \item $N(t) \ge 0$: número de átomos excitados en el medio activo.
\end{itemize}

El balance de fotones se rige por:
\begin{equation}
    \dot{n} = \text{ganancia} - \text{pérdidas} = G n N - k n
    \label{eq:balance_fotones}
\end{equation}
donde:
\begin{itemize}
    \item $G > 0$ es el \textbf{coeficiente de ganancia} por emisión estimulada (la probabilidad por unidad de tiempo de que un fotón interactúe con un átomo excitado).
    \item $k > 0$ es la \textbf{tasa de pérdidas} de la cavidad (escape a través del espejo semi-reflector y disipación térmica). La vida media de un fotón en la cavidad es $\tau = 1/k$.
\end{itemize}

\paragraph{Agotamiento atómico (\emph{gain saturation}).}
Tras emitir un fotón, el átomo decae a su nivel fundamental. En un pulso de bombeo, el número inicial de átomos excitados es $N_0$. A medida que la población de fotones $n(t)$ crece, los átomos excitados se agotan a una tasa proporcional a la intensidad del campo:
\begin{equation}
    N(t) = N_0 - \alpha n(t), \quad \alpha > 0
    \label{eq:saturacion_atomos}
\end{equation}
donde $\alpha$ cuantifica el factor de agotamiento de la inversión de población.

Sustituyendo \eqref{eq:saturacion_atomos} en la ecuación \eqref{eq:balance_fotones}:
\begin{align}
    \dot{n} &= G n (N_0 - \alpha n) - k n \nonumber \\
    &= (G N_0 - k)n - \alpha G n^2
    \label{eq:laser_transcritica_ode}
\end{align}
¡Esta ecuación tiene exactamente la estructura de una \textbf{bifurcación transcrítica} $\dot{n} = r n - c n^2$ con parámetro efectivo $r = G N_0 - k$ y $c = \alpha G > 0$!

\subsection{Umbral $N_{0,c} = k/G$: régimen lámpara vs. régimen láser}

Los puntos fijos del sistema satisfacen:
\begin{equation}
    n \left[(G N_0 - k) - \alpha G n\right] = 0
\end{equation}
cuyas soluciones son:
\begin{equation}
    n_1^* = 0, \quad n_2^* = \frac{G N_0 - k}{\alpha G}
\end{equation}

\begin{figure}[htbp]
\centering
\begin{tikzpicture}[>=Stealth, scale=0.82]
    % Panel 1: N_0 < k/G
    \begin{scope}[shift={(-6.2, 0)}]
        \draw[->, thick, black!60] (-0.3,0) -- (3.2,0) node[right, font=\tiny] {$n$};
        \draw[->, thick, black!60] (0,-1.5) -- (0,1.8) node[above, font=\tiny] {$\dot{n}$};
        \draw[thick, black!85, domain=0:2.5, samples=50] plot (\x, {-0.8*\x - 0.5*\x*\x});
        \filldraw[black] (0,0) circle (2pt) node[above right=1pt, font=\tiny] {$n^*=0$ (estable)};
        \draw[->, thick, black!90] (2.2,0) -- (1.2,0);
        \draw[->, thick, black!90] (1.0,0) -- (0.3,0);
        \node[font=\footnotesize, align=center] at (1.5, -1.8) {\textbf{(a)} $N_0 < k/G$\\\emph{Régimen lámpara}};
    \end{scope}

    % Panel 2: N_0 = k/G
    \begin{scope}[shift={(0, 0)}]
        \draw[->, thick, black!60] (-0.3,0) -- (3.2,0) node[right, font=\tiny] {$n$};
        \draw[->, thick, black!60] (0,-1.5) -- (0,1.8) node[above, font=\tiny] {$\dot{n}$};
        \draw[thick, black!85, domain=0:2.2, samples=50] plot (\x, {-0.7*\x*\x});
        \filldraw[black!60] (0,0) circle (2pt) node[above right=1pt, font=\tiny] {$n^*=0$};
        \draw[->, thick, black!90] (1.8,0) -- (0.8,0);
        \node[font=\footnotesize, align=center] at (1.5, -1.8) {\textbf{(b)} $N_0 = k/G$\\\emph{Umbral crítico}};
    \end{scope}

    % Panel 3: N_0 > k/G
    \begin{scope}[shift={(6.2, 0)}]
        \draw[->, thick, black!60] (-0.3,0) -- (3.2,0) node[right, font=\tiny] {$n$};
        \draw[->, thick, black!60] (0,-1.5) -- (0,1.8) node[above, font=\tiny] {$\dot{n}$};
        \draw[thick, black!85, domain=0:2.4, samples=50] plot (\x, {1.2*\x - 0.7*\x*\x});
        \draw[black, fill=white, thick] (0,0) circle (2pt) node[below left=1pt, font=\tiny] {$0$};
        \filldraw[black] (1.71,0) circle (2pt) node[above=2pt, font=\tiny] {$n_2^*$ (estable)};
        \draw[->, thick, black!90] (0.3,0) -- (1.0,0);
        \draw[->, thick, black!90] (2.4,0) -- (1.9,0);
        \node[font=\footnotesize, align=center] at (1.5, -1.8) {\textbf{(c)} $N_0 > k/G$\\\emph{Régimen láser}};
    \end{scope}
\end{tikzpicture}
\caption{Línea de fase del modelo de Haken para tres intensidades de bombeo atómico $N_0$. (a) Bajo el umbral, los fotones decaen a cero. (b) En el umbral crítico. (c) Sobre el umbral, la radiación láser coherente se estabiliza en $n_2^* > 0$.}
\label{fig:fases_laser}
\end{figure}

El umbral de encendido del láser (\emph{laser threshold}) ocurre en el valor crítico:
\begin{equation}
    G N_0 - k = 0 \iff \boxed{N_{0,c} = \frac{k}{G}}
    \label{eq:umbral_laser_critico}
\end{equation}

\paragraph{Interpretación física de los regímenes:}
\begin{enumerate}
    \item \textbf{Régimen de lámpara ($N_0 < N_{0,c}$):} La ganancia óptica no compensa las pérdidas de la cavidad ($G N_0 < k$). El único equilibrio físicamente admisible ($n \ge 0$) es $n_1^* = 0$, que es asintóticamente estable. Cualquier destello inicial de fotones se extingue exponencialmente: el dispositivo emite luz incoherente débil.
    \item \textbf{Régimen láser ($N_0 > N_{0,c}$):} La ganancia por bombeo supera las pérdidas ($G N_0 > k$). El estado apagado $n_1^* = 0$ se vuelve inestable y el sistema converge hacia el atractor positivo estable:
    \begin{equation}
        n_2^* = \frac{G N_0 - k}{\alpha G} = \frac{1}{\alpha}\left(N_0 - N_{0,c}\right)
    \end{equation}
    La intensidad del campo láser aumenta linealmente con el bombeo adicional por encima del umbral (véase la \cref{fig:diagrama_bifurcacion_laser})\aside{Existe una notable analogía matemática entre el umbral del láser y el número reproductivo básico $R_0$ en la epidemiología matemática (modelos SIR): $R_0 = \frac{G N_0}{k}$. Si $R_0 < 1$, la tasa de recuperación domina y la infección desaparece ($I^* = 0$, análogo a la lámpara); si $R_0 > 1$, la tasa de transmisión supera la recuperación y se establece un estado endémico persistente ($I^* > 0$, análogo al haz láser).}.
\end{enumerate}

\begin{figure}[htbp]
\centering
\begin{tikzpicture}[>=Stealth, scale=0.9]
    \draw[->, thick, black!60] (0,0) -- (6.0,0) node[right, font=\small] {$N_0$ (bombeo)};
    \draw[->, thick, black!60] (0,0) -- (0,3.5) node[above, font=\small] {$n^*$ (fotones)};
    
    % Umbral crítico N_{0,c} = 2.5
    \draw[dashed, black!50, thick] (2.5, 0) node[below, font=\footnotesize] {$N_{0,c} = k/G$} -- (2.5, 3.0);
    
    % Rama n* = 0
    \draw[ultra thick, black!85] (0,0) -- (2.5,0);
    \draw[thick, dashed, black!60] (2.5,0) -- (5.5,0);
    
    % Rama láser n* = (N_0 - N_c)/alpha
    \draw[ultra thick, black!85] (2.5,0) -- (5.2, 2.7);
    \draw[thick, dashed, black!60] (0.8, -1.7) -- (2.5,0);
    
    \filldraw[black] (2.5,0) circle (2.5pt) node[above left=3pt, font=\footnotesize] {Umbral láser};
    
    \node[font=\footnotesize, black!75, align=center] at (1.2, 1.2) {\textbf{Régimen Lámpara}\\$n^* = 0$ estable};
    \node[font=\footnotesize, black!75, align=center] at (4.5, 1.2) {\textbf{Régimen Láser}\\$n^* > 0$ coherente};
\end{tikzpicture}
\caption{Diagrama de bifurcación transcrítica del número de fotones en estado estacionario $n^*$ en función del parámetro de bombeo $N_0$. El encendido del láser representa una transición de fase continua de no equilibrio.}
\label{fig:diagrama_bifurcacion_laser}
\end{figure}


% ============================================================
% FUENTES: "Notas Biomate 1.pdf" (español) — NO está en el draft.
% ============================================================
\section{Bifurcación pitchfork}

La \textbf{bifurcación pitchfork} (también llamada bifurcación de diapasón o tridente) es la tercera clase fundamental de bifurcación local en una dimensión. Surge de manera natural y genérica en sistemas físicos y biológicos que poseen una simetría espacial de reflexión respecto al origen.

\subsection{Simetría y equivariancia}

Consideremos un flujo unidimensional $\dot{x} = f(x, r)$. Decimos que el sistema es simétrico ante reflexiones espaciales $x \to -x$ si el campo vectorial es una función impar respecto a $x$:
\begin{equation}
    f(-x, r) = -f(x, r)
    \label{eq:imparidad_campo}
\end{equation}
Bajo esta condición, el campo vectorial se denomina \textbf{equivariante}: si $x(t)$ es una solución de la ecuación diferencial, entonces la trayectoria reflejada $-x(t)$ también es forzosamente una solución válida.

Una consecuencia geométrica directa de la equivariancia es que el origen $x^* = 0$ es siempre un punto de equilibrio para cualquier valor del parámetro ($f(0, r) = -f(0, r) \implies f(0, r) = 0$), y cualquier nuevo punto fijo no trivial debe nacer o destruirse obligatoriamente en parejas simétricas $\pm x^*$.

Existen dos variedades fundamentales de esta bifurcación: la \emph{supercrítica} y la \emph{subcrítica}.

\subsection{Pitchfork supercrítica $\dot{x} = rx - x^{3}$; slowing down crítico}

La forma normal canónica de la bifurcación pitchfork supercrítica es:
\begin{equation}
    \dot{x} = rx - x^3 = x(r - x^2)
    \label{eq:pitchfork_supercritica}
\end{equation}
donde el término cúbico no lineal $-x^3$ ejerce un efecto estabilizador y confinante a gran escala.

\paragraph{Análisis de los regímenes dinámicos:}
\begin{enumerate}
    \item \textbf{Para $r < 0$:} El único punto fijo real es el origen $x^* = 0$. Como $f'(0) = r < 0$, el origen es un \textbf{atractor globalmente estable}; las trayectorias decaen exponencialmente hacia cero como $x(t) \sim e^{rt}$.
    \item \textbf{Para $r = 0$:} El origen $x^* = 0$ sigue siendo el único punto de equilibrio y es asintóticamente estable, pero el término lineal se anula ($f'(0) = 0$). La dinámica se rige por $\dot{x} = -x^3$, cuya solución exacta es:
    \begin{equation}
        \int \frac{dx}{x^3} = -\int dt \implies x(t) = \frac{x_0}{\sqrt{1 + 2x_0^2 t}}
    \end{equation}
    Para tiempos largos, $x(t) \sim t^{-1/2}$. La convergencia ya no es exponencial rápida sino algebraica y sumamente lenta. Este fenómeno universal se denomina \textbf{ralentización crítica o desaceleración crítica (\emph{critical slowing down})}. Cerca de la bifurcación, la escala de tiempo característica $\tau = 1/|f'(0)|$ diverge a infinito ($\tau \to \infty$), por lo que el sistema pierde capacidad de recuperarse rápidamente ante perturbaciones externas.
    \item \textbf{Para $r > 0$:} El origen $x_0^* = 0$ pierde estabilidad ($f'(0) = r > 0$, repulsor), y emergen simétricamente dos nuevos puntos fijos estables:
    \begin{equation}
        x_\pm^* = \pm\sqrt{r}
    \end{equation}
    Evaluando la estabilidad: $f'(\pm\sqrt{r}) = r - 3(\pm\sqrt{r})^2 = -2r < 0 \implies$ \textbf{Estables}.
\end{enumerate}

\begin{figure}[htbp]
\centering
\begin{tikzpicture}[>=Stealth, scale=0.82]
    % Panel a: r < 0
    \begin{scope}[shift={(-6.2, 3.2)}]
        \draw[->, thick, black!60] (-2.2,0) -- (2.2,0) node[right, font=\tiny] {$x$};
        \draw[->, thick, black!60] (0,-1.5) -- (0,1.8) node[above, font=\tiny] {$\dot{x}$};
        \draw[thick, black!85, domain=-1.4:1.4, samples=50] plot (\x, {-0.8*\x - \x*\x*\x});
        \filldraw[black] (0,0) circle (2pt) node[above right=1pt, font=\tiny] {$0$ (estable)};
        \draw[->, thick, black!90] (-1.3,0) -- (-0.5,0);
        \draw[->, thick, black!90] (1.3,0) -- (0.5,0);
        \node[font=\footnotesize, align=center] at (0,-1.8) {\textbf{(a)} $r < 0$ (pozo único)};
    \end{scope}

    % Panel b: r = 0
    \begin{scope}[shift={(0, 3.2)}]
        \draw[->, thick, black!60] (-2.2,0) -- (2.2,0) node[right, font=\tiny] {$x$};
        \draw[->, thick, black!60] (0,-1.5) -- (0,1.8) node[above, font=\tiny] {$\dot{x}$};
        \draw[thick, black!85, domain=-1.3:1.3, samples=50] plot (\x, {-\x*\x*\x});
        \filldraw[black] (0,0) circle (2pt) node[above right=1pt, font=\tiny] {$0$ (\emph{slowing down})};
        \draw[->, thick, black!90] (-1.2,0) -- (-0.4,0);
        \draw[->, thick, black!90] (1.2,0) -- (0.4,0);
        \node[font=\footnotesize, align=center] at (0,-1.8) {\textbf{(b)} $r = 0$ (decae $\sim t^{-1/2}$)};
    \end{scope}

    % Panel c: r > 0
    \begin{scope}[shift={(6.2, 3.2)}]
        \draw[->, thick, black!60] (-2.2,0) -- (2.2,0) node[right, font=\tiny] {$x$};
        \draw[->, thick, black!60] (0,-1.5) -- (0,1.8) node[above, font=\tiny] {$\dot{x}$};
        \draw[thick, black!85, domain=-1.4:1.4, samples=50] plot (\x, {1.2*\x - \x*\x*\x});
        \filldraw[black] (-1.095, 0) circle (2pt) node[below=2pt, font=\tiny] {$-\sqrt{r}$};
        \draw[black, fill=white, thick] (0, 0) circle (2pt) node[above right=1pt, font=\tiny] {$0$};
        \filldraw[black] (1.095, 0) circle (2pt) node[above=2pt, font=\tiny] {$+\sqrt{r}$};
        \draw[->, thick, black!90] (-1.8,0) -- (-1.3,0);
        \draw[->, thick, black!90] (-0.3,0) -- (-0.8,0);
        \draw[->, thick, black!90] (0.3,0) -- (0.8,0);
        \draw[->, thick, black!90] (1.8,0) -- (1.3,0);
        \node[font=\footnotesize, align=center] at (0,-1.8) {\textbf{(c)} $r > 0$ (biestabilidad)};
    \end{scope}

    % Panel d: Diagrama de bifurcación
    \begin{scope}[shift={(0, -2.8)}]
        \draw[->, thick, black!60] (-3.0,0) -- (3.0,0) node[right, font=\small] {$r$};
        \draw[->, thick, black!60] (0,-2.2) -- (0,2.2) node[above, font=\small] {$x^*$};
        
        % Rama trivial
        \draw[very thick, black!85] (-2.5,0) -- (0,0);
        \draw[very thick, dashed, black!75] (0,0) -- (2.5,0);
        
        % Parábolas simétricas x* = \pm \sqrt{r} => r = (x*)^2
        \draw[very thick, black!85, domain=0:1.8, samples=50] plot ({\x*\x*0.7}, \x);
        \draw[very thick, black!85, domain=0:1.8, samples=50] plot ({\x*\x*0.7}, {-\x});
        
        \filldraw[black] (0,0) circle (2.5pt) node[below right=2pt, font=\footnotesize] {Pitchfork $(0,0)$};
        \node[above, font=\footnotesize] at (1.8, 1.4) {$x^* = +\sqrt{r}$ (estable)};
        \node[below, font=\footnotesize] at (1.8, -1.4) {$x^* = -\sqrt{r}$ (estable)};
        \node[font=\bfseries\small] at (0, -2.6) {\textbf{(d)} Diagrama de bifurcación pitchfork supercrítica};
    \end{scope}
\end{tikzpicture}
\caption{Bifurcación pitchfork supercrítica en $\dot{x} = rx - x^3$. Para $r > 0$, el origen pierde estabilidad y nacen dos ramas simétricas estables $x^* = \pm\sqrt{r}$ que se abren como un diapasón.}
\label{fig:pitchfork_supercritica_fases}
\end{figure}

\subsection{Potencial de la pitchfork}

Integrando la ecuación $\dot{x} = -V'(x) = rx - x^3$, obtenemos la función potencial:
\begin{equation}
    V(x) = -\frac{r}{2}x^2 + \frac{1}{4}x^4
    \label{eq:potencial_pitchfork_paneles}
\end{equation}
En la \cref{fig:potenciales_pitchfork_paneles} observamos la deformación geométrica del potencial al incrementar $r$:
\begin{itemize}
    \item Para $r = -1 < 0$: $V(x) = \frac{1}{2}x^2 + \frac{1}{4}x^4$. El potencial tiene un único pozo cuadrático profundo en $x^* = 0$.
    \item Para $r = 0$: $V(x) = \frac{1}{4}x^4$. El fondo del pozo se aplana por completo ($V''(0) = 0$), explicando la extrema lentitud de relajación del sistema (\emph{critical slowing down}).
    \item Para $r = 2 > 0$: $V(x) = -x^2 + \frac{1}{4}x^4$. El fondo central se levanta transformándose en una barrera inestable (máximo local), mientras se crean dos pozos simétricos de mínima energía en $x^* = \pm\sqrt{2}$. La partícula rueda hacia uno de los dos pozos, rompiendo espontáneamente la simetría (\emph{spontaneous symmetry breaking}).
\end{itemize}

\begin{figure}[htbp]
\centering
\begin{tikzpicture}[>=Stealth, scale=0.82]
    % Panel 1: r = -1
    \begin{scope}[shift={(-6.2, 0)}]
        \draw[->, thick, black!60] (-2.2,0) -- (2.2,0) node[right, font=\tiny] {$x$};
        \draw[->, thick, black!60] (0,-0.8) -- (0,2.5) node[above, font=\tiny] {$V(x)$};
        \draw[very thick, black!85, domain=-1.6:1.6, samples=50] plot (\x, {0.6*\x*\x + 0.3*\x*\x*\x*\x});
        \filldraw[black] (0,0) circle (2pt);
        \node[font=\footnotesize, align=center] at (0,-1.2) {\textbf{(a)} $r = -1$\\Pozo simple};
    \end{scope}

    % Panel 2: r = 0
    \begin{scope}[shift={(0, 0)}]
        \draw[->, thick, black!60] (-2.2,0) -- (2.2,0) node[right, font=\tiny] {$x$};
        \draw[->, thick, black!60] (0,-0.8) -- (0,2.5) node[above, font=\tiny] {$V(x)$};
        \draw[very thick, black!85, domain=-1.6:1.6, samples=50] plot (\x, {0.4*\x*\x*\x*\x});
        \filldraw[black] (0,0) circle (2pt);
        \node[font=\footnotesize, align=center] at (0,-1.2) {\textbf{(b)} $r = 0$\\Fondo plano};
    \end{scope}

    % Panel 3: r = 2
    \begin{scope}[shift={(6.2, 0)}]
        \draw[->, thick, black!60] (-2.2,0) -- (2.2,0) node[right, font=\tiny] {$x$};
        \draw[->, thick, black!60] (0,-1.5) -- (0,2.0) node[above, font=\tiny] {$V(x)$};
        \draw[very thick, black!85, domain=-1.9:1.9, samples=60] plot (\x, {-1.0*\x*\x + 0.4*\x*\x*\x*\x});
        \draw[black, fill=white, thick] (0,0) circle (2pt);
        \filldraw[black] (-1.118, -0.625) circle (2pt);
        \filldraw[black] (1.118, -0.625) circle (2pt);
        \node[font=\footnotesize, align=center] at (0,-1.8) {\textbf{(c)} $r = 2$\\Doble pozo simétrico};
    \end{scope}
\end{tikzpicture}
\caption{Evolución del paisaje de potencial $V(x) = -\frac{r}{2}x^2 + \frac{1}{4}x^4$ en la bifurcación pitchfork supercrítica para $r=-1$, $r=0$ y $r=2$.}
\label{fig:potenciales_pitchfork_paneles}
\end{figure}

\subsection{Pitchfork subcrítica $\dot{x} = rx + x^{3}$ e histéresis}

En la \textbf{bifurcación pitchfork subcrítica}, el término cúbico no lineal tiene signo positivo ($+x^3$), actuando como una fuerza desestabilizadora que expulsa a las trayectorias hacia el infinito:
\begin{equation}
    \dot{x} = rx + x^3 = x(r + x^2)
    \label{eq:pitchfork_subcritica}
\end{equation}

\paragraph{Análisis dinámico y divergencia en tiempo finito:}
\begin{itemize}
    \item \textbf{Para $r < 0$:} Existen tres equilibrios. El origen $x^* = 0$ es localmente estable ($f'(0) = r < 0$), pero está flanqueado por dos puntos de equilibrio simétricos \textbf{inestables} en:
    \begin{equation}
        x_\pm^* = \pm\sqrt{-r}
    \end{equation}
    Estos dos puntos inestables fijan la frontera de la cuenca de atracción del origen. Si una perturbación saca al sistema fuera del intervalo $(-\sqrt{-r}, +\sqrt{-r})$, el término $+x^3$ domina y provoca un crecimiento explosivo que alcanza infinito en tiempo finito (\emph{finite-time blow-up}):
    \begin{equation}
        \dot{x} \approx x^3 \implies \int_{x_0}^\infty \frac{dx}{x^3} = \int_0^{t_\infty} dt \implies t_\infty = \frac{1}{2x_0^2} < \infty
    \end{equation}
    \item \textbf{Para $r \ge 0$:} El origen se vuelve inestable ($f'(0) \ge 0$) y no existe ningún equilibrio cercano: cualquier condición inicial $x_0 \neq 0$ diverge catastróficamente.
\end{itemize}

\paragraph{Estabilización física de orden superior e histéresis.}
En cualquier sistema físico o biológico real, las trayectorias no escapan a infinito sino que quedan confinadas por términos no lineales de orden superior estabilizadores, tales como un término quíntico $-x^5$:
\begin{equation}
    \dot{x} = rx + x^3 - x^5
    \label{eq:pitchfork_subcritica_estabilizada}
\end{equation}
La \cref{fig:pitchfork_subcritica_histeresis} ilustra el diagrama de bifurcación resultante: al aumentar $r$ lentamente desde valores negativos, el sistema permanece en el estado $x^* = 0$ hasta cruzar $r = 0$, momento en el cual salta bruscamente (\emph{hard jump}) a la rama de gran amplitud $x^* \approx +1$. Al reducir $r$ nuevamente, el sistema no regresa a cero en $r=0$, sino que permanece en la rama excitada hasta alcanzar un punto de pliegue saddle-node $r_s < 0$, donde colapsa de vuelta al reposo. Este fenómeno de dependencia de la historia previa se denomina \textbf{histéresis}.

\begin{figure}[htbp]
\centering
\begin{tikzpicture}[>=Stealth, scale=0.85]
    \draw[->, thick, black!60] (-3.2,0) -- (3.2,0) node[right, font=\small] {$r$};
    \draw[->, thick, black!60] (0,-2.2) -- (0,2.2) node[above, font=\small] {$x^*$};
    
    % Rama trivial
    \draw[very thick, black!85] (-2.8,0) -- (0,0);
    \draw[very thick, dashed, black!75] (0,0) -- (2.8,0);
    
    % Ramas subcríticas inestables que retroceden (dashed)
    \draw[thick, dashed, black!75, domain=0:1.0, samples=40] plot ({-\x*\x*0.9}, \x);
    \draw[thick, dashed, black!75, domain=0:1.0, samples=40] plot ({-\x*\x*0.9}, {-\x});
    
    % Pliegue y ramas exteriores estables (saturadas por -x^5)
    \draw[very thick, black!85, domain=1.0:1.6, samples=40] plot ({-0.9 + (\x-1.0)*(\x-1.0)*3.5}, \x);
    \draw[very thick, black!85, domain=1.0:1.6, samples=40] plot ({-0.9 + (\x-1.0)*(\x-1.0)*3.5}, {-\x});
    
    % Flechas de salto catastrófico e histéresis
    \draw[->, double, thick, black!90] (0.05, 0.1) -- (0.05, 1.15) node[right=2pt, font=\tiny] {Salto en $r=0$};
    \draw[->, double, thick, black!90] (-0.85, 0.95) -- (-0.85, 0.1) node[left=2pt, font=\tiny] {Colapso en $r_s$};
    
    \node[font=\footnotesize] at (0, -2.5) {Ciclo de histéresis y biestabilidad};
\end{tikzpicture}
\caption{Bifurcación pitchfork subcrítica estabilizada por términos de orden superior $\dot{x} = rx + x^3 - x^5$. Muestra biestabilidad entre el estado de reposo y los estados de gran amplitud, junto con saltos catastróficos y un ciclo de histéresis.}
\label{fig:pitchfork_subcritica_histeresis}
\end{figure}


% ============================================================
% FUENTES: "Notas Biomate 1.pdf" (español) — NO está en el draft.
% ============================================================
\section{Aplicación: redes neuronales y función de activación sigmoidal}

Como aplicación biológica y computacional directa de la bifurcación pitchfork, consideremos el modelo dinámico de una neurona artificial o ensamble neuronal auto-recurrente con memoria.

\subsection{Motivación: redes recurrentes (RNN) y gradiente desvaneciente}

En el aprendizaje profundo y en los modelos continuos de neurociencia computacional, una de las funciones de activación más emblemáticas para mitigar el desvanecimiento del gradiente en redes recurrentes (\emph{Recurrent Neural Networks}, RNN) es la \textbf{tangente hiperbólica}:
\begin{equation}
    \tanh x = \frac{e^x - e^{-x}}{e^x + e^{-x}}
\end{equation}
Consideremos la dinámica de activación de una neurona recurrente con decaimiento pasivo lineal (tipo \emph{leaky integrator}):
\begin{equation}
    \dot{x} = -x + \beta \tanh x
    \label{eq:neurona_ode}
\end{equation}
donde:
\begin{itemize}
    \item $x(t) \in \mathbb{R}$ representa el potencial o nivel de activación de la neurona.
    \item $-x$ modela la fuga pasiva hacia el potencial de reposo en ausencia de estímulo.
    \item $\beta > 0$ es el \textbf{peso sináptico recurrente} o ganancia de auto-excitación.
    \item $\tanh x$ es la función de activación no lineal sigmoidal que satura en $\pm 1$ para valores extremos de excitación.
\end{itemize}

Notemos que el sistema posee simetría impar estricta ante reflexiones: si $f(x) = -x + \beta\tanh x$, entonces $f(-x) = -(-x) + \beta\tanh(-x) = x - \beta\tanh x = -f(x)$. Por tanto, el campo es equivariante.

\subsection{Análisis geométrico: intersección de $x$ y $\beta\tanh x$}

Los puntos fijos del sistema satisfacen la condición de equilibrio:
\begin{equation}
    -x + \beta \tanh x = 0 \iff x^* = \beta \tanh(x^*)
    \label{eq:interseccion_neurona}
\end{equation}
Geométricamente, los equilibrios corresponden a los puntos de intersección entre la recta identidad $y = x$ y la función sigmoidal escalada $y = \beta \tanh x$.

Dado que $\tanh(0) = 0$, el origen $x^* = 0$ es siempre una solución para cualquier valor de $\beta$. Para saber si existen soluciones adicionales no nulas, comparamos las pendientes de ambas curvas en el origen:
\begin{equation}
    \left.\frac{d}{dx}(\beta \tanh x)\right|_{x=0} = \beta \operatorname{sech}^2(0) = \beta (1) = \beta
\end{equation}

\begin{figure}[htbp]
\centering
\begin{tikzpicture}[>=Stealth, scale=0.82]
    % Panel a: beta < 1 (beta = 0.6)
    \begin{scope}[shift={(-6.2, 0)}]
        \draw[->, thick, black!60] (-2.2,0) -- (2.2,0) node[right, font=\tiny] {$x$};
        \draw[->, thick, black!60] (0,-1.8) -- (0,1.8) node[above, font=\tiny] {$y$};
        \draw[dashed, black!50, thick] (-1.5,-1.5) -- (1.5,1.5) node[right, font=\tiny] {$y=x$};
        \draw[very thick, black!85, domain=-1.8:1.8, samples=50] plot (\x, {0.6*tanh(\x)});
        \filldraw[black] (0,0) circle (2pt);
        \node[font=\footnotesize, align=center] at (0,-2.1) {\textbf{(a)} $\beta < 1$ ($\beta=0.6$)\\1 equilibrio estable};
    \end{scope}

    % Panel b: beta = 1
    \begin{scope}[shift={(0, 0)}]
        \draw[->, thick, black!60] (-2.2,0) -- (2.2,0) node[right, font=\tiny] {$x$};
        \draw[->, thick, black!60] (0,-1.8) -- (0,1.8) node[above, font=\tiny] {$y$};
        \draw[dashed, black!50, thick] (-1.5,-1.5) -- (1.5,1.5) node[right, font=\tiny] {$y=x$};
        \draw[very thick, black!85, domain=-1.8:1.8, samples=50] plot (\x, {1.0*tanh(\x)});
        \filldraw[black] (0,0) circle (2pt);
        \node[font=\footnotesize, align=center] at (0,-2.1) {\textbf{(b)} $\beta = 1$ (tangencia)\\Bifurcación crítica};
    \end{scope}

    % Panel c: beta > 1 (beta = 1.8)
    \begin{scope}[shift={(6.2, 0)}]
        \draw[->, thick, black!60] (-2.2,0) -- (2.2,0) node[right, font=\tiny] {$x$};
        \draw[->, thick, black!60] (0,-1.8) -- (0,1.8) node[above, font=\tiny] {$y$};
        \draw[dashed, black!50, thick] (-1.5,-1.5) -- (1.5,1.5) node[right, font=\tiny] {$y=x$};
        \draw[very thick, black!85, domain=-1.8:1.8, samples=50] plot (\x, {1.8*tanh(\x)});
        \filldraw[black] (-1.4, -1.4) circle (2pt) node[below left, font=\tiny] {$-x^*$};
        \draw[black, fill=white, thick] (0,0) circle (2pt);
        \filldraw[black] (1.4, 1.4) circle (2pt) node[above right, font=\tiny] {$+x^*$};
        \node[font=\footnotesize, align=center] at (0,-2.1) {\textbf{(c)} $\beta > 1$ ($\beta=1.8$)\\3 equilibrios (biestabilidad)};
    \end{scope}
\end{tikzpicture}
\caption{Determinación geométrica de los puntos de equilibrio de la neurona recurrente $\dot{x} = -x + \beta\tanh x$ mediante la intersección de $y=x$ con $y=\beta\tanh x$.}
\label{fig:intersecciones_red_neuronal}
\end{figure}

Surgen tres comportamientos cualitativos según la ganancia $\beta$:
\begin{enumerate}
    \item \textbf{Caso $\beta < 1$ (subcrítico):} La pendiente en el origen es menor que 1. Como $\tanh x$ es estrictamente cóncava para $x > 0$, la curva $y = \beta\tanh x$ permanece por debajo de la recta $y = x$ en todo el semieje positivo. El único punto de intersección es $x^* = 0$.
    La derivada es $f'(0) = -1 + \beta < 0$, por lo que el origen es un \textbf{atractor asintóticamente estable global}.
    \item \textbf{Caso $\beta = 1$ (umbral crítico):} La curva $y = \tanh x$ es exactamente tangente a la identidad en el origen. El origen es el único punto fijo y es estable con decaimiento subexponencial (\emph{critical slowing down}).
    \item \textbf{Caso $\beta > 1$ (supercrítico):} La pendiente en el origen es mayor que 1. La curva $\beta\tanh x$ comienza por encima de $y = x$ y, debido a que satura asintóticamente hacia el valor finito $\beta$, forzosamente debe cruzar a la identidad en un valor $x_+^* > 0$. Por simetría impar, existe una intersección idéntica en $-x_+^* < 0$.
    El origen $x^* = 0$ pierde estabilidad ($f'(0) = \beta - 1 > 0$, inestable), mientras que los dos nuevos estados activos $\pm x_+^*$ son \textbf{atractores estables}.
\end{enumerate}

\subsection{Verificación analítica con serie de Taylor de $\tanh x$}

Para demostrar analíticamente que el sistema experimenta una bifurcación pitchfork supercrítica canónica en $\beta_c = 1$, expandimos $\tanh x$ en serie de Maclaurin cerca de $x = 0$:
\begin{equation}
    \tanh x = x - \frac{x^3}{3} + \frac{2x^5}{15} - \mathcal{O}(x^7)
\end{equation}
Sustituyendo en la ecuación diferencial \eqref{eq:neurona_ode}:
\begin{align}
    \dot{x} &= -x + \beta \left(x - \frac{x^3}{3} + \mathcal{O}(x^5)\right) \nonumber \\
    &= (\beta - 1)x - \frac{\beta}{3}x^3 + \mathcal{O}(x^5)
    \label{eq:taylor_neurona}
\end{align}
Cerca del punto de bifurcación ($\beta \approx 1$), podemos aproximar $\beta/3 \approx 1/3$:
\begin{equation}
    \dot{x} \approx (\beta - 1)x - \frac{1}{3}x^3
\end{equation}
Haciendo el cambio de parámetro $r = \beta - 1$, recuperamos de forma exacta la forma normal supercrítica $\dot{x} = rx - \frac{1}{3}x^3$. Las ramas de equilibrio no nulas para $\beta > 1$ vienen dadas localmente por:
\begin{equation}
    (\beta - 1)x^* - \frac{1}{3}(x^*)^3 \approx 0 \implies x_\pm^* \approx \pm\sqrt{3(\beta - 1)}
    \label{eq:ramas_analiticas_neurona}
\end{equation}

\begin{figure}[htbp]
\centering
\begin{tikzpicture}[>=Stealth, scale=0.85]
    \draw[->, thick, black!60] (0,0) -- (4.5,0) node[right, font=\small] {$\beta$ (ganancia)};
    \draw[->, thick, black!60] (0,-2.2) -- (0,2.2) node[above, font=\small] {$x^*$};
    
    % Rama trivial
    \draw[ultra thick, black!85] (0,0) -- (1.5,0);
    \draw[very thick, dashed, black!75] (1.5,0) -- (4.0,0);
    
    % Ramas de la pitchfork x* = \pm \sqrt{3(\beta - 1)} => \beta = 1 + (x*)^2 / 3
    \draw[ultra thick, black!85, domain=0:1.6, samples=50] plot ({1.5 + \x*\x*0.9}, \x);
    \draw[ultra thick, black!85, domain=0:1.6, samples=50] plot ({1.5 + \x*\x*0.9}, {-\x});
    
    \filldraw[black] (1.5,0) circle (2.5pt) node[below=4pt, font=\footnotesize] {$\beta_c = 1$};
    \node[above, font=\footnotesize] at (3.2, 1.4) {$x_+^*$ (memoria $+1$)};
    \node[below, font=\footnotesize] at (3.2, -1.4) {$x_-^*$ (memoria $-1$)};
    \node[above, font=\footnotesize] at (0.8, 0.1) {Reposo inactivo};
\end{tikzpicture}
\caption{Diagrama de bifurcación pitchfork de la neurona recurrente en función del parámetro de ganancia sináptica $\beta$. Para $\beta > 1$, la red almacena un bit de memoria binaria biestable.}
\label{fig:bifurcacion_red_neuronal}
\end{figure}

Esta bifurcación es el mecanismo fundamental mediante el cual las redes recurrentes implementan \textbf{memoria asociativa y toma de decisiones biestable}: para $\beta > 1$, el sistema puede conmutar entre los atractores $+x^*$ (estado binario $1$) y $-x^*$ (estado binario $0$).

% ============================================================
% FUENTES: "Notas Biomate 1.pdf" (español) — NO está en el draft.
% ============================================================
\section{Paréntesis numérico: el método de Newton--Raphson}

En el capítulo anterior utilizamos series de Taylor para aproximar las ramas de equilibrio de la neurona recurrente en las inmediaciones del punto crítico ($\beta \approx 1$). Sin embargo, cuando el parámetro de ganancia se aleja del umbral (por ejemplo, $\beta = 2$ o $\beta = 5$), las aproximaciones polinomiales locales pierden precisión cuantitativa.

\subsection{La ecuación implícita $x^{*} = \beta\tanh x^{*}$}

Para determinar con exactitud arbitraria la ubicación de los atractores biestables para cualquier valor de $\beta > 1$, debemos resolver la ecuación de punto fijo:
\begin{equation}
    x^* = \beta \tanh(x^*) \iff g(x^*) \equiv \beta \tanh(x^*) - x^* = 0
    \label{eq:raiz_implicita_nr}
\end{equation}
La ecuación \eqref{eq:raiz_implicita_nr} es una \textbf{ecuación trascendente e implícita}: no existe ningún procedimiento algebraico finito que permita despejar $x^*(\beta)$ en términos de funciones elementales cerradas.

Ante la imposibilidad de un despeje analítico exacto, recurrimos al cálculo numérico. El método iterativo más eficiente y extendido para hallar raíces reales de ecuaciones no lineales es el \textbf{método de Newton--Raphson}.

\subsection{Algoritmo de Newton--Raphson y criterio de parada}

Sea $g(x)$ una función continuamente diferenciable en un intervalo que contiene una raíz $x^*$ tal que $g(x^*) = 0$. Supongamos que disponemos de una estimación inicial (o ``adivinanza'') $x_0 \approx x^*$.

\paragraph{Derivación geométrica por rectas tangentes.}
Aproximamos la curva $y = g(x)$ por su recta tangente en el punto $(x_n, g(x_n))$:
\begin{equation}
    y - g(x_n) = g'(x_n)(x - x_n)
\end{equation}
Definimos la siguiente aproximación $x_{n+1}$ como la intersección de dicha recta tangente con el eje horizontal ($y = 0$):
\begin{equation}
    0 - g(x_n) = g'(x_n)(x_{n+1} - x_n) \implies \boxed{x_{n+1} = x_n - \frac{g(x_n)}{g'(x_n)}}
    \label{eq:formula_newton_raphson}
\end{equation}

\begin{figure}[htbp]
\centering
\begin{tikzpicture}[>=Stealth, scale=0.9]
    \draw[->, thick, black!60] (-0.5,0) -- (5.0,0) node[right, font=\small] {$x$};
    \draw[->, thick, black!60] (0,-1.5) -- (0,3.0) node[above, font=\small] {$g(x)$};
    
    % Curva g(x)
    \draw[very thick, black!85, domain=0.5:4.5, samples=50] plot (\x, {1.8*tanh(0.9*\x) - 0.5*\x - 0.2});
    
    % Raíz verdadera x*
    \filldraw[black] (2.78, 0) circle (2.5pt) node[below left=2pt, font=\footnotesize] {Raíz $x^*$};
    
    % Estimación inicial x_0
    \draw[dashed, black!50] (4.2, 0) node[below, font=\footnotesize] {$x_0$} -- (4.2, 1.25);
    \filldraw[black] (4.2, 1.25) circle (2pt);
    
    % Tangente en x_0
    \draw[thick, black!70] (4.6, 1.55) -- (2.3, -0.18);
    \filldraw[black] (2.55, 0) circle (2pt) node[below=3pt, font=\footnotesize] {$x_1$};
    
    % Tangente en x_1
    \draw[dashed, black!50] (2.55, 0) -- (2.55, -0.15);
    \filldraw[black] (2.55, -0.15) circle (1.5pt);
    
    \node[font=\footnotesize, align=center] at (1.2, 2.2) {Convergencia cuadrática\\$x_{n+1} = x_n - \frac{g(x_n)}{g'(x_n)}$};
\end{tikzpicture}
\caption{Interpretación geométrica del método de Newton--Raphson. Cada iteración proyecta la recta tangente local hacia el eje de abscisas, acelerando cuadráticamente la convergencia hacia la raíz verdadera $x^*$.}
\label{fig:newton_raphson_geometrico}
\end{figure}

\paragraph{Aplicación al sistema neuronal:}
Para nuestra función $g(x) = \beta \tanh x - x$, calculamos su primera derivada analítica:
\begin{equation}
    g'(x) = \beta \operatorname{sech}^2 x - 1
\end{equation}
Sustituyendo en la regla de actualización \eqref{eq:formula_newton_raphson}:
\begin{equation}
    x_{n+1} = x_n - \frac{\beta \tanh x_n - x_n}{\beta \operatorname{sech}^2 x_n - 1}
    \label{eq:iteracion_neurona_nr}
\end{equation}

\paragraph{Criterio de convergencia y tolerancia.}
Dado que una computadora opera con aritmética de precisión finita, el proceso iterativo debe detenerse cuando la diferencia entre dos estimaciones sucesivas cae por debajo de una tolerancia predefinida $\varepsilon$ (típicamente $\varepsilon = 10^{-8}$):
\begin{equation}
    |x_{n+1} - x_n| < \varepsilon \quad \text{o bien} \quad |g(x_{n+1})| < \varepsilon
\end{equation}

\paragraph{Sensibilidad a la semilla inicial $x_0$:}
\begin{itemize}
    \item Si elegimos $x_0 > 0$, el método converge rápidamente al atractor activo positivo $+x^*$.
    \item Si elegimos $x_0 < 0$, el algoritmo converge al atractor activo negativo simétrico $-x^*$.
    \item Si elegimos $x_0 = 0$, dado que $g(0) = 0$, el método se detiene en una sola iteración en el punto fijo inestable del reposo.
    \item Si la semilla inicial satisface $g'(x_0) = 0$ (es decir, $\operatorname{sech}^2 x_0 = 1/\beta$), la pendiente se anula y el método falla al producirse una división entre cero.
\end{itemize}

Con este paréntesis numérico concluimos el estudio sistemático de los flujos y bifurcaciones en espacios de fases unidimensionales.


% ---- Parte IV — Sistemas en dos dimensiones ----
\part{Sistemas en dos dimensiones}
% ============================================================
% FUENTES: Notas Biomate 1; raw_plain_nonlinear_dynamics.txt.
% Strogatz (2015), caps. 5 y 6; Edelstein-Keshet (2007).
% ============================================================
\section{Sistemas lineales en dos dimensiones}
\label{sec:lineales2d}

En los capítulos precedentes exploramos la dinámica de flujos sobre la recta real $\mathbb{R}$. Una de las conclusiones fundamentales fue la imposibilidad topológica de oscilaciones periódicas en una dimensión: sobre una línea, el flujo sólo puede avanzar o retroceder monótonamente hacia un punto fijo o hacia el infinito.

Al dar el salto a dos dimensiones ($n = 2$), el espacio de estados adquiere suficiente libertad geométrica para que las trayectorias giren, roten y se entrelacen. En este capítulo desarrollaremos la teoría completa de los \textbf{sistemas lineales en el plano}, que constituye la piedra angular para clasificar y comprender el comportamiento local de cualquier sistema dinámico bidimensional mediante linealización~\cite{strogatz2015,edelstein2007}.

\subsection{Formulación matricial y ecuación característica}

Un sistema lineal bidimensional autónomo homogéneo se escribe en forma vectorial como:
\begin{equation}
    \dot{\vec{x}} = A \vec{x}, \quad \vec{x} = \begin{pmatrix} x \\ y \end{pmatrix} \in \mathbb{R}^2, \quad A = \begin{pmatrix} a & b \\ c & d \end{pmatrix} \in \mathbb{R}^{2\times 2}
    \label{eq:sistema_lineal_2d}
\end{equation}
o en componentes escalares:
\begin{equation}
\begin{aligned}
    \dot{x} &= ax + by, \\
    \dot{y} &= cx + dy.
\end{aligned}
\end{equation}
El origen $\vec{x}^* = (0,0)^\top$ es siempre un punto de equilibrio ($A\vec{0} = \vec{0}$). Si $\det A \neq 0$, el origen es el único punto fijo del sistema.

Buscamos soluciones exponenciales de la forma $\vec{x}(t) = \vec{v} e^{\lambda t}$, donde $\vec{v} \neq \vec{0}$ es un vector propio constante y $\lambda \in \mathbb{C}$ es el valor propio asociado:
\begin{equation}
    \dot{\vec{x}} = \lambda \vec{v} e^{\lambda t} = A \vec{v} e^{\lambda t} \implies (A - \lambda I)\vec{v} = \vec{0}
\end{equation}
Para que existan soluciones no triviales ($\vec{v} \neq \vec{0}$), se requiere que la matriz $(A - \lambda I)$ sea singular:
\begin{equation}
    \det(A - \lambda I) = 0 \iff \det \begin{pmatrix} a - \lambda & b \\ c & d - \lambda \end{pmatrix} = (a - \lambda)(d - \lambda) - bc = 0
\end{equation}
Desarrollando este determinante, obtenemos la \textbf{ecuación característica}:
\begin{equation}
    \lambda^2 - \tau \lambda + \Delta = 0
    \label{eq:ec_caracteristica_2d}
\end{equation}
donde hemos definido dos invariantes algebraicos fundamentales de la matriz $A$:
\begin{align}
    \tau &\equiv \operatorname{tr}(A) = a + d \quad \text{(\textbf{traza})}, \label{eq:def_traza} \\
    \Delta &\equiv \det(A) = ad - bc \quad \text{(\textbf{determinante})}. \label{eq:def_det}
\end{align}

Las raíces de la \cref{eq:ec_caracteristica_2d} son los \textbf{autovalores} de la matriz $A$:
\begin{equation}
    \lambda_{1,2} = \frac{\tau \pm \sqrt{\tau^2 - 4\Delta}}{2}
    \label{eq:autovalores_formula}
\end{equation}
Asimismo, por las relaciones de Cardano--Vieta:
\begin{equation}
    \lambda_1 + \lambda_2 = \tau, \qquad \lambda_1 \lambda_2 = \Delta.
    \label{eq:vieta_autovalores}
\end{equation}
El comportamiento cualitativo del sistema lineal queda enteramente determinado por el signo y la naturaleza de $\lambda_1$ y $\lambda_2$, o de forma equivalente, por la posición del punto $(\tau, \Delta)$ en el plano traza--determinante.

\subsection{Clasificación completa de los equilibrios}

El discriminante de la ecuación cuadrática es:
\begin{equation}
    \mathcal{D} = \tau^2 - 4\Delta
\end{equation}
La parábola $\Delta = \frac{1}{4}\tau^2$ (o $\tau^2 - 4\Delta = 0$) divide el semiplano superior $\Delta > 0$ en dos regiones fundamentales: autovalores reales $(\mathcal{D} > 0)$ y autovalores complejos conjugados $(\mathcal{D} < 0)$.

\begin{definition}[Tipos de puntos fijos en 2D]
Clasificamos el punto fijo en el origen según los valores de $\tau$ y $\Delta$:
\begin{enumerate}
    \item \textbf{Punto de ensilladura o silla (\emph{Saddle point}) ($\Delta < 0$):}
    Como $\Delta = \lambda_1 \lambda_2 < 0$, los autovalores son reales y de signos opuestos: $\lambda_1 < 0 < \lambda_2$. Las trayectorias a lo largo de la dirección del autovector $\vec{v}_1$ convergen exponencialmente hacia el origen (\emph{variedad estable}), mientras que a lo largo de $\vec{v}_2$ divergen exponencialmente (\emph{variedad inestable}). Una silla es siempre \textbf{inestable}.
    
    \item \textbf{Nodos (\emph{Nodes}) ($\Delta > 0$ y $\tau^2 - 4\Delta > 0$):}
    Los autovalores son reales, distintos y del mismo signo:
    \begin{itemize}
        \item \textbf{Nodo estable ($\tau < 0$):} $\lambda_1 < \lambda_2 < 0$. Todas las trayectorias convergen asintóticamente al origen cuando $t \to +\infty$.
        \item \textbf{Nodo inestable ($\tau > 0$):} $0 < \lambda_1 < \lambda_2$. Todas las trayectorias divergen alejándose del origen.
    \end{itemize}

    \item \textbf{Focos o espirales (\emph{Spirals / Foci}) ($\Delta > 0$ y $\tau^2 - 4\Delta < 0$):}
    Los autovalores son complejos conjugados: $\lambda_{1,2} = \alpha \pm i\omega$, donde la parte real es $\alpha = \frac{\tau}{2}$ y la frecuencia angular es $\omega = \frac{\sqrt{4\Delta - \tau^2}}{2} > 0$. La solución tiene la forma $\vec{x}(t) \sim e^{\alpha t}(\vec{c}_1 \cos\omega t + \vec{c}_2 \sin\omega t)$:
    \begin{itemize}
        \item \textbf{Espiral estable ($\tau < 0 \implies \alpha < 0$):} Las trayectorias oscilan y giran en espiral convergiendo hacia el origen (atractor).
        \item \textbf{Espiral inestable ($\tau > 0 \implies \alpha > 0$):} Las trayectorias oscilan en espiral con amplitud creciente alejándose del origen (repulsor).
    \end{itemize}

    \item \textbf{Centro (\emph{Center}) ($\tau = 0$ y $\Delta > 0$):}
    Los autovalores son puramente imaginarios: $\lambda_{1,2} = \pm i\sqrt{\Delta}$ ($\alpha = 0$). El origen está rodeado por una familia continua de órbitas cerradas elípticas concéntricas de periodo $T = \frac{2\pi}{\sqrt{\Delta}}$. El centro es \textbf{neutramente estable} en el sentido de Lyapunov, pero no es asintóticamente estable.

    \item \textbf{Casos degenerados y de frontera:}
    \begin{itemize}
        \item \textbf{Nodos estrella y nodos impropios ($\Delta = \frac{1}{4}\tau^2$):} Autovalores reales repetidos $\lambda_1 = \lambda_2 = \tau/2$. Si $A = \lambda I$, todas las direcciones son autovectores (\emph{star node}); si $A$ no es diagonalizable, existe un único autovector independiente y las trayectorias son tangentes a él (\emph{degenerate node}).
        \item \textbf{Línea de equilibrios no aislados ($\Delta = 0$):} Al menos un autovalor es cero ($\lambda_1 = 0$), por lo que existe una recta continua de puntos fijos correspondiente al núcleo de $A$.
    \end{itemize}
\end{enumerate}
\end{definition}

\begin{figure}[htbp]
\centering
\begin{tikzpicture}[>=Stealth, scale=1.05]
    % Regiones sombreadas
    % Semiplano inferior: Sillas (Delta < 0)
    \fill[black!6] (-4.5,-2.2) rectangle (4.5,0);
    
    % Nodos estables: Delta > 0, tau < 0, Delta < tau^2/4
    \fill[black!12, domain=-4.2:0, variable=\t]
        (-4.2,0) -- plot (\t, {0.25*\t*\t}) -- (0,0) -- cycle;
        
    % Nodos inestables: Delta > 0, tau > 0, Delta < tau^2/4
    \fill[black!12, domain=0:4.2, variable=\t]
        (0,0) -- plot (\t, {0.25*\t*\t}) -- (4.2,0) -- cycle;

    % Espirales estables: tau < 0, Delta > tau^2/4
    \fill[black!20, domain=-3.5:0, variable=\t]
        (-3.5, 3.2) -- plot (\t, {max(0.25*\t*\t, 0)}) -- (0, 3.2) -- cycle;

    % Espirales inestables: tau > 0, Delta > tau^2/4
    \fill[black!20, domain=0:3.5, variable=\t]
        (0, 3.2) -- plot (\t, {max(0.25*\t*\t, 0)}) -- (3.5, 3.2) -- cycle;

    % Ejes coordenados
    \draw[->, thick, black!70] (-4.6,0) -- (4.6,0) node[right, font=\small] {$\tau = \operatorname{tr}(A)$};
    \draw[->, thick, black!70] (0,-2.3) -- (0,3.5) node[above, font=\small] {$\Delta = \det(A)$};

    % Parábola Delta = tau^2 / 4
    \draw[very thick, black!90, domain=-3.55:3.55, samples=100] plot (\x, {0.25*\x*\x});
    \node[above, font=\footnotesize] at (-2.8, 2.6) {$\Delta = \frac{1}{4}\tau^2$};

    % Etiquetas de regiones
    \node[font=\bfseries\small, align=center] at (0, -1.1) {SILLAS\\(\emph{Saddles})\\$\Delta < 0$};
    
    \node[font=\footnotesize, align=center] at (-2.8, 0.6) {\textbf{Nodos}\\estables\\$\tau < 0$};
    \node[font=\footnotesize, align=center] at (2.8, 0.6) {\textbf{Nodos}\\inestables\\$\tau > 0$};
    
    \node[font=\footnotesize, align=center] at (-1.4, 2.2) {\textbf{Espirales}\\estables\\$\tau < 0$};
    \node[font=\footnotesize, align=center] at (1.4, 2.2) {\textbf{Espirales}\\inestables\\$\tau > 0$};
    
    % Eje vertical positivo: Centros
    \draw[very thick, black!90] (0,0) -- (0,3.3);
    \filldraw[black] (0, 1.6) circle (2pt);
    \node[left=3pt, font=\bfseries\footnotesize, align=right] at (0, 1.6) {Centros\\$\tau = 0, \Delta > 0$};

    % Eje horizontal: Equilibrios degenerados
    \node[below, font=\tiny, black!80] at (-3.2, -0.05) {$\Delta = 0$ (Línea de equilibrios)};
    \node[below, font=\tiny, black!80] at (3.2, -0.05) {$\Delta = 0$};
    
    % Frontera parabólica: Nodos degenerados / estrella
    \node[font=\tiny, rotate=38, align=center] at (2.1, 1.35) {Nodos estrella / impropios};
    \node[font=\tiny, rotate=-38, align=center] at (-2.1, 1.35) {Nodos estrella / impropios};
\end{tikzpicture}
\caption{Diagrama maestro de clasificación en el plano traza--determinante $(\tau, \Delta)$. La parábola $\Delta = \frac{1}{4}\tau^2$ separa los nodos reales de las espirales oscilatorias; el eje $\tau = 0$ marca la frontera conservativa de centros; y el semiplano inferior $\Delta < 0$ corresponde a sillas inestables.}
\label{fig:mapa_tr_det}
\end{figure}

\subsection{Linealización de sistemas no lineales en 2D: la matriz Jacobiana}

Consideremos ahora un sistema dinámico bidimensional no lineal general:
\begin{equation}
\begin{aligned}
    \dot{x} &= f(x, y), \\
    \dot{y} &= g(x, y),
\end{aligned}
\label{eq:sistema_no_lineal_2d}
\end{equation}
con $f, g \in \mathcal{C}^1$. Sea $(x^*, y^*)$ un punto fijo de equilibrio, de modo que $f(x^*, y^*) = 0$ y $g(x^*, y^*) = 0$.

Para estudiar el comportamiento en las inmediaciones del equilibrio, introducimos pequeñas perturbaciones $u(t) = x(t) - x^*$ y $v(t) = y(t) - y^*$. Desarrollando en serie de Taylor bidimensional en torno a $(x^*, y^*)$:
\begin{equation}
\begin{aligned}
    \dot{u} &= \dot{x} = f(x^*+u, y^*+v) \\
            &= f(x^*,y^*) + \frac{\partial f}{\partial x}(x^*,y^*) u + \frac{\partial f}{\partial y}(x^*,y^*) v + \mathcal{O}(u^2, v^2), \\
    \dot{v} &= \dot{y} = g(x^*+u, y^*+v) \\
            &= g(x^*,y^*) + \frac{\partial g}{\partial x}(x^*,y^*) u + \frac{\partial g}{\partial y}(x^*,y^*) v + \mathcal{O}(u^2, v^2).
\end{aligned}
\end{equation}
Despreciando los términos de segundo orden y superiores, obtenemos el sistema linealizado:
\begin{equation}
    \begin{pmatrix} \dot{u} \\ \dot{v} \end{pmatrix} \approx J(x^*, y^*) \begin{pmatrix} u \\ v \end{pmatrix}
    \label{eq:linealizacion_jacobiano}
\end{equation}
donde $J(x^*, y^*)$ es la \textbf{matriz jacobiana} evaluada en el equilibrio:
\begin{equation}
    J(x^*, y^*) = \begin{pmatrix}
        \dfrac{\partial f}{\partial x} & \dfrac{\partial f}{\partial y} \\[8pt]
        \dfrac{\partial g}{\partial x} & \dfrac{\partial g}{\partial y}
    \end{pmatrix}_{(x^*, y^*)}
    \label{eq:def_jacobiano}
\end{equation}

\begin{theorem}[Teorema de Hartman--Grobman (interpretación intuitiva)]
Si el punto fijo $(x^*, y^*)$ es \textbf{hiperbólico} (es decir, ningún autovalor de la matriz jacobiana $J(x^*, y^*)$ tiene parte real nula: $\operatorname{Re}(\lambda_i) \neq 0$), entonces existe un homeomorfismo local que mapea las trayectorias del sistema no lineal \eqref{eq:sistema_no_lineal_2d} en las trayectorias del sistema linealizado \eqref{eq:linealizacion_jacobiano}.
\end{theorem}

En palabras sencillas: para sillas, nodos y espirales, la linealización predice exactamente la topología local del retrato de fases no lineal. La única advertencia ocurre cuando $\operatorname{Re}(\lambda) = 0$ (centros lineales con $\tau = 0, \Delta > 0$ o puntos degenerados con $\Delta = 0$): en tales casos no hiperbólicos, los términos no lineales de orden superior pueden romper la estructura lineal y convertir un centro aparente en una espiral débil atrayente o repelente\aside{Esta es precisamente la razón por la que un centro lineal requiere un análisis más sutil mediante integrales primeras conservativas o funciones de Lyapunov, como estudiaremos en el modelo de Lotka--Volterra y en los sistemas hamiltonianos.}.

\subsection{Ejemplo: clasificación analítica de un sistema en el plano}

\begin{example}[Clasificación de un equilibrio no hiperbólico vs hiperbólico]
Consideremos el sistema lineal:
\begin{equation}
\begin{aligned}
    \dot{x} &= x - y, \\
    \dot{y} &= x + 3y.
\end{aligned}
\end{equation}
La matriz del sistema es $A = \begin{pmatrix} 1 & -1 \\ 1 & 3 \end{pmatrix}$. Calculamos su traza y su determinante:
\begin{equation}
    \tau = \operatorname{tr}(A) = 1 + 3 = 4, \qquad \Delta = \det(A) = (1)(3) - (-1)(1) = 3 + 1 = 4.
\end{equation}
Evaluamos el discriminante:
\begin{equation}
    \tau^2 - 4\Delta = 4^2 - 4(4) = 16 - 16 = 0.
\end{equation}
Como $\Delta = \frac{1}{4}\tau^2 > 0$ y $\tau = 4 > 0$, los autovalores son reales repetidos: $\lambda_1 = \lambda_2 = \tau/2 = 2 > 0$. Para determinar si es un nodo estrella o un nodo impropio, calculamos los autovectores $(A - 2I)\vec{v} = \vec{0}$:
\begin{equation}
    \begin{pmatrix} -1 & -1 \\ 1 & 1 \end{pmatrix} \begin{pmatrix} v_1 \\ v_2 \end{pmatrix} = \begin{pmatrix} 0 \\ 0 \end{pmatrix} \implies v_1 + v_2 = 0 \implies \vec{v}_1 = \begin{pmatrix} 1 \\ -1 \end{pmatrix}.
\end{equation}
Al existir sólo una dirección propia independiente, el origen es un \textbf{nodo impropio inestable (\emph{degenerate unstable node})}.
\end{example}

\subsection{Ejercicios}

\paragraph{Ejercicio 1 (mecánico; clasificar por $\tau$--$\Delta$).} Clasifica el origen de $\dot{\vec{x}}=A\vec{x}$ para (i) $A=\begin{pmatrix}1&2\\2&1\end{pmatrix}$ y (ii) $A=\begin{pmatrix}-3&1\\1&-3\end{pmatrix}$, calculando $\tau$, $\Delta$, el discriminante y los autovalores.
\aside{Pista: usa \eqref{eq:autovalores_formula} y la \cref{fig:mapa_tr_det}.}
% SOL: (i) tau=2, Delta=1-4=-3<0 -> silla (lambda=1±2 => 3,-1). (ii) tau=-6, Delta=8, D=36-32=4>0 -> nodo estable (lambda=-2,-4).

\paragraph{Ejercicio 2 (análisis; linealizar un no lineal).} Para el oscilador de Duffing amortiguado $\dot{x}=y$, $\dot{y}=x-x^3-0.2\,y$, halla los tres puntos fijos, escribe la jacobiana \eqref{eq:def_jacobiano} y clasifica cada equilibrio.
\aside{Pista: $y=0$ y $x-x^3=0\Rightarrow x=0,\pm1$; $J=\begin{pmatrix}0&1\\1-3x^2&-0.2\end{pmatrix}$.}
% SOL: FP (0,0),(±1,0). En (0,0): Delta=-1<0 silla. En (±1,0): 1-3=-2 -> Delta=2, tau=-0.2, D=0.04-8<0 -> espirales estables. (Doble pozo amortiguado.)

\paragraph{Ejercicio 3 (interpretación; nodo vs foco).} Un equilibrio estable puede relajarse de forma monótona (nodo) o con oscilaciones amortiguadas (foco). Explica, en términos de cruzar la parábola $\Delta=\tfrac14\tau^2$, qué cambio físico convierte una recuperación monótona en una con sobreimpulso, y da un ejemplo biológico donde importe la diferencia.
\aside{Pista: por debajo de la parábola, autovalores reales (nodo); por encima, complejos (foco con $\omega=\tfrac12\sqrt{4\Delta-\tau^2}$).}
% SOL: Aumentar Delta (o reducir |tau|) cruza la parabola: reales->complejos, nodo->foco. Ejemplo: recuperacion de hemoglobina tras hemorragia con o sin overshoot (cf. cap. 20); o glucosa-insulina.

% ============================================================
% FUENTES: Notas Biomate 1; raw_plain_nonlinear_dynamics.txt.
% Strogatz (2015), caps. 5 y 6; Edelstein-Keshet (2007).
% ============================================================
\section{Planos de fases y nullclines}

En el capítulo anterior establecimos las herramientas algebraicas para caracterizar el comportamiento local de un sistema en la vecindad inmediata de sus puntos de equilibrio mediante la linealización jacobiana. Sin embargo, para comprender la dinámica global de un sistema bidimensional ---cómo se conectan las cuencas de atracción, hacia dónde escapan las trayectorias y cómo se organizan los flujos a gran escala--- necesitamos métodos geométricos no lineales.

El método más potente y visual para desentrañar la geometría global del espacio de estados en dos dimensiones es el análisis mediante \textbf{isoclinas nulas} o \textbf{nullclines}~\cite{strogatz2015,edelstein2007}.

\subsection{El método de las isoclinas nulas (\emph{Nullclines})}

Consideremos un sistema bidimensional autónomo acoplado:
\begin{equation}
\begin{aligned}
    \dot{x} &= f(x, y), \\
    \dot{y} &= g(x, y).
\end{aligned}
\label{eq:sistema_2d_general_nullclines}
\end{equation}
En cada punto $(x,y)$ del plano, el campo vectorial $\vec{v}(x,y) = (\dot{x}, \dot{y})^\top$ define una dirección y una velocidad instantánea. Las trayectorias del sistema son tangentes en todo instante a este campo.

\begin{definition}[Isoclinas nulas o Nullclines]
Las \textbf{nullclines} son las curvas geométricas en el plano de fases a lo largo de las cuales una de las componentes del campo vectorial se anula:
\begin{itemize}
    \item \textbf{$x$-nullcline:} Es el conjunto de puntos donde $\dot{x} = f(x,y) = 0$. Sobre esta curva, la velocidad horizontal es nula; por ende, el flujo del sistema es \textbf{estrictamente vertical} ($\uparrow$ o $\downarrow$).
    \item \textbf{$y$-nullcline:} Es el conjunto de puntos donde $\dot{y} = g(x,y) = 0$. Sobre esta curva, la velocidad vertical es nula; por ende, el flujo del sistema es \textbf{estrictamente horizontal} ($\rightarrow$ o $\leftarrow$).
\end{itemize}
\end{definition}

Las nullclines ofrecen dos ventajas fundamentales para el análisis cualitativo:
\begin{enumerate}
    \item \textbf{Localización inmediata de equilibrios:} Un punto fijo $(x^*, y^*)$ satisface simultáneamente $\dot{x} = 0$ y $\dot{y} = 0$. Por consiguiente, los puntos fijos del sistema son precisamente las \textbf{intersecciones} entre las $x$-nullclines y las $y$-nullclines.
    \item \textbf{Partición del plano en regiones de signo constante:} Las nullclines dividen el plano de fases en sectores disjuntos donde los signos de $f(x,y)$ y $g(x,y)$ permanecen invariantes. En cada sector, el campo vectorial apunta consistentemente hacia uno de los cuatro cuadrantes direccionales:
    \begin{align*}
        (+,+) &\implies \text{noreste } (\nearrow), &
        (-,+) &\implies \text{noroeste } (\nwarrow), \\
        (-,-) &\implies \text{suroeste } (\swarrow), &
        (+,-) &\implies \text{sureste } (\searrow).
    \end{align*}
\end{enumerate}

\begin{figure}[htbp]
\centering
\begin{tikzpicture}[>=Stealth, scale=0.9]
    % Ejes coordenados
    \draw[->, thick, black!60] (-4,0) -- (4,0) node[right, font=\small] {$x$};
    \draw[->, thick, black!60] (0,-3.2) -- (0,3.2) node[above, font=\small] {$y$};

    % x-nullcline (línea discontinua)
    \draw[very thick, black!85, dashed, domain=-3.5:3.5, samples=60] 
        plot (\x, {0.2*\x*\x - 1.2}) node[right, font=\footnotesize] {$f(x,y)=0$ ($x$-nullcline)};
    
    % y-nullcline (línea punteada)
    \draw[very thick, black!85, dotted, domain=-3.5:3.5, samples=60] 
        plot (\x, {-0.3*\x + 0.8}) node[above right, font=\footnotesize] {$g(x,y)=0$ ($y$-nullcline)};

    % Puntos fijos de intersección
    \filldraw[black] (2.5, 0.05) circle (2.5pt) node[below right=2pt, font=\footnotesize] {$(x_1^*, y_1^*)$};
    \filldraw[black] (-3.8, 1.94) circle (2.5pt) node[above left=2pt, font=\footnotesize] {$(x_2^*, y_2^*)$};

    % Vectores sobre x-nullcline (flujo puramente vertical)
    \draw[->, very thick, black!80] (1.0, -1.0) -- (1.0, -0.4);
    \draw[->, very thick, black!80] (3.2, 0.85) -- (3.2, 0.25);
    \node[font=\tiny, black!70, right] at (1.0, -0.7) {$\dot{x}=0$ (flujo vertical)};

    % Vectores sobre y-nullcline (flujo puramente horizontal)
    \draw[->, very thick, black!80] (-1.5, 1.25) -- (-0.9, 1.25);
    \draw[->, very thick, black!80] (1.0, 0.5) -- (1.6, 0.5);
    \node[font=\tiny, black!70, above] at (-1.2, 1.35) {$\dot{y}=0$ (flujo horizontal)};

    % Flechas direccionales en regiones
    \node[font=\large] at (0, -2.0) {$\nearrow$};
    \node[font=\large] at (0, 1.8) {$\nwarrow$};
    \node[font=\large] at (3.0, -1.5) {$\searrow$};
    \node[font=\large] at (-2.5, 0) {$\nearrow$};
\end{tikzpicture}
\caption{Esquema conceptual del método de las nullclines. Las curvas $f(x,y)=0$ (línea discontinua) y $g(x,y)=0$ (línea punteada) dividen el plano en regiones de flujo homogéneo; sus intersecciones corresponden a los equilibrios del sistema.}
\label{fig:esquema_nullclines}
\end{figure}

\subsection{El péndulo no lineal: análisis global en el plano de fases}

Retomemos el problema fundamental del péndulo simple sin amortiguamiento presentado en el capítulo 3. La ecuación diferencial de segundo orden que gobierna la dinámica del ángulo $\theta(t)$ respecto a la vertical es:
\begin{equation}
    \ddot{\theta} + \omega_0^2 \sin\theta = 0
    \label{eq:pendulo_segundo_orden}
\end{equation}
donde $\omega_0 = \sqrt{g/L}$ es la frecuencia natural de oscilación para pequeñas amplitudes.

Definiendo la velocidad angular $v \equiv \dot{\theta}$, convertimos la ecuación de segundo orden en un sistema de dos ecuaciones diferenciales de primer orden en el plano de fases $(\theta, v)$:
\begin{equation}
\begin{aligned}
    \dot{\theta} &= v, \\
    \dot{v} &= -\omega_0^2 \sin\theta.
\end{aligned}
\label{eq:sistema_pendulo_plano}
\end{equation}

\paragraph{Nullclines del péndulo.}
Las isoclinas nulas son inmediatamente identificables:
\begin{itemize}
    \item \textbf{$\theta$-nullcline ($\dot{\theta} = 0$):} Ocurre cuando $v = 0$ (el eje horizontal $\theta$). Sobre este eje, las trayectorias cruzan en dirección puramente vertical: hacia abajo si $\sin\theta > 0$ y hacia arriba si $\sin\theta < 0$.
    \item \textbf{$v$-nullcline ($\dot{v} = 0$):} Ocurre cuando $\sin\theta = 0$, lo que implica $\theta = k\pi$ para todo $k \in \mathbb{Z}$ (una familia infinita de rectas verticales equidistantes). Sobre estas rectas, el flujo es puramente horizontal: hacia la derecha si $v > 0$ y hacia la izquierda si $v < 0$.
\end{itemize}

\paragraph{Puntos fijos y linealización local.}
Los equilibrios resultan de la intersección entre $v=0$ y $\sin\theta=0$:
\begin{equation}
    (\theta^*, v^*) = (k\pi, 0), \quad k \in \mathbb{Z}.
\end{equation}
La matriz jacobiana del sistema evaluada en un punto general $(\theta, v)$ es:
\begin{equation}
    J(\theta, v) = \begin{pmatrix}
        \dfrac{\partial \dot{\theta}}{\partial \theta} & \dfrac{\partial \dot{\theta}}{\partial v} \\[6pt]
        \dfrac{\partial \dot{v}}{\partial \theta} & \dfrac{\partial \dot{v}}{\partial v}
    \end{pmatrix} = \begin{pmatrix} 0 & 1 \\ -\omega_0^2 \cos\theta & 0 \end{pmatrix}.
\end{equation}
Distingamos las dos clases de equilibrios:
\begin{enumerate}
    \item \textbf{Equilibrio inferior colgado ($\theta^* = 2k\pi$, $v^* = 0$):}
    Aquí $\cos(2k\pi) = +1$, de modo que
    \begin{equation}
        J(2k\pi, 0) = \begin{pmatrix} 0 & 1 \\ -\omega_0^2 & 0 \end{pmatrix}, \quad \tau = 0, \quad \Delta = \omega_0^2 > 0.
    \end{equation}
    Los autovalores son puramente imaginarios: $\lambda_{1,2} = \pm i\omega_0$. La linealización predice un \textbf{centro}, correspondiente a pequeñas oscilaciones armónicas elípticas de periodo $T_0 = 2\pi/\omega_0$.
    
    \item \textbf{Equilibrio superior invertido ($\theta^* = (2k+1)\pi$, $v^* = 0$):}
    Aquí $\cos((2k+1)\pi) = -1$, de modo que
    \begin{equation}
        J((2k+1)\pi, 0) = \begin{pmatrix} 0 & 1 \\ \omega_0^2 & 0 \end{pmatrix}, \quad \tau = 0, \quad \Delta = -\omega_0^2 < 0.
    \end{equation}
    Los autovalores son reales y opuestos: $\lambda_{1,2} = \pm \omega_0$. El equilibrio es una \textbf{silla (\emph{saddle point}) inestable}. Las direcciones propias asociadas a $\lambda = \pm \omega_0$ corresponden a los vectores $\vec{v}_{\pm} = (1, \pm\omega_0)^\top$.
\end{enumerate}

\subsection{Libraciones, rotaciones y la separatrix}

Multiplicando la ecuación $\ddot{\theta} + \omega_0^2 \sin\theta = 0$ por $\dot{\theta}$ e integrando respecto al tiempo, obtenemos la ley de conservación de la energía mecánica total:
\begin{equation}
    \frac{d}{dt}\left[ \frac{1}{2}v^2 - \omega_0^2 \cos\theta \right] = v \dot{v} + \omega_0^2 \sin\theta \, v = v(-\omega_0^2 \sin\theta) + \omega_0^2 \sin\theta \, v = 0.
\end{equation}
Por tanto, la energía adimensionalizada por unidad de masa $E$ es una constante del movimiento:
\begin{equation}
    E(\theta, v) = \frac{1}{2}v^2 - \omega_0^2 \cos\theta = \text{constante}.
    \label{eq:energia_pendulo}
\end{equation}
Las curvas de nivel $E(\theta, v) = \text{const}$ en el plano $(\theta, v)$ corresponden a las trayectorias exactas del sistema. El valor de la energía clasifica globalmente los movimientos posibles:

\begin{figure}[htbp]
\centering
\begin{tikzpicture}[>=Stealth, scale=1.0]
    \begin{axis}[
        xmin=-4.2, xmax=4.2, ymin=-3.2, ymax=3.2,
        axis lines=middle,
        xlabel={$\theta$ (ángulo)}, ylabel={$v$ (velocidad angular)},
        xlabel style={at={(ticklabel* cs:1)},anchor=west},
        ylabel style={at={(ticklabel* cs:1)},anchor=south},
        xtick={-3.1416, 0, 3.1416},
        xticklabels={$-\pi$, $0$, $\pi$},
        ytick={-2, 0, 2},
        yticklabels={$-2\omega_0$, $0$, $2\omega_0$},
        width=0.88\linewidth, height=6.5cm,
        grid=both, grid style={gray!20}
    ]
        % Libraciones (órbitas cerradas para E < omega0^2)
        \addplot[thick, black!85, domain=-1.04:1.04, samples=80] {sqrt(2*(cos(deg(x)) - 0.5))};
        \addplot[thick, black!85, domain=-1.04:1.04, samples=80] {-sqrt(2*(cos(deg(x)) - 0.5))};

        \addplot[thick, black!85, domain=-2.09:2.09, samples=100] {sqrt(2*(cos(deg(x)) + 0.5))};
        \addplot[thick, black!85, domain=-2.09:2.09, samples=100] {-sqrt(2*(cos(deg(x)) + 0.5))};

        % Separatriz (E = omega0^2 => v = \pm 2 cos(theta/2))
        \addplot[very thick, black!95, domain=-3.1415:3.1415, samples=120] {2*cos(deg(x/2))};
        \addplot[very thick, black!95, domain=-3.1415:3.1415, samples=120] {-2*cos(deg(x/2))};

        % Rotaciones superiores (E = 2.0 omega0^2 => v = + \sqrt{2*(cos(theta) + 2)})
        \addplot[thick, black!65, domain=-4.2:4.2, samples=120] {sqrt(2*(cos(deg(x)) + 2.0))};
        \addplot[thick, black!65, domain=-4.2:4.2, samples=120] {sqrt(2*(cos(deg(x)) + 3.2))};

        % Rotaciones inferiores (v < 0)
        \addplot[thick, black!65, domain=-4.2:4.2, samples=120] {-sqrt(2*(cos(deg(x)) + 2.0))};
        \addplot[thick, black!65, domain=-4.2:4.2, samples=120] {-sqrt(2*(cos(deg(x)) + 3.2))};

        % Flechas de dirección sobre la separatriz
        \draw[->, very thick] (axis cs:0, 2) -- (axis cs:0.3, 2);
        \draw[->, very thick] (axis cs:0, -2) -- (axis cs:-0.3, -2);
        \draw[->, thick] (axis cs:0, 1.22) -- (axis cs:0.3, 1.22);
        \draw[->, thick] (axis cs:0, -1.22) -- (axis cs:-0.3, -1.22);
        \draw[->, thick] (axis cs:0, 2.45) -- (axis cs:0.4, 2.45);
        \draw[->, thick] (axis cs:0, -2.45) -- (axis cs:-0.4, -2.45);

        % Puntos fijos
        \filldraw[black] (axis cs:0,0) circle (2.2pt) node[above right=1pt, font=\tiny] {Centro $(0,0)$};
        \draw[black, fill=white, thick] (axis cs:3.1416,0) circle (2.2pt) node[above right=1pt, font=\tiny] {Silla $(\pi,0)$};
        \draw[black, fill=white, thick] (axis cs:-3.1416,0) circle (2.2pt) node[above left=1pt, font=\tiny] {Silla $(-\pi,0)$};

        % Etiquetas de régimen
        \node[font=\footnotesize, align=center] at (axis cs:0, 0.4) {Libraciones\\($E < \omega_0^2$)};
        \node[font=\footnotesize, align=center] at (axis cs:0, 2.8) {Rotaciones hacia la derecha ($v > 0$)};
        \node[font=\footnotesize, align=center] at (axis cs:0, -2.8) {Rotaciones hacia la izquierda ($v < 0$)};
        \node[font=\footnotesize, fill=white, inner sep=1pt] at (axis cs:1.8, 1.5) {Separatriz ($E = \omega_0^2$)};
    \end{axis}
\end{tikzpicture}
\caption{Retrato de fases global del péndulo no lineal en el plano $(\theta, v)$. Las curvas cerradas interiores representan libraciones periódicas; las curvas ondulantes exteriores representan rotaciones completas; y la curva gruesa es la \textbf{separatriz}, que conecta las sillas inestables en $\pm\pi$.}
\label{fig:retrato_fases_pendulo_completo}
\end{figure}

\begin{enumerate}
    \item \textbf{Libraciones ($-\omega_0^2 < E < \omega_0^2$):}
    La energía cinética no es suficiente para que el péndulo alcance la cúspide vertical $\theta = \pm \pi$. El movimiento queda confinado a un intervalo de oscilación $[-\theta_{\max}, \theta_{\max}]$, donde $\cos\theta_{\max} = -E/\omega_0^2$. Las órbitas son curvas cerradas periódicas alrededor del centro $(0,0)$. Para oscilaciones grandes, el periodo $T(E)$ crece monótonamente debido a los efectos no lineales, divergiendo cuando $E \to \omega_0^2$.
    
    \item \textbf{Rotaciones ($E > \omega_0^2$):}
    El péndulo posee energía cinética sobrante en el punto más alto ($v(\pm\pi) = \pm\sqrt{2(E - \omega_0^2)} \neq 0$). Por ende, el ángulo $\theta(t)$ crece o decrece indefinidamente, girando en vueltas continuas sobre su eje sin detenerse jamás.
    
    \item \textbf{La separatriz ($E = \omega_0^2$):}
    Es la curva de nivel crítica que separa geométricamente las oscilaciones cerradas (libraciones) de las rotaciones abiertas. Despejando $v$ de la \cref{eq:energia_pendulo} con $E = \omega_0^2$:
    \begin{equation}
        \frac{1}{2}v^2 - \omega_0^2 \cos\theta = \omega_0^2 \implies v(\theta) = \pm 2\omega_0 \cos\left(\frac{\theta}{2}\right).
        \label{eq:separatriz_formula}
    \end{equation}
    Esta trayectoria parte en $t \to -\infty$ desde el equilibrio inestable invertido en $-\pi$ y llega asintóticamente en $t \to +\infty$ al equilibrio inestable en $+\pi$. Como $\theta$ es una variable angular en el círculo $\mathbb{S}^1$, los puntos $\theta = -\pi$ y $\theta = +\pi$ representan el mismo estado físico; por tanto, la separatriz es una \textbf{órbita homoclínica} que tarda un tiempo infinito en recorrerse\aside{Si concebimos el espacio de estados no sobre el plano $\mathbb{R}^2$ sino sobre la variedad cilíndrica natural $\mathbb{S}^1 \times \mathbb{R}$, las rotaciones son curvas cerradas no contraíbles que envuelven el cilindro, mientras que las libraciones son lazos cerrados topológicamente contraíbles a un punto.}.
\end{enumerate}

% ============================================================
% FUENTES: Notas Biomate 1; raw_plain_nonlinear_dynamics.txt.
% Strogatz (2015), caps. 6 y 7; Murray (2002); Edelstein-Keshet (2007).
% ============================================================
\section{Sistemas conservativos y teorema de Poincaré--Bendixson}

En el estudio de los sistemas dinámicos en dos dimensiones, surge una distinción física y matemática de primer orden: la frontera entre los sistemas \textbf{conservativos} (comunes en la mecánica clásica sin fricción) y los sistemas \textbf{disipativos} (universales en la biología, la química y la neurociencia)~\cite{strogatz2015,murray2002}.

En este capítulo analizaremos el papel de las funciones de energía como integrales primeras para certificar la existencia de centros no lineales y presentaremos uno de los resultados más célebres y profundos de la topología dinámica: el \textbf{Teorema de Poincaré--Bendixson}.

\subsection{Sistemas conservativos e integrales primeras}

Consideremos un sistema bidimensional autónomo $\dot{\vec{x}} = \vec{f}(\vec{x})$, con $\vec{x} = (x,y)^\top \in \mathbb{R}^2$.

\begin{definition}[Cantidad conservada / Integral primera]
Una función continuamente diferenciable de valores reales $E(x,y)$ es una \textbf{cantidad conservada} (o integral primera) del sistema si no es constante en ningún conjunto abierto y su derivada temporal a lo largo de cualquier trayectoria solución es idénticamente nula:
\begin{equation}
    \dot{E} = \frac{d}{dt}E(x(t), y(t)) = \nabla E \cdot \dot{\vec{x}} = \frac{\partial E}{\partial x} \dot{x} + \frac{\partial E}{\partial y} \dot{y} = 0.
    \label{eq:derivada_energia_cero}
\end{equation}
Un sistema que admite una cantidad conservada se denomina \textbf{sistema conservativo}.
\end{definition}

La existencia de una integral primera $E(x,y)$ reduce en una unidad los grados de libertad efectivos del sistema: cada trayectoria individual $\vec{x}(t)$ queda estrictamente confinada a una curva de nivel $E(x,y) = C$, donde la constante $C = E(x(0), y(0))$ está fijada por la condición inicial.

\begin{theorem}[Criterio no lineal de centros conservativos]
\label{thm:centros_conservativos}
Sea $\dot{\vec{x}} = \vec{f}(\vec{x})$ un sistema conservativo con energía $E(x,y)$ y sea $\vec{x}^*$ un punto fijo aislado. Si $\vec{x}^*$ es un \textbf{mínimo o máximo local estricto} de $E(x,y)$, entonces $\vec{x}^*$ es un \textbf{centro no lineal}: está rodeado por una familia continua de órbitas periódicas cerradas cerrándose concéntricamente sobre él.
\end{theorem}

\begin{theorem}[Imposibilidad de atractores en sistemas conservativos]
En un sistema conservativo continuo en $\mathbb{R}^2$, \textbf{no pueden existir atractores}: ningún punto fijo puede ser un nodo asintóticamente estable ni un foco atractor, y no pueden existir ciclos límite aislados.
\end{theorem}

La intuición geométrica es transparente: si un punto $\vec{x}^*$ fuese un atractor asintótico, todas las trayectorias de una vecindad $U$ convergerían hacia él cuando $t \to \infty$. Pero a lo largo de cada trayectoria $E(\vec{x}(t))$ es estrictamente constante ($E(x_0, y_0)$); por continuidad, si todas convergieran a $\vec{x}^*$, la energía $E$ tendría que ser idéntica en toda la vecindad $U$, contradiciendo la definición de que $E$ no es constante.

\subsection{Centros vs Focos: el fallo de la linealización no hiperbólica}

Recordemos del capítulo 15 que el Teorema de Hartman--Grobman garantiza la fidelidad topológica de la linealización jacobiana únicamente para equilibrios \textbf{hiperbólicos} ($\operatorname{Re}(\lambda_i) \neq 0$). Cuando la matriz jacobiana tiene autovalores puramente imaginarios $\lambda = \pm i\omega$ ($\tau = 0, \Delta > 0$), la linealización predice un centro lineal, pero el equilibrio no es hiperbólico.

¿Qué ocurre en el sistema no lineal completo? Los términos no lineales de orden superior pueden alterar radicalmente el retrato de fases local:

\begin{example}[Destrucción no lineal de un centro lineal]
Consideremos el sistema en coordenadas polares $(r, \theta)$:
\begin{equation}
\begin{aligned}
    \dot{r} &= a r^3, \\
    \dot{\theta} &= 1,
\end{aligned}
\end{equation}
con $a \in \mathbb{R}$. En coordenadas cartesianas $(x = r\cos\theta, y = r\sin\theta)$, el sistema linealizado en el origen es:
\begin{equation}
    \begin{pmatrix} \dot{x} \\ \dot{y} \end{pmatrix} = \begin{pmatrix} 0 & -1 \\ 1 & 0 \end{pmatrix} \begin{pmatrix} x \\ y \end{pmatrix},
\end{equation}
cuya matriz jacobiana tiene traza $\tau = 0$ y determinante $\Delta = 1 > 0$, prediciendo un centro lineal perfecto para cualquier valor de $a$.

Sin embargo, la ecuación radial $\dot{r} = a r^3$ revela la dinámica real:
\begin{itemize}
    \item Si $a < 0$: $\dot{r} < 0$ para todo $r > 0$. El radio decae monótonamente a cero ($r(t) \sim t^{-1/2}$); el origen es en realidad un \textbf{foco o espiral asintóticamente estable} (atractor débil).
    \item Si $a > 0$: $\dot{r} > 0$ para todo $r > 0$. Las trayectorias se alejan en espiral creciente; el origen es un \textbf{foco inestable} (repulsor).
    \item Si $a = 0$: $\dot{r} = 0$, las órbitas son circunferencias cerradas concéntricas $r(t) = r_0$; el origen es un \textbf{centro verdadero}.
\end{itemize}
\end{example}

Por tanto, un centro lineal predicho por la matriz jacobiana es una condición estructuralmente marginal. Se requiere una propiedad global ---como la existencia demostrada de una integral primera $E(x,y)$ mediante el \cref{thm:centros_conservativos} o una simetría temporal reversible--- para confirmar que el equilibrio no lineal es un centro genuino.

\subsection{Sistemas disipativos y atractores biológicos}

A diferencia de los planetas orbitando en el vacío o los péndulos ideales sin rozamiento, los sistemas biológicos son sistemas termodinámicamente abiertos que intercambian continuamente materia y energía con su entorno. En estos sistemas el volumen del espacio de fases no se conserva, sino que se contrae en promedio: son \textbf{sistemas disipativos}~\cite{edelstein2007}.

En un sistema disipativo, los centros neutros son aberraciones que no sobreviven a perturbaciones realistas. En su lugar, emergen \textbf{atractores}: estructuras asintóticas robustas hacia las cuales el sistema converge sin importar pequeñas fluctuaciones en las condiciones iniciales.

\subsection{El Teorema de Poincaré--Bendixson}

El resultado culminante para la dinámica en el plano continuo $\mathbb{R}^2$ es el teorema demostrado por Henri Poincaré en 1892 y refinado por Ivar Bendixson en 1901:

\begin{theorem}[Teorema de Poincaré--Bendixson (formulación intuitiva)]
\label{thm:poincare_bendixson}
Sea $\dot{\vec{x}} = \vec{f}(\vec{x})$ un sistema dinámico autónomo continuo con campo vectorial $\mathcal{C}^1$ en el plano $\mathbb{R}^2$. Sea $R \subset \mathbb{R}^2$ una región cerrada y acotada (\textbf{compacta}) que satisface dos condiciones:
\begin{enumerate}
    \item $R$ es una \textbf{región de atrapamiento} (\emph{trapping region}): en todos los puntos de la frontera $\partial R$, el campo vectorial apunta hacia el interior de $R$ (o es tangente), de modo que ninguna trayectoria que entre a $R$ puede salir de ella jamás.
    \item $R$ \textbf{no contiene ningún punto fijo de equilibrio} ($\vec{f}(\vec{x}) \neq \vec{0}$ para todo $\vec{x} \in R$).
\end{enumerate}
Entonces, toda trayectoria que inicie dentro de $R$ permanece en $R$ para todo $t \ge 0$ y, cuando $t \to \infty$, \textbf{debe ser una órbita periódica cerrada o converger asintóticamente hacia un ciclo límite periódico contenido en $R$}.
\end{theorem}

\begin{figure}[htbp]
\centering
\begin{tikzpicture}[>=Stealth, scale=0.95]
    % Región anular sombreada
    \fill[black!8] (0,0) circle (3.2cm);
    \fill[white] (0,0) circle (1.1cm);

    % Frontera exterior e interior
    \draw[very thick, black!85, dashed] (0,0) circle (3.2cm);
    \draw[very thick, black!85, dashed] (0,0) circle (1.1cm);

    % Punto fijo inestable en el origen
    \draw[black, fill=white, thick] (0,0) circle (2.5pt) node[below=4pt, font=\tiny] {$\vec{x}^*$ (repulsor inestable)};

    % Vectores de flujo entrante en la frontera exterior
    \foreach \ang in {0, 45, 90, 135, 180, 225, 270, 315} {
        \draw[->, very thick, black!75] (\ang:3.2cm) -- (\ang:2.6cm);
    }
    \node[font=\footnotesize, black!85] at (0, 3.5) {Frontera exterior: flujo entra a $R$};

    % Vectores de flujo saliente del núcleo inestable en la frontera interior
    \foreach \ang in {20, 65, 110, 155, 200, 245, 290, 335} {
        \draw[->, very thick, black!75] (\ang:1.1cm) -- (\ang:1.6cm);
    }
    \node[font=\tiny, black!85, align=center] at (0, -0.6) {Flujo expulsado\\hacia $R$};

    % Ciclo límite periódico atractor dentro de R
    \draw[very thick, black!95] (0,0) circle (2.1cm);
    \node[font=\bfseries\footnotesize, black!90] at (0, 2.3) {Ciclo límite periódico $\Gamma$};

    % Trayectoria espiral exterior convergiendo a Gamma
    \draw[thick, black!70, domain=0:9.42, samples=150, variable=\t] 
        plot ({ (2.1 + 0.9*exp(-0.35*\t))*cos(\t r) }, { (2.1 + 0.9*exp(-0.35*\t))*sin(\t r) });
    \draw[->, thick, black!70] ({ (2.1 + 0.9*exp(-0.35*2.5))*cos(2.5 r) }, { (2.1 + 0.9*exp(-0.35*2.5))*sin(2.5 r) }) 
        -- ({ (2.1 + 0.9*exp(-0.35*2.7))*cos(2.7 r) }, { (2.1 + 0.9*exp(-0.35*2.7))*sin(2.7 r) });

    % Trayectoria espiral interior convergiendo a Gamma
    \draw[thick, black!70, domain=0:9.42, samples=150, variable=\t] 
        plot ({ (2.1 - 0.7*exp(-0.35*\t))*cos(\t r) }, { (2.1 - 0.7*exp(-0.35*\t))*sin(\t r) });

    % Etiqueta de la región de atrapamiento
    \node[font=\bfseries\small] at (-2.4, -1.8) {Región de atrapamiento $R$};
\end{tikzpicture}
\caption{Construcción clásica del Teorema de Poincaré--Bendixson mediante un anillo de atrapamiento (\emph{trapping annulus}). El campo vectorial apunta hacia el interior en ambas fronteras de $R$; al no existir equilibrios en $R$, las trayectorias quedan atrapadas y convergen forzosamente hacia el ciclo límite $\Gamma$.}
\label{fig:poincare_bendixson_annulus}
\end{figure}

\paragraph{Imposibilidad topológica de caos en flujos 2D.}
Una consecuencia trascendental del Teorema de Poincaré--Bendixson es la siguiente:
\begin{theorem}[No existencia de caos en $\mathbb{R}^2$ continuo]
Un sistema de ecuaciones diferenciales autónomas continuas $\dot{\vec{x}} = \vec{f}(\vec{x})$ con $f \in \mathcal{C}^1$ en dos dimensiones ($n = 2$) \textbf{jamás puede exhibir caos determinista ni atractores extraños}.
\end{theorem}

La demostración reside en la combinación de dos principios:
\begin{enumerate}
    \item Por el teorema de existencia y unicidad de Picard--Lindelöf, dos trayectorias en el plano de fases no pueden cruzarse jamás en un tiempo finito.
    \item En dos dimensiones, una curva cerrada simple divide el plano en dos regiones desconectadas (Teorema de la curva de Jordan). Una trayectoria confinada en una región compacta sin puntos fijos no tiene espacio topológico para doblarse, entrelazarse y estirarse sin autointersecarse; está forzada a enrolarse asintóticamente sobre una curva cerrada periódica.
\end{enumerate}

Para que un flujo autónomo continuo produzca trayectorias caóticas con sensibilidad a condiciones iniciales (como el atractor de Lorenz o el atractor de Rössler), se requieren forzosamente \textbf{al menos tres dimensiones} ($n \ge 3$) en el espacio de fases\aside{Esta es la razón profunda por la que el sistema atmosférico de Lorenz simplificado requirió tres variables dinámicas acopladas $(X, Y, Z)$ para generar su célebre mariposa caótica.}.

% ============================================================
% FUENTES: body/body.tex; Notas Biomate 1; raw_plain_nonlinear_dynamics.txt.
% Murray (2002); Strogatz (2015), caps. 6 y 7; Edelstein-Keshet (2007).
% ============================================================
\section{El modelo de Lotka--Volterra y ciclos límite}

En este capítulo analizaremos el modelo clásico de interacción biológica: el sistema depredador--presa de Lotka--Volterra~\cite{volterra1926,murray2002}. A partir de su estructura matemática descubriremos la fragilidad de los centros conservativos y daremos el paso hacia los \textbf{ciclos límite}, que constituyen el mecanismo real tras los ritmos biológicos, las oscilaciones neuronales y el latido cardíaco.

\subsection{El modelo depredador--presa de Lotka--Volterra}

Propuesto independientemente por Alfred J. Lotka (1925) en el contexto de reacciones químicas oscilantes y por Vito Volterra (1926) para explicar las fluctuaciones estadísticas de peces en el mar Adriático tras la Primera Guerra Mundial, el modelo describe la interacción entre una población de presas $x(t) \ge 0$ y una de depredadores $y(t) \ge 0$:
\begin{align}
  \frac{dx}{dt} &= \alpha\, x - \beta\, x y, \label{eq:lv-x}\\[2pt]
  \frac{dy}{dt} &= -\gamma\, y + \delta\, x y, \label{eq:lv-y}
\end{align}
donde los parámetros $\alpha, \beta, \gamma, \delta > 0$ tienen interpretaciones biológicas directas:
\begin{itemize}
    \item $\alpha$: Tasa de crecimiento malthusiano exponencial de las presas en ausencia de depredadores.
    \item $\beta$: Tasa de depredación por encuentro (interacción de ley de acción de masas $xy$).
    \item $\gamma$: Tasa de mortalidad natural de los depredadores en ausencia de alimento.
    \item $\delta$: Eficiencia de conversión de presas consumidas en nuevos depredadores.
\end{itemize}

\paragraph{Nullclines y puntos fijos.}
Las isoclinas nulas del sistema en el primer cuadrante son:
\begin{itemize}
    \item \textbf{$x$-nullclines ($\dot{x} = x(\alpha - \beta y) = 0$):} Consisten en dos rectas: el eje $x = 0$ (extinción de presas) y la recta horizontal $y = \alpha/\beta$.
    \item \textbf{$y$-nullclines ($\dot{y} = y(-\gamma + \delta x) = 0$):} Consisten en dos rectas: el eje $y = 0$ (extinción de depredadores) y la recta vertical $x = \gamma/\delta$.
\end{itemize}

\begin{definition}[Equilibrios del sistema]
El sistema~\eqref{eq:lv-x}--\eqref{eq:lv-y} admite dos puntos de equilibrio en el cuadrante no negativo:
\begin{enumerate}
    \item El \textbf{equilibrio trivial de extinción total} $(0,0)$.
    \item El \textbf{equilibrio de coexistencia no trivial}:
    \begin{equation}
        (x^{*}, y^{*}) = \left(\frac{\gamma}{\delta},\, \frac{\alpha}{\beta}\right).
        \label{eq:lv_coexistencia}
    \end{equation}
\end{enumerate}
\end{definition}

\begin{theorem}[Estabilidad lineal del equilibrio de coexistencia]
\label{thm:lv_estabilidad}
Para cualquier combinación de parámetros estrictamente positivos $\alpha, \beta, \gamma, \delta > 0$, la matriz jacobiana del sistema evaluada en el punto de coexistencia $(x^{*}, y^{*})$ tiene traza nula $\tau = 0$ y determinante positivo $\Delta = \alpha\gamma > 0$. Sus autovalores son puramente imaginarios:
\begin{equation}
    \lambda_{\pm} = \pm i \sqrt{\alpha \gamma}.
\end{equation}
El equilibrio linealizado es un centro con frecuencia angular $\omega_0 = \sqrt{\alpha\gamma}$.
\end{theorem}

\begin{proof}
La matriz jacobiana general es:
\begin{equation}
    J(x,y) = \begin{pmatrix} \alpha - \beta y & -\beta x \\ \delta y & -\gamma + \delta x \end{pmatrix}.
\end{equation}
Evaluando en el punto de coexistencia $(x^*, y^*) = (\gamma/\delta, \alpha/\beta)$:
\begin{equation}
    J(x^{*}, y^{*}) = \begin{pmatrix} \alpha - \beta(\alpha/\beta) & -\beta(\gamma/\delta) \\ \delta(\alpha/\beta) & -\gamma + \delta(\gamma/\delta) \end{pmatrix} = \begin{pmatrix} 0 & -\frac{\beta\gamma}{\delta} \\ \frac{\delta\alpha}{\beta} & 0 \end{pmatrix}.
\end{equation}
Calculando sus invariantes algebraicos:
\begin{align}
    \tau &= \operatorname{tr} J(x^*, y^*) = 0 + 0 = 0, \\
    \Delta &= \det J(x^*, y^*) = (0)(0) - \left(-\frac{\beta\gamma}{\delta}\right)\left(\frac{\delta\alpha}{\beta}\right) = \alpha\gamma > 0.
\end{align}
Por tanto, la ecuación característica es $\lambda^2 + \alpha\gamma = 0 \implies \lambda_{\pm} = \pm i\sqrt{\alpha\gamma}$.
\end{proof}

\subsection{Integral primera y órbitas cerradas en Lotka--Volterra}

Como aprendimos en el capítulo 17, el hecho de que la linealización prediga autovalores puramente imaginarios ($\operatorname{Re}(\lambda) = 0$) no basta para asegurar que el sistema no lineal posea órbitas periódicas cerradas. Sin embargo, el sistema de Lotka--Volterra posee una integral primera exacta.

Dividiendo la ecuación~\eqref{eq:lv-y} entre la ecuación~\eqref{eq:lv-x}:
\begin{equation}
    \frac{dy}{dx} = \frac{y(-\gamma + \delta x)}{x(\alpha - \beta y)} \implies \frac{\alpha - \beta y}{y}\, dy = \frac{-\gamma + \delta x}{x}\, dx.
\end{equation}
Separando variables e integrando:
\begin{equation}
    \int \left( \frac{\alpha}{y} - \beta \right) dy = \int \left( -\frac{\gamma}{x} + \delta \right) dx \implies \alpha \ln y - \beta y = -\gamma \ln x + \delta x - C.
\end{equation}
Reordenando los términos, definimos la función $V(x,y)$:
\begin{equation}
    V(x,y) = \delta x - \gamma \ln x + \beta y - \alpha \ln y = C.
    \label{eq:lv_primer_integral}
\end{equation}

\begin{theorem}[Órbitas cerradas globales en Lotka--Volterra]
La función $V(x,y)$ es una integral primera invariante a lo largo de las trayectorias ($\dot{V} = 0$). Además, $V(x,y)$ tiene un mínimo global estricto aislado en el punto de coexistencia $(x^*, y^*)$. Por el Teorema~\ref{thm:centros_conservativos}, el equilibrio $(x^*, y^*)$ es un \textbf{centro no lineal verdadero} rodeado por una familia continua y anidada de órbitas periódicas cerradas.
\end{theorem}

\begin{proof}
Calculando la derivada temporal de $V(x,y)$ a lo largo de cualquier solución:
\begin{align}
    \dot{V} &= \frac{\partial V}{\partial x}\dot{x} + \frac{\partial V}{\partial y}\dot{y} = \left(\delta - \frac{\gamma}{x}\right)(\alpha x - \beta x y) + \left(\beta - \frac{\alpha}{y}\right)(-\gamma y + \delta x y) \notag\\
    &= \frac{\delta x - \gamma}{x} x (\alpha - \beta y) + \frac{\beta y - \alpha}{y} y (\delta x - \gamma) \notag\\
    &= (\delta x - \gamma)(\alpha - \beta y) + (\beta y - \alpha)(\delta x - \gamma) = 0.
\end{align}
Dado que $\dot{V} \equiv 0$, el valor de $V$ se conserva idénticamente en el tiempo. Las curvas de nivel $V(x,y) = C$ son curvas cerradas convexas en el primer cuadrante (véase la \cref{fig:lv_phase}).
\end{proof}

\begin{figure}[htbp]
  \centering
  \begin{tikzpicture}
    \begin{axis}[
        view={0}{90}, domain=0.05:2.2, y domain=0.05:2.2,
        xmin=0, xmax=2.3, ymin=0, ymax=2.3,
        samples=13, width=0.75\linewidth,
        xlabel={$x$ (presas)}, ylabel={$y$ (depredadores)},
        axis lines=middle, grid=both, grid style={gray!15},
        xlabel near ticks, ylabel near ticks,
      ]
      \addplot3[gray!70, -latex, quiver={
          u={x - x*y}, v={-y + x*y}, scale arrows=0.055}] (x,y,0);
      \draw[red!70!black, thick, dashed] (axis cs:1,0) -- (axis cs:1,2.2)
        node[below right, font=\footnotesize]{$x^{*}=1$};
      \draw[blue!70!black, thick, dashed] (axis cs:0,1) -- (axis cs:2.2,1)
        node[above right, font=\footnotesize]{$y^{*}=1$};
      
      % Curvas de nivel cerradas de V(x,y) = x - ln x + y - ln y
      \draw[thick, black!85] (axis cs:1,1) circle [x radius=0.4, y radius=0.4];
      \draw[thick, black!85] (axis cs:1,1) circle [x radius=0.75, y radius=0.7];
      \draw[thick, black!85] (axis cs:1,1) circle [x radius=1.05, y radius=0.95];

      \fill (axis cs:1,1) circle (2.2pt) node[above right, font=\footnotesize]{$(x^{*},y^{*})$};
    \end{axis}
  \end{tikzpicture}
  \caption{Retrato de fases del modelo de Lotka--Volterra con $\alpha=\beta=\gamma=\delta=1$. Las líneas discontinuas representan las nullclines ($x^*=1, y^*=1$) y las curvas cerradas concéntricas son las trayectorias periódicas generadas por la integral primera $V(x,y) = C$.}
  \label{fig:lv_phase}
\end{figure}

\subsection{Simulación computacional: Python y NetLogo}

\paragraph{Integración numérica en Python.}
El sistema~\eqref{eq:lv-x}--\eqref{eq:lv-y} puede integrarse numéricamente de forma directa mediante esquemas estándar como Euler o Runge--Kutta:

\begin{minted}[fontsize=\small, framesep=4pt, frame=lines]{python}
import numpy as np

alpha, beta, gamma, delta = 1.0, 1.0, 1.0, 1.0

def lv(u, t):
    """Campo vectorial de Lotka-Volterra."""
    x, y = u
    return np.array([alpha*x - beta*x*y,
                     -gamma*y + delta*x*y])

dt, T = 0.01, 20.0        # paso temporal y horizonte
u = np.array([1.2, 0.6])  # condición inicial
for t in np.arange(0, T, dt):
    u = u + dt * lv(u, t)
print(f"estado final: x = {u[0]:.3f}, y = {u[1]:.3f}")
\end{minted}

\paragraph{Simulación basada en agentes en NetLogo.}
La misma dinámica poblacional puede explorarse a nivel individual en un modelo basado en agentes~\cite{netlogo}, donde cada presa y depredador es una tortuga discreta con reglas de búsqueda, alimentación y reproducción:

\begin{minted}[fontsize=\small, framesep=4pt, frame=lines]{lexer/netlogo.py:NetLogoLexer}
globals [presas-iniciales depredadores-iniciales]
turtles-own [energia]

to setup
  clear-all
  create-turtles presas-iniciales [
    setxy random-xcor random-ycor
    set color green  set energia 10
  ]
  create-turtles depredadores-iniciales [
    setxy random-xcor random-ycor
    set color red
  ]
  reset-ticks
end

to go
  ask turtles [
    ifelse color = red
      [ let presa min-one-of turtles with [color = green]
                        [distance-to self]
        if presa != nobody [ face presa  fd 1.2 ] ]
      [ fd 1  set energia energia - 1
        if energia < 2 [ die ] ]
  ]
  tick
end
\end{minted}

\paragraph{Conexión con tiempo discreto: mapa logístico y caos.}
Si en lugar de tiempo continuo se modela la población con generaciones discretas no superpuestas $N_{t+1} = r N_t (1 - N_t/K)$~\cite{may1976}, aparece el célebre mapa logístico cuadrático. Como demostró Robert May en 1976, al incrementar la tasa de crecimiento intrínseca el sistema pasa por una cascada infinita de duplicación de periodos y desemboca en caos determinista, ilustrando la riqueza que adquieren las no linealidades en dominios iterados.

\subsection{Centros vs Ciclos Límite: Estabilidad estructural}

A pesar de su elegancia histórica, el modelo de Lotka--Volterra adolece de una severa limitación desde el punto de vista biológico: es \textbf{estructuralmente inestable}~\cite{strogatz2015,murray2002}.

En un centro conservativo como el de Lotka--Volterra:
\begin{enumerate}
    \item La amplitud y el periodo de la oscilación dependen por completo de la condición inicial $(x_0, y_0)$.
    \item Si una perturbación ambiental transitoria desplaza al sistema (por ejemplo, una sequía que diezme momentáneamente las presas), el sistema salta a una órbita cerrada diferente y \textbf{jamás regresa a la oscilación original}. No existe ninguna fuerza restauradora sobre la amplitud.
    \item Cualquier modificación biológicamente realista ---como agregar una capacidad de carga $K$ para las presas ($\dot{x} = \alpha x (1 - x/K) - \beta x y$) o una respuesta funcional de saturación de Holling tipo II--- destruye inmediatamente el centro, convirtiéndolo en un foco espiral amortiguado o en un ciclo límite.
\end{enumerate}

Los osciladores biológicos reales (el marcapasos del nódulo sinoauricular del corazón, las oscilaciones circadianas de proteínas \emph{Clock/Period}, el ritmo respiratorio) no funcionan como centros frágiles: son notablemente robustos. Si el corazón sufre una perturbación (un sobresalto o una arritmia transitoria), la trayectoria es atraída de regreso al mismo ritmo y a la misma amplitud exacta. Este comportamiento no es un centro, sino un \textbf{ciclo límite atractor}.

\begin{definition}[Ciclo Límite]
Un \textbf{ciclo límite} $\Gamma$ es una trayectoria periódica cerrada e \textbf{aislada} en el espacio de fases de un sistema continuo. El término \emph{aislado} significa que en una vecindad tubular de $\Gamma$ no existen otras órbitas cerradas:
\begin{itemize}
    \item \textbf{Estable (atractor):} Todas las trayectorias vecinas espiralan hacia $\Gamma$ cuando $t \to +\infty$.
    \item \textbf{Inestable (repulsor):} Todas las trayectorias vecinas se alejan de $\Gamma$ cuando $t \to +\infty$.
    \item \textbf{Semi-estable:} Las trayectorias se aproximan a $\Gamma$ por un lado y se alejan por el otro.
\end{itemize}
\end{definition}

\subsection{El oscilador de Van der Pol y osciladores biológicos}

El prototipo fundamental de un oscilador no lineal con ciclo límite fue formulado en 1926 por el ingeniero y físico holandés Balthasar van der Pol para describir circuitos con tubos de vacío (tríodos) y, posteriormente, para modelar el ritmo autónomo del corazón humano:
\begin{equation}
    \ddot{x} - \mu (1 - x^2)\dot{x} + x = 0, \quad \mu > 0.
    \label{eq:vanderpol_ode}
\end{equation}
Escribiendo en variables de estado $(x, y = \dot{x})$:
\begin{equation}
\begin{aligned}
    \dot{x} &= y, \\
    \dot{y} &= -x + \mu(1 - x^2)y.
\end{aligned}
\label{eq:vanderpol_sistema}
\end{equation}

\paragraph{Mecanismo físico de disipación no lineal.}
El término $-\mu(1 - x^2)\dot{x}$ actúa como un coeficiente de amortiguamiento dependiente del estado:
\begin{itemize}
    \item Para pequeñas amplitudes ($|x| < 1$): El coeficiente es negativo ($-\mu(1-x^2) < 0$). El sistema experimenta amortiguamiento negativo, inyectando energía y amplificando cualquier perturbación lejos del origen inestable.
    \item Para grandes amplitudes ($|x| > 1$): El coeficiente es positivo ($+\mu(x^2-1) > 0$). El sistema disipa energía, frenando la oscilación.
\end{itemize}
El balance dinámico exacto entre la inyección de energía para $|x| < 1$ y la disipación para $|x| > 1$ fuerza la existencia de una única órbita periódica cerrada aislada de amplitud finita $x_{\max} \approx 2$: el \textbf{ciclo límite de Van der Pol} (véase la \cref{fig:vanderpol_portrait}).

\begin{figure}[htbp]
\centering
\begin{tikzpicture}[>=Stealth, scale=0.95]
    \begin{axis}[
        xmin=-3.2, xmax=3.2, ymin=-3.5, ymax=3.5,
        axis lines=middle,
        xlabel={$x$}, ylabel={$y = \dot{x}$},
        xlabel style={at={(ticklabel* cs:1)},anchor=west},
        ylabel style={at={(ticklabel* cs:1)},anchor=south},
        width=0.82\linewidth, height=6.8cm,
        grid=both, grid style={gray!20}
    ]
        % Ciclo límite de Van der Pol para mu=1 (aproximación numérica suave)
        \draw[very thick, black!95] plot [smooth cycle, tension=0.8] coordinates {
            (2.0, 0) (1.8, 1.2) (0.8, 2.2) (-0.5, 2.5) (-2.0, 0.2)
            (-1.8, -1.2) (-0.8, -2.2) (0.5, -2.5) (1.9, -0.2)
        };

        % Trayectoria espiral interna que sale del origen y converge al ciclo límite
        \draw[thick, black!65, domain=0:12.56, samples=150, variable=\t]
            plot ({ (0.2 + 0.14*\t)*cos(\t r) }, { (0.2 + 0.14*\t)*sin(\t r) });

        % Trayectoria exterior que entra al ciclo límite
        \draw[thick, black!65, domain=0:8.0, samples=120, variable=\t]
            plot ({ (3.0 - 0.12*\t)*cos(\t r) }, { (3.0 - 0.12*\t)*sin(\t r) });

        % Origen inestable
        \draw[black, fill=white, thick] (axis cs:0,0) circle (2.5pt) node[below right=2pt, font=\tiny] {Foco inestable};

        % Flecha de convergencia
        \node[font=\bfseries\footnotesize, fill=white, inner sep=1pt] at (axis cs:1.4, 2.6) {Ciclo límite $\Gamma$};
        \node[font=\footnotesize, black!75] at (axis cs:-2.0, -2.8) {Atractor asintótico global en $\mathbb{R}^2 \setminus \{\vec{0}\}$};
    \end{axis}
\end{tikzpicture}
\caption{Retrato de fases del oscilador de Van der Pol con $\mu = 1$. El origen es un foco inestable que repele las trayectorias; todas las soluciones (interiores y exteriores) son atraídas asintóticamente hacia el ciclo límite aislado $\Gamma$.}
\label{fig:vanderpol_portrait}
\end{figure}

\paragraph{Universalidad biológica de los ciclos límite.}
Como anticipamos en la introducción histórica del curso (capítulo 2), los ciclos límite son la base teórica de:
\begin{itemize}
    \item \textbf{Marcapasos cardíacos:} El potencial de membrana de las células marcapasos oscila espontáneamente sin necesidad de estímulo externo.
    \item \textbf{Excitabilidad neuronal:} El modelo de FitzHugh--Nagumo (una simplificación bidimensional del sistema de 4 variables de Hodgkin--Huxley) utiliza la geometría tipo Van der Pol para explicar cómo una neurona pasa del reposo a la emisión periódica de espigas (\emph{spiking}).
    \item \textbf{Relojes circadianos y biológicos:} Arthur Winfree~\cite{winfree2001} demostró que la geometría de fases de los ciclos límite permite entender la sincronización y el desfase (\emph{jet lag}) de los ritmos biológicos circadianos bajo pulsos de luz externos.
\end{itemize}

% integrado de body/body.tex

% ============================================================
% FUENTES: Notas Biomate 1; raw_plain_nonlinear_dynamics.txt.
% Strogatz (2015), cap. 8; Murray (2002); Edelstein-Keshet (2007).
% ============================================================
\section{Bifurcación de Hopf}

En los capítulos de la Parte III estudiamos las bifurcaciones en flujos unidimensionales (saddle-node, transcrítica y pitchfork). En todas ellas, la pérdida de estabilidad de un equilibrio ocurría necesariamente a través de un autovalor real que cruzaba por cero ($\lambda = 0$).

Al expandir nuestro análisis a dos dimensiones, emerge un mecanismo de inestabilidad cualitativamente superior: un par de autovalores complejos conjugados puede cruzar el eje imaginario con velocidad no nula. Este fenómeno, descubierto por Aleksandr Andronov en la década de 1930 y popularizado por Eberhard Hopf en 1942, se conoce como la \textbf{bifurcación de Hopf} (o de Andronov--Hopf)~\cite{strogatz2015,murray2002}. La bifurcación de Hopf es la vía canónica por la cual un estado estacionario pierde su estabilidad y engendra espontáneamente una oscilación periódica autosostenida (un ciclo límite).

\subsection{El Teorema de Hopf en el plano}

Consideremos un sistema bidimensional parametrizado por $\mu \in \mathbb{R}$:
\begin{equation}
    \dot{\vec{x}} = \vec{f}(\vec{x}; \mu), \quad \vec{x} \in \mathbb{R}^2.
\end{equation}
Supongamos que el sistema posee un punto fijo de equilibrio $\vec{x}^*(\mu)$ para todos los valores de $\mu$ cercanos a un valor crítico $\mu_c = 0$, de modo que $\vec{f}(\vec{x}^*(\mu); \mu) = \vec{0}$.

\begin{theorem}[Teorema de la bifurcación de Hopf (Andronov--Hopf)]
\label{thm:hopf_teorema}
Sea $J(\mu) = D\vec{f}(\vec{x}^*(\mu); \mu)$ la matriz jacobiana evaluada en el equilibrio. Supongamos que:
\begin{enumerate}
    \item Los autovalores de $J(\mu)$ son un par de complejos conjugados:
    \begin{equation}
        \lambda_{1,2}(\mu) = \alpha(\mu) \pm i\omega(\mu).
    \end{equation}
    \item En el valor crítico $\mu = 0$, los autovalores cruzan el eje imaginario:
    \begin{equation}
        \alpha(0) = 0, \qquad \omega_0 \equiv \omega(0) > 0.
    \end{equation}
    \item \textbf{Condición de transversalidad:} La parte real de los autovalores cruza el eje imaginario con velocidad estrictamente no nula:
    \begin{equation}
        \left. \frac{d\alpha}{d\mu} \right|_{\mu = 0} \neq 0.
    \end{equation}
\end{enumerate}
Entonces, al variar $\mu$ a través de $\mu = 0$, el punto fijo $\vec{x}^*(\mu)$ cambia de estabilidad (de foco estable a inestable o viceversa) y \textbf{nace o muere una familia de órbitas periódicas cerradas (ciclos límite)} en una vecindad del equilibrio, con frecuencia angular inicial $\omega_0$ y periodo $T_0 \approx 2\pi/\omega_0$.
\end{theorem}

Dependiendo del signo de los términos no lineales de orden superior en el punto crítico, la bifurcación de Hopf se clasifica en dos tipos fundamentales: \textbf{supercrítica} y \textbf{subcrítica}.

\subsection{Bifurcación de Hopf supercrítica}

En la bifurcación supercrítica, la transición hacia la oscilación es continua, suave y reversible (\emph{soft onset}).

\paragraph{Forma normal en coordenadas polares.}
La forma normal canónica en coordenadas polares $(r, \theta)$ se escribe:
\begin{align}
    \dot{r} &= \mu r - r^3, \label{eq:hopf_super_r} \\
    \dot{\theta} &= \omega + b r^2, \label{eq:hopf_super_theta}
\end{align}
donde $\mu$ es el parámetro de bifurcación, $\omega > 0$ es la frecuencia básica y $b \in \mathbb{R}$ modula el cambio de frecuencia con la amplitud. Observemos que la ecuación radial $\dot{r} = \mu r - r^3$ tiene formalmente la estructura de una bifurcación pitchfork supercrítica en el radio $r \ge 0$:

\begin{itemize}
    \item \textbf{Régimen $\mu < 0$:} El único equilibrio radial es $r^* = 0$ (el origen). Dado que $\dot{r} = \mu r - r^3 < 0$ para todo $r > 0$, todas las trayectorias espiralan convergiendo exponencialmente al origen. El origen es un \textbf{foco estable atractor}.
    
    \item \textbf{Punto crítico $\mu = 0$:} La ecuación radial se reduce a $\dot{r} = -r^3$. El origen sigue siendo asintóticamente estable, pero la relajación es lenta y algebraica ($r(t) \sim t^{-1/2}$).
    
    \item \textbf{Régimen $\mu > 0$:} En el origen, $\left.\frac{d\dot{r}}{dr}\right|_{r=0} = \mu > 0$, por lo que el equilibrio $r^* = 0$ se vuelve un \textbf{foco inestable repulsor}. Simultáneamente, aparece un nuevo equilibrio radial no trivial dado por $\mu r - r^3 = 0 \implies r^*(\mu) = \sqrt{\mu}$.
\end{itemize}

Como $\dot{\theta} \approx \omega > 0$, la solución en $r = \sqrt{\mu}$ corresponde a una trayectoria periódica circular cerrada de radio constante: un \textbf{ciclo límite estable} (véase la \cref{fig:hopf_supercritica}).

\begin{figure}[htbp]
\centering
\begin{tikzpicture}[>=Stealth, scale=0.9]
    % Panel izquierdo: Espacio de fases para mu < 0
    \begin{scope}[shift={(-4.5, 0)}]
        \draw[->, thick, black!60] (-2,0) -- (2,0) node[right, font=\tiny] {$x$};
        \draw[->, thick, black!60] (0,-2) -- (0,2) node[above, font=\tiny] {$y$};
        \filldraw[black] (0,0) circle (2.2pt);
        % Espiral hacia adentro
        \draw[thick, black!85, domain=0:9.42, samples=120, variable=\t]
            plot ({ 1.6*exp(-0.35*\t)*cos(\t r) }, { 1.6*exp(-0.35*\t)*sin(\t r) });
        \node[font=\footnotesize, align=center] at (0,-2.4) {\textbf{(a)} $\mu < 0$\\Foco estable};
    \end{scope}

    % Panel central: Espacio de fases para mu > 0
    \begin{scope}[shift={(0, 0)}]
        \draw[->, thick, black!60] (-2,0) -- (2,0) node[right, font=\tiny] {$x$};
        \draw[->, thick, black!60] (0,-2) -- (0,2) node[above, font=\tiny] {$y$};
        \draw[black, fill=white, thick] (0,0) circle (2.2pt);
        % Ciclo límite estable
        \draw[very thick, black!95] (0,0) circle (1.25cm);
        % Espiral interna
        \draw[thick, black!65, domain=0:8.0, samples=100, variable=\t]
            plot ({ (0.2 + 0.12*\t)*cos(\t r) }, { (0.2 + 0.12*\t)*sin(\t r) });
        % Espiral externa
        \draw[thick, black!65, domain=0:6.28, samples=100, variable=\t]
            plot ({ (1.9 - 0.10*\t)*cos(\t r) }, { (1.9 - 0.10*\t)*sin(\t r) });
        \node[font=\footnotesize, align=center] at (0,-2.4) {\textbf{(b)} $\mu > 0$\\Ciclo límite ($r=\sqrt{\mu}$)};
    \end{scope}

    % Panel derecho: Diagrama de bifurcación
    \begin{scope}[shift={(4.5, 0)}]
        \draw[->, thick, black!60] (-2,0) -- (2.2,0) node[right, font=\tiny] {$\mu$};
        \draw[->, thick, black!60] (0,0) -- (0,2.2) node[above, font=\tiny] {$r$};
        % Origen: estable para mu < 0, inestable para mu > 0
        \draw[very thick, black!90] (-1.8,0) -- (0,0);
        \draw[very thick, black!90, dashed] (0,0) -- (1.8,0);
        % Rama de ciclo límite r = sqrt(mu)
        \draw[very thick, black!95, domain=0:1.8, samples=60] plot (\x, {1.2*sqrt(\x)});
        \filldraw[black] (0,0) circle (2.2pt);
        \node[font=\footnotesize, align=center] at (1.1, 1.7) {Ciclo límite\\$r^* = \sqrt{\mu}$};
        \node[font=\footnotesize, align=center] at (0,-2.4) {\textbf{(c)} Diagrama de\\bifurcación};
    \end{scope}
\end{tikzpicture}
\caption{Bifurcación de Hopf supercrítica. (a) Para $\mu < 0$, el origen es un foco estable. (b) Para $\mu > 0$, el origen se vuelve inestable y nace un ciclo límite atractor. (c) Diagrama de bifurcación mostrando el crecimiento continuo de la amplitud como $r \sim \sqrt{\mu}$.}
\label{fig:hopf_supercritica}
\end{figure}

\subsection{Bifurcación de Hopf subcrítica e histéresis}

En la bifurcación de Hopf subcrítica, el término cúbico no lineal tiene signo opuesto ($+r^3$) y actúa como desestabilizador. Para modelar un sistema físicamente confinado se incluye un término estabilizador de orden superior (quíntico, $-r^5$):
\begin{align}
    \dot{r} &= \mu r + r^3 - r^5, \label{eq:hopf_sub_r} \\
    \dot{\theta} &= \omega + b r^2. \label{eq:hopf_sub_theta}
\end{align}

\paragraph{Dinámica y catástrofe dinámica (\emph{Hard onset}).}
El comportamiento cualitativo revela una rica estructura de biestabilidad:
\begin{itemize}
    \item Para $\mu < 0$: El origen $r = 0$ es un foco estable, pero está rodeado por un \textbf{ciclo límite inestable} de radio $r_1^* \approx \sqrt{-\mu}$. Este ciclo inestable define la frontera de la cuenca de atracción del origen: si una perturbación saca al sistema fuera de este ciclo, la trayectoria es expulsada hacia un \textbf{ciclo límite estable de gran amplitud} ($r_2^* \approx 1$).
    \item En $\mu = 0$: El ciclo límite inestable colisiona con el origen y lo destruye.
    \item Para $\mu > 0$: El origen es inestable y no existen oscilaciones pequeñas; el sistema experimenta un salto brusco e instantáneo (\emph{hard onset}) hacia la oscilación de gran amplitud.
\end{itemize}

Si una vez encendida la oscilación se reduce nuevamente $\mu$ por debajo de cero ($\mu < 0$), el sistema no se apaga en $\mu = 0$, sino que permanece en la rama de oscilación grande hasta alcanzar el punto de silla-nodo de ciclos límite en $\mu_c < 0$. Este fenómeno es la \textbf{histéresis oscilatoria}.

\begin{figure}[htbp]
\centering
\begin{tikzpicture}[>=Stealth, scale=0.95]
    \draw[->, thick, black!60] (-2.5,0) -- (2.5,0) node[right, font=\small] {$\mu$};
    \draw[->, thick, black!60] (0,0) -- (0,2.6) node[above, font=\small] {$r$};

    % Eje r=0: Estable para mu < 0, inestable para mu > 0
    \draw[very thick, black!90] (-2.2,0) -- (0,0);
    \draw[very thick, black!90, dashed] (0,0) -- (2.2,0);
    \filldraw[black] (0,0) circle (2.2pt);

    % Ciclo límite inestable (línea discontinua) que nace subcríticamente hacia la izquierda
    \draw[thick, black!85, dashed, domain=-1.2:0, samples=60] plot (\x, {1.1*sqrt(-\x)});

    % Ciclo límite estable de gran amplitud (línea sólida superior)
    \draw[very thick, black!95, domain=-1.2:2.2, samples=80] plot (\x, {1.7 + 0.15*\x});

    % Pliegue en mu_c = -1.2
    \draw[thick, black!85, dashed] (-1.2, 1.2) -- (-1.2, 1.52);

    % Flechas de salto por histéresis
    \draw[->, thick, black!80] (0.05, 0.1) -- (0.05, 1.6);
    \node[right, font=\tiny] at (0.05, 0.85) {Salto brusco (\emph{hard onset})};

    \draw[->, thick, black!80] (-1.25, 1.4) -- (-1.25, 0.1);
    \node[left, font=\tiny] at (-1.25, 0.75) {Colapso};

    \node[font=\footnotesize, align=center] at (-1.5, 0.8) {Ciclo inestable\\(frontera cuenca)};
    \node[font=\footnotesize, align=center] at (1.2, 2.1) {Ciclo estable\\de gran amplitud};
    \node[font=\footnotesize, align=center] at (-1.4, -0.4) {Biestabilidad};
\end{tikzpicture}
\caption{Diagrama de bifurcación de Hopf subcrítica. La rama inestable (línea discontinua) retrocede hacia $\mu < 0$ y se pliega conectando con un ciclo límite estable de gran amplitud, generando histéresis y biestabilidad entre el reposo y la oscilación.}
\label{fig:hopf_subcritica}
\end{figure}

\subsection{Aplicaciones biológicas: excitabilidad neuronal y relojes bioquímicos}

\paragraph{Excitabilidad neuronal Tipo II (Hodgkin--Huxley).}
En el célebre modelo electrofisiológico del axón gigante de calamar formulado por Alan Hodgkin y Andrew Huxley (1952), el potencial de membrana $V(t)$ y las variables de compuerta iónica de sodio y potasio $(m, n, h)$ responden a una corriente externa inyectada $I_{\text{ext}}$.

Para corrientes bajas ($I_{\text{ext}} < I_c$), la neurona permanece en un estado de reposo hiperpolarizado estable. Al superar la corriente crítica ($I_{\text{ext}} > I_c$), el equilibrio de reposo pierde estabilidad a través de una \textbf{bifurcación de Hopf subcrítica}. El potencial salta bruscamente a un tren de potenciales de acción periódicos (\emph{spiking}) con una frecuencia de disparo no nula y finita desde el encendido. Esta respuesta se conoce en neurociencia como \textbf{excitabilidad neuronal de Tipo II}~\cite{edelstein2007}.

\paragraph{Oscilaciones glicolíticas y relojes genéticos.}
En la bioquímica celular, las oscilaciones periódicas en la concentración de intermediarios metabólicos como el NADH en levaduras y las oscilaciones sintéticas en circuitos genéticos (como el \emph{Repressilator}) emergen cuando la tasa de sustrato o la inhibición por retroalimentación negativa supera el umbral crítico establecido por el teorema de Hopf\aside{Para detectar numéricamente bifurcaciones de Hopf en sistemas no lineales de alta dimensión, se emplean algoritmos de continuación numérica como AUTO o XPPAUT, que rastrean los ceros de $\operatorname{Re}(\lambda)$ a lo largo de ramas de equilibrios.}.


% ---- Parte V — Estudio de caso (cierre del curso) ----
\part{Estudio de caso: del modelo completo a la reducción}
% ============================================================
% FUENTE: Dor & Alon (2026), PLoS Comput Biol 22(3):e1014111,
%         papers_local/journal_pcbi_1014111.txt. Ecuaciones (1)-(8)
%         y Tabla 10 (parametros). Cierre del curso: reduccion
%         pedagogica 4D -> 1D -> 2D reutilizando todo el aparato.
% ============================================================
\section{Estudio de caso: homeostasis de glóbulos rojos y anemia}

Cerramos el curso con un modelo reciente y honesto de fisiología: el modelo mínimo de eritropoyesis de Dor y Alon~\cite{dor2026}. No es un ejemplo de juguete; son cuatro ecuaciones diferenciales acopladas cuyos parámetros provienen, casi todos, de mediciones experimentales, y que reproducen la relación EPO--hemoglobina observada en $\sim 1830$ pacientes de $36$ estudios. Nuestro objetivo no es reproducir su ajuste estadístico, sino \emph{desmontarlo}: partir del sistema 4D completo y, mediante hipótesis de separación de escalas temporales enunciadas explícitamente, destilarlo primero a una única ecuación de línea de fase (Parte II) y luego a un plano de fases (Parte IV). Cada paso reutiliza una herramienta concreta de un capítulo previo; la \cref{tab:mapa_herramientas} al final es el mapa de esa reutilización.

La biología es la de un termostato. Los glóbulos rojos (RBC) transportan oxígeno; si su número cae, los tejidos renales detectan hipoxia y secretan la hormona \textbf{eritropoyetina} (EPO), que rescata a los progenitores eritroides de la apoptosis y acelera su maduración hacia nuevos glóbulos rojos. Más glóbulos rojos $\Rightarrow$ más oxígeno $\Rightarrow$ menos EPO. Es un lazo de \textbf{retroalimentación negativa}, exactamente como el PLL del \cref{fig:pll_block_diagram} del \cref{sec:historia} pero implementado en médula ósea y riñón.

\subsection{El modelo completo (4D)}
\label{sec:eritro_4d}

El sistema sigue cuatro poblaciones, cada una gobernada por la plantilla universal $\dot X = (\text{producción}) - (\text{remoción})$:
\begin{align}
    \dot{H} &= \underbrace{a(E)\,H\left(1 - \frac{H}{H_{\max}}\right)}_{\text{proliferación logística}} - \underbrace{d(E)\,H}_{\text{diferenciación}}, \label{eq:eritro_H}\\[4pt]
    \dot{R} &= \underbrace{d(E)\,H}_{\text{entrada desde }H} - \underbrace{\gamma_R(C)\,R}_{\text{maduración a RBC}}, \label{eq:eritro_R}\\[4pt]
    \dot{C} &= \underbrace{\gamma_R(C)\,R}_{\text{liberación de RBC}} - \underbrace{\gamma_C\,C}_{\text{muerte por senescencia}}, \label{eq:eritro_C}\\[4pt]
    \dot{E} &= \underbrace{\sigma_E\,e^{-C/D}}_{\text{síntesis renal (hipoxia)}} - \underbrace{\gamma_E\,H\,E}_{\text{aclaramiento por endocitosis}}. \label{eq:eritro_E}
\end{align}

Las variables son: $H(t)$ = progenitores eritroides CFU-E en médula (células), $R(t)$ = reticulocitos medulares (células), $C(t)$ = glóbulos rojos circulantes (células, proporcional a la hemoglobina), $E(t)$ = EPO en plasma (U/L). Las funciones reguladoras son de tipo Michaelis--Menten (proliferación y diferenciación crecen con la EPO y saturan) y una tasa de maduración dependiente de la carga:
\begin{equation}
    a(E) = \frac{a_{\max}\,E}{K_a + E}, \qquad
    d(E) = \frac{d_{\max}\,E}{K_d + E}, \qquad
    \gamma_R(C) = \frac{1}{m_R\,C + n_R}.
    \label{eq:eritro_MM}
\end{equation}

\paragraph{Interpretación mecanística término a término.} La EPO no crea progenitores de la nada: modula la \emph{fracción} de progenitores que prolifera ($a(E)$) frente a la que se diferencia y abandona el compartimento ($d(E)$). El factor logístico $(1-H/H_{\max})$ impone una capacidad de carga medular $H_{\max}$: la médula no es infinita. La ecuación \eqref{eq:eritro_E} contiene la sutileza clave del paper: la EPO no se degrada a tasa constante, sino que se \emph{consume} al unirse a sus receptores sobre los progenitores $H$; por eso su aclaramiento es $\gamma_E H E$ y no $\gamma_E E$. Esto tendrá una consecuencia clínica sorprendente (Ejercicio~5). La síntesis $\sigma_E e^{-C/D}$ codifica el sensor de hipoxia: decae exponencialmente cuando hay muchos glóbulos rojos y se dispara cuando $C\to 0$.

\begin{table}[htbp]
\centering
\footnotesize
\caption{Símbolos, significado biológico, valor y unidad (Dor--Alon~\cite{dor2026}, Tabla 10). Todos salvo $d_{\max}$ y $K_d$ provienen de medición directa.}
\label{tab:eritro_simbolos}
\begin{tabular}{llll}
\hline
Símbolo & Significado & Valor & Unidad \\
\hline
$H_{\max}$   & Capacidad medular de progenitores       & $10^{13}$              & células \\
$a_{\max}$   & Tasa máx. de proliferación CFU-E        & $5.28\times10^{-6}$   & s$^{-1}$ \\
$d_{\max}$   & Tasa máx. de diferenciación             & $1.5\times10^{-6}$    & s$^{-1}$ \\
$K_a$        & EPO a media proliferación               & $30$                  & mU/mL \\
$K_d$        & EPO a media diferenciación              & $23$                  & mU/mL \\
$\gamma_C$   & Tasa de muerte de RBC                   & $1.0\times10^{-7}$    & s$^{-1}$ \\
$\sigma_E$   & Producción máx. de EPO                  & $0.035$               & U\,s$^{-1}$L$^{-1}$ \\
$D$          & Escala de RBC del sensor de hipoxia     & $5.3\times10^{12}$    & células \\
$\gamma_E$   & Aclaramiento de EPO por receptor        & $6.1\times10^{-18}$   & s$^{-1}$célula$^{-1}$ \\
$\tau_R$     & Retardo de maduración CFU-E$\to$reticulocito & $2.16\times10^{5}$ & s \\
$m_R,\,n_R$  & Pendiente e intercepto del tránsito medular & $10^{-8};\ 5.8\times10^{4}$ & (cél$\cdot$s)$^{-1}$; s \\
\hline
\end{tabular}
\end{table}

De $\gamma_C = 10^{-7}\,\mathrm{s}^{-1}$ leemos la escala de tiempo característica del \cref{sec:estabilidad_lineal}: la vida media del glóbulo rojo es $\tau_C = 1/\gamma_C = 10^{7}\,\mathrm{s}\approx 116$ días, cercana al valor canónico de $120$ días. Este número es el protagonista de la reducción.

\begin{figure}[htbp]
\centering
\begin{tikzpicture}[>=Stealth, auto, >=Latex, font=\small]
    \tikzstyle{block} = [draw, rectangle, rounded corners=2pt, minimum height=1.0cm, minimum width=1.9cm, fill=black!5, align=center, font=\small]
    \node[block] (H) {Progenitores\\$H$};
    \node[block, right=1.9cm of H] (R) {Reticulocitos\\$R$};
    \node[block, right=1.9cm of R] (C) {RBC circul.\\$C$};
    \node[block, above=1.5cm of R, fill=black!10] (E) {Riñón / EPO\\$E$};
    \draw[->, thick] (H) -- node[above,font=\tiny]{$d(E)H$} (R);
    \draw[->, thick] (R) -- node[above,font=\tiny]{$\gamma_R(C)R$} (C);
    \draw[->, thick] (E) -| node[near start,above,font=\tiny]{$a(E),d(E)$} (H.north);
    % feedback: C reduce hipoxia -> reduce E
    \draw[->, thick, dashed] (C.north) |- node[near end,above,font=\tiny]{$-\,e^{-C/D}$ (O$_2$)} (E.east);
    % consumo de EPO por H
    \draw[->, thick, dotted] (H.north) -- node[left,font=\tiny]{$-\gamma_E H E$} (E.west);
    % salida y muerte
    \draw[->, thick] (C.east) -- ++(1.1,0) node[right,font=\tiny]{$\gamma_C C$ (muerte)};
\end{tikzpicture}
\caption{Cascada eritropoyética como lazo de retroalimentación negativa (compárese con el diagrama de bloques del PLL, \cref{fig:pll_block_diagram}). Flechas continuas: flujos de maduración. Discontinua: sensor de hipoxia (más $C$ $\Rightarrow$ menos EPO). Punteada: la EPO se consume al ser captada por los progenitores $H$.}
\label{fig:eritro_cascada}
\end{figure}

\subsection{Reducción a 1D: separación de escalas y línea de fase}
\label{sec:eritro_1d}

\paragraph{Hipótesis de reducción (explícitas).} Comparemos las escalas de tiempo de las cuatro variables:
\begin{itemize}
    \item EPO: vida media de \emph{unas horas} (aclaramiento rápido).
    \item Reticulocitos y progenitores: tránsito medular $m_R C + n_R \approx 1.1\times10^{5}\,\mathrm{s}\approx 1.3$ días; retardo $\tau_R\approx 2.5$ días.
    \item Glóbulos rojos: vida media $\tau_C\approx 116$ días.
\end{itemize}
Hay una jerarquía nítida: EPO $\ll$ (reticulocitos, progenitores) $\ll$ RBC, con más de un orden de magnitud de separación en cada salto. Por el principio de \textbf{eliminación adiabática} (una variable mucho más rápida se ``esclaviza'' instantáneamente al valor de equilibrio dictado por las lentas), fijamos $\dot H=\dot R=\dot E=0$ y dejamos a $C$ como \emph{la única} variable dinámica lenta.

\paragraph{Un regalo de la saturación: $H^*$ casi constante.} De $\dot H=0$ con $H\neq 0$ se obtiene, tras dividir \eqref{eq:eritro_H} entre $H$,
\begin{equation}
    a(E)\left(1 - \frac{H}{H_{\max}}\right) = d(E)
    \quad\Longrightarrow\quad
    H^* = H_{\max}\left(1 - \frac{d(E)}{a(E)}\right).
    \label{eq:eritro_Hstar}
\end{equation}
A niveles fisiológicos de EPO, tanto $a(E)$ como $d(E)$ están saturadas ($E\gg K_a,K_d$), de modo que el cociente $d(E)/a(E)\to d_{\max}/a_{\max}$ deja de depender de $E$. Con los valores de la \cref{tab:eritro_simbolos}, $\rho \equiv d_{\max}/a_{\max}=1.5/5.28\approx 0.284$, y por tanto
\begin{equation}
    H^*\approx H_{\max}(1-\rho)\approx 7.2\times10^{12}\ \text{células},
\end{equation}
una constante. Esto es lo que permite cerrar la reducción de forma limpia.

\paragraph{La ecuación 1D.} Con $\dot R=0$, el flujo de reticulocitos que sale iguala al que entra: $\gamma_R(C)R = d(E)H^*$. Sustituyendo en \eqref{eq:eritro_C},
\begin{equation}
    \dot{C} = d(E)\,H^* - \gamma_C\,C.
\end{equation}
Falta esclavizar la EPO. De $\dot E=0$ en \eqref{eq:eritro_E} con $H\approx H^*$:
\begin{equation}
    E(C) = \frac{\sigma_E}{\gamma_E H^*}\,e^{-C/D} \equiv \hat E\,e^{-C/D},
    \qquad \hat E = \frac{\sigma_E}{\gamma_E H^*}\approx 8\times10^{2}\ \mathrm{U/L}.
    \label{eq:eritro_EC}
\end{equation}
$E(C)$ es \emph{monótona decreciente}: más glóbulos rojos, menos EPO. Insertándola queda una única ecuación de línea de fase,
\begin{equation}
    \boxed{\ \dot{C} = P(C) - \gamma_C\,C, \qquad P(C) = H^*\,d\bigl(E(C)\bigr) = \frac{H^* d_{\max}\,\hat E\,e^{-C/D}}{K_d + \hat E\,e^{-C/D}}.\ }
    \label{eq:eritro_1d}
\end{equation}
La producción $P(C)$ es positiva, acotada y \emph{estrictamente decreciente} en $C$ (composición de $d(\cdot)$ creciente con $E(C)$ decreciente): es la firma de la retroalimentación negativa.

\paragraph{Equilibrios y estabilidad.} Los equilibrios resuelven $P(C^*) = \gamma_C C^*$: intersección de una curva decreciente y acotada con una recta creciente que pasa por el origen. Existe \textbf{exactamente uno}, $C^*>0$. Aplicando el criterio del \cref{sec:estabilidad_lineal} a $g(C)=P(C)-\gamma_C C$,
\begin{equation}
    g'(C^*) = \underbrace{P'(C^*)}_{<0} - \gamma_C < 0,
\end{equation}
así que \emph{siempre} es estable, con tiempo de relajación $\tau = 1/|g'(C^*)|$. El sistema es \textbf{monoestable}: hay un único estado sano, globalmente atractor. Esto no es una simplificación excesiva nuestra, sino exactamente lo que reporta el paper: ``the model's steady state is highly robust to parameter changes\ldots\ and does not exhibit oscillatory behavior''~\cite{dor2026}.

\paragraph{¿Dónde está entonces la anemia?} No es un segundo atractor, sino el \emph{mismo} equilibrio desplazado cuando un parámetro cambia. La \cref{fig:eritro_1d} lo muestra: reducir la capacidad productiva (menor $H^*$, anemia aplásica) o aumentar la destrucción ($\gamma_C\to 4.8\gamma_C$, anemia hemolítica) baja el punto de corte $C^*$. Este es el mecanismo unificador del paper (su Tabla~9): distintas anemias son distintos \emph{deslizamientos suaves} del único equilibrio, no catástrofes.

\begin{figure}[htbp]
\centering
\begin{tikzpicture}[>=Stealth, scale=1.0]
    % ejes
    \draw[->, thick, black!60] (-0.3,0) -- (6.3,0) node[right,font=\small] {$C$};
    \draw[->, thick, black!60] (0,-1.6) -- (0,1.8) node[above,font=\small] {$\dot{C}=g(C)$};
    % curva sana g(C): empieza positiva, cruza en C*=3.5, luego negativa
    \draw[very thick, black!85, domain=0.05:6.0, samples=80]
        plot (\x, {1.5*exp(-0.55*\x) - 0.30*\x});
    % curva anemia (desplazada, cruza antes)
    \draw[very thick, black!45, dashed, domain=0.05:6.0, samples=80]
        plot (\x, {0.85*exp(-0.55*\x) - 0.30*\x});
    % equilibrios aprox: sano ~ 3.05, anemia ~ 1.6
    \filldraw[black] (3.05,0) circle (2pt) node[below=3pt,font=\tiny] {$C^*_{\text{sano}}$};
    \filldraw[black!55] (1.55,0) circle (2pt) node[above=3pt,font=\tiny] {$C^*_{\text{anemia}}$};
    % flechas de flujo sobre el eje
    \draw[->, thick, black!80] (0.6,0.28) -- (1.4,0.28);
    \draw[->, thick, black!80] (5.2,-0.28) -- (4.4,-0.28);
    \node[font=\tiny, align=center] at (5.0,1.3) {sano};
    \node[font=\tiny, align=center, black!55] at (2.6,-1.1) {menor $H^*$\\o mayor $\gamma_C$};
\end{tikzpicture}
\caption{Línea de fase de la ecuación reducida \eqref{eq:eritro_1d}. La curva de producción neta $g(C)=P(C)-\gamma_C C$ (negra) cruza el eje una sola vez, en el equilibrio sano estable $C^*_{\text{sano}}$ ($g'<0$). Una pérdida de capacidad medular o un aumento de destrucción (curva gris) desplaza el único equilibrio hacia la anemia sin crear nuevos puntos fijos.}
\label{fig:eritro_1d}
\end{figure}

\paragraph{Una bifurcación real: el establecimiento del progenitor.} Hay un lugar donde sí aparece una bifurcación genuina, y es la ecuación del progenitor \eqref{eq:eritro_H} tomada aisladamente. Escribiéndola con $a,d$ congeladas en sus valores saturados,
\begin{equation}
    \dot{H} = (a-d)\,H - \frac{a}{H_{\max}}\,H^2,
    \label{eq:eritro_transcritica}
\end{equation}
que es \emph{exactamente} la forma normal de la bifurcación transcrítica del \cref{sec:transcritica} con parámetro $r=a-d$. Los equilibrios son $H=0$ (médula vacía) y $H^*=H_{\max}(1-d/a)$. Linealizando: en $H=0$, $\,g'(0)=a-d$; en $H^*$, $\,g'(H^*)=-(a-d)$. Cuando el cociente $\rho=d/a$ cruza $1$ (la diferenciación supera a la proliferación) las dos ramas \emph{intercambian estabilidad}: por debajo de $\rho=1$ la médula sostiene un pool sano; por encima, colapsa a $H^*=0$, la aplasia total. La \cref{fig:eritro_bif} resume ambos fenómenos.

\begin{figure}[htbp]
\centering
\begin{tikzpicture}[>=Stealth, scale=0.95]
  % Panel (a): transcritica de H
  \begin{scope}
    \draw[->, thick, black!60] (-0.2,0) -- (4.3,0) node[right,font=\tiny] {$\rho=d/a$};
    \draw[->, thick, black!60] (0,-1.3) -- (0,1.9) node[above,font=\tiny] {$H^*$};
    \draw[dashed, black!40] (2.0,-1.3) -- (2.0,1.9);
    \node[font=\tiny] at (2.0,-1.55) {$\rho=1$};
    % rama H*=Hmax(1-rho): estable para rho<1
    \draw[very thick, black!85, domain=0.1:2.0, samples=40] plot (\x, {1.6*(1-\x/2.0)});
    % rama H=0 inestable para rho<1 (dashed), estable rho>1 (solid)
    \draw[very thick, black!85, dashed] (0.1,0) -- (2.0,0);
    \draw[very thick, black!85] (2.0,0) -- (4.1,0);
    % continuacion negativa (no fisica) discontinua
    \draw[very thick, black!45, dashed, domain=2.0:4.1, samples=30] plot (\x, {1.6*(1-\x/2.0)});
    \node[font=\tiny, align=center] at (2.2,1.5) {sano};
    \node[font=\tiny, align=center] at (3.2,0.25) {aplasia};
    \node[font=\footnotesize] at (2.0,-2.0) {\textbf{(a)} progenitor: transcrítica};
  \end{scope}
  % Panel (b): C* suave vs 1/gammaC
  \begin{scope}[shift={(6.5,0)}]
    \draw[->, thick, black!60] (-0.2,0) -- (4.3,0) node[right,font=\tiny] {$1/\gamma_C$ (vida media)};
    \draw[->, thick, black!60] (0,-1.3) -- (0,1.9) node[above,font=\tiny] {$C^*$};
    \draw[very thick, black!85, domain=0.1:4.1, samples=60] plot (\x, {1.7*(1-exp(-0.9*\x))});
    \filldraw[black] (0.7,{1.7*(1-exp(-0.63))}) circle (1.6pt) node[below right,font=\tiny] {hemólisis};
    \filldraw[black] (3.0,{1.7*(1-exp(-2.7))}) circle (1.6pt) node[below right,font=\tiny] {sano};
    \node[font=\footnotesize] at (2.0,-2.0) {\textbf{(b)} RBC: deslizamiento suave};
  \end{scope}
\end{tikzpicture}
\caption{(a) La subred del progenitor sufre una bifurcación \emph{transcrítica} genuina (\cref{sec:transcritica}) al cruzar $\rho=d_{\max}/a_{\max}=1$: por debajo hay médula sana ($H^*>0$ estable), por encima colapsa a aplasia ($H^*=0$). (b) En cambio, el pool de glóbulos rojos $C^*$ depende \emph{de forma suave y monótona} del parámetro de destrucción $\gamma_C$: no hay bifurcación, sólo desplazamiento del único equilibrio. La honestidad exige distinguir ambos casos.}
\label{fig:eritro_bif}
\end{figure}

\subsection{Reducción a 2D: reintroducir la dinámica de la EPO}
\label{sec:eritro_2d}

La reducción 1D descansaba en esclavizar la EPO ($\dot E=0$). Relajemos \emph{esa} hipótesis --- es la de mayor valor pedagógico, porque la EPO es la variable reguladora y su acoplamiento con $C$ genera un plano de fases con la estructura activador--inhibidor que estudiamos en la Parte IV. Mantenemos $H\approx H^*$ y $\dot R=0$, pero conservamos $C$ y $E$ como dinámicas:
\begin{equation}
    \dot{C} = H^* d(E) - \gamma_C C, \qquad
    \dot{E} = \sigma_E e^{-C/D} - \gamma_E H^* E.
    \label{eq:eritro_2d}
\end{equation}

\paragraph{Nullclines (\cref{sec:nullclines}).} Las isoclinas nulas son:
\begin{align}
    \dot C = 0:\quad & C = \frac{H^* d(E)}{\gamma_C} = \frac{H^* d_{\max}}{\gamma_C}\cdot\frac{E}{K_d+E} && \text{(creciente y saturante en } E),\\
    \dot E = 0:\quad & E = \frac{\sigma_E}{\gamma_E H^*}\,e^{-C/D} = \hat E\,e^{-C/D} && \text{(decreciente en } C).
\end{align}
La primera dice ``más EPO $\Rightarrow$ más glóbulos rojos en equilibrio''; la segunda, ``más glóbulos rojos $\Rightarrow$ menos EPO''. Una curva creciente y una decreciente se cortan \textbf{una sola vez}: un único equilibrio $(C^*,E^*)$, coherente con el análisis 1D.

\paragraph{Jacobiano y clasificación (\cref{sec:lineales2d}).} Evaluando en el equilibrio, y usando en $\partial\dot E/\partial C$ la relación de equilibrio $\sigma_E e^{-C^*/D}=\gamma_E H^* E^*$:
\begin{equation}
    J = \begin{pmatrix}
        -\gamma_C & H^* d'(E^*) \\[3pt]
        -\dfrac{\gamma_E H^* E^*}{D} & -\gamma_E H^*
    \end{pmatrix},
    \qquad d'(E^*) = \frac{d_{\max}K_d}{(K_d+E^*)^2}>0.
\end{equation}
Leemos traza y determinante:
\begin{align}
    \operatorname{tr}J &= -\gamma_C - \gamma_E H^* < 0 \quad\text{(siempre)},\\
    \det J &= \gamma_C\gamma_E H^* + \underbrace{H^* d'(E^*)}_{>0}\cdot\underbrace{\frac{\gamma_E H^* E^*}{D}}_{>0} > 0 \quad\text{(siempre)}.
\end{align}
Con $\operatorname{tr}J<0$ y $\det J>0$, el equilibrio es \textbf{incondicionalmente estable}: nodo si $(\operatorname{tr}J)^2>4\det J$ (recuperación monótona), foco espiral si $(\operatorname{tr}J)^2<4\det J$ (recuperación con sobreimpulso amortiguado). Ambos regímenes se ilustran en la \cref{fig:eritro_2d}.

\paragraph{Honestidad sobre las oscilaciones.} Un lector que venga del \cref{sec:hopf} preguntará: ¿puede este sistema sufrir una bifurcación de Hopf y oscilar? \textbf{No}, y la razón es estructural: la traza es una suma de términos negativos que \emph{nunca} puede cruzar cero. Sin cruce de traza no hay Hopf. Esto concuerda con lo que afirma el paper: su modelo (sin retardo) es robustamente monoestable~\cite{dor2026}. Las oscilaciones periódicas de la hematopoyesis --- leucemia mieloide crónica (CML) periódica, neutropenia cíclica, anemia hemolítica autoinmune cíclica --- \emph{sí existen} clínicamente, pero \emph{no} son una predicción de este modelo minimal. Son el comportamiento clásico de modelos de eritropoyesis con \textbf{retardo}: la maduración $\tau_R$ del \cref{tab:eritro_simbolos} convierte \eqref{eq:eritro_2d} en una ecuación con retardo,
\begin{equation}
    \dot{C}(t) = H^* d\bigl(E(t-\tau_R)\bigr) - \gamma_C C(t),
    \label{eq:eritro_retardo}
\end{equation}
y un lazo de retroalimentación negativa con retardo se desestabiliza vía Hopf cuando $\tau_R$ supera un umbral crítico $\tau_c$ (mayor retardo $\Rightarrow$ el sistema ``corrige tarde'' y se pasa de rosca). Este es el resultado clásico de Mackey y Glass~\cite{mackeyglass1977} y de Bélair, Mackey y Mahaffy~\cite{belair1995} --- este último es la referencia [16] que el propio Dor--Alon cita. En resumen: la monoestabilidad es un resultado del modelo; las oscilaciones son una \emph{analogía} rigurosa con modelos de retardo bien establecidos, no una afirmación del paper. Distinguir ambas cosas es el hábito honesto que este curso quiere dejar.

\begin{figure}[htbp]
\centering
\begin{tikzpicture}[>=Stealth, scale=1.0]
    \draw[->, thick, black!60] (-0.2,0) -- (5.2,0) node[right,font=\small] {$C$};
    \draw[->, thick, black!60] (0,-0.2) -- (0,4.2) node[above,font=\small] {$E$};
    % E-nullcline decreciente
    \draw[very thick, black!85, domain=0.15:5.0, samples=80] plot (\x, {3.6*exp(-0.5*\x)});
    \node[black!85, font=\tiny] at (4.4,0.55) {$\dot E=0$};
    % C-nullcline creciente saturante (C funcion de E) -> dibujamos como E funcion de C invertida
    \draw[very thick, black!55, domain=0.15:5.0, samples=80] plot (\x, {0.75*\x/(3.0-0.5*\x + 0.001)});
    \node[black!55, font=\tiny] at (3.6,3.7) {$\dot C=0$};
    % equilibrio interseccion aprox (1.7, 1.55)
    \filldraw[black] (1.7,1.55) circle (2pt) node[above right,font=\tiny] {$(C^*,E^*)$};
    % trayectoria espiral hacia el equilibrio
    \draw[->, thick, black!70, domain=0:5.5, samples=90, variable=\t]
        plot ({1.7 + 1.4*exp(-0.28*\t)*cos(150*\t)}, {1.55 + 1.4*exp(-0.28*\t)*sin(150*\t)});
    \node[font=\tiny, align=center] at (3.9,2.6) {foco estable\\(sobreimpulso)};
\end{tikzpicture}
\caption{Plano de fases de la reducción 2D \eqref{eq:eritro_2d}. La nullcline $\dot E=0$ (negra, decreciente) y la $\dot C=0$ (gris, creciente y saturante) se cortan en un único equilibrio. Como $\operatorname{tr}J<0$ y $\det J>0$ siempre, toda trayectoria converge; cuando $(\operatorname{tr}J)^2<4\det J$ lo hace en espiral (recuperación con sobreimpulso amortiguado tras una hemorragia). No hay ciclo límite: la traza no puede cambiar de signo.}
\label{fig:eritro_2d}
\end{figure}

\subsection{Conexión hacia atrás: mapa de herramientas}
\label{sec:eritro_mapa}

Todo lo anterior fue reciclaje. La \cref{tab:mapa_herramientas} explicita qué herramienta de qué capítulo hizo cada paso: es el cierre del curso.

\begin{table}[htbp]
\centering
\footnotesize
\caption{Mapa de reutilización: cada paso de la reducción invoca una herramienta previa.}
\label{tab:mapa_herramientas}
\begin{tabular}{lll}
\hline
Herramienta & Capítulo & Uso en el estudio de caso \\
\hline
Plantilla producción$-$remoción & \cref{sec:flujos1d} & Estructura de las 4 ODEs \eqref{eq:eritro_H}--\eqref{eq:eritro_E} \\
Línea de fase, punto fijo estable & \cref{sec:puntos_fijos} & Equilibrio sano de \eqref{eq:eritro_1d} \\
Criterio $g'(x^*)$, tiempo $\tau=1/|g'|$ & \cref{sec:estabilidad_lineal} & Estabilidad y relajación de $C^*$; vida media $1/\gamma_C$ \\
Eliminación adiabática / escalas & \cref{sec:estabilidad_lineal} & Esclavizar $E,R,H$ frente a $C$ \\
Forma normal transcrítica $rx-x^2$ & \cref{sec:transcritica} & Establecimiento vs.\ aplasia del progenitor \eqref{eq:eritro_transcritica} \\
Diagrama de bifurcación & \cref{sec:saddlenode} & $H^*(\rho)$ y deslizamiento suave $C^*(\gamma_C)$ \\
Nullclines & \cref{sec:nullclines} & Isoclinas del sistema 2D \eqref{eq:eritro_2d} \\
Jacobiano, traza--determinante & \cref{sec:lineales2d} & Clasificación nodo/foco de $(C^*,E^*)$ \\
Poincaré--Bendixson & \cref{sec:conservativos} & Descartar caos; sólo equilibrio o ciclo en 2D \\
Bifurcación de Hopf y retardo & \cref{sec:hopf} & Origen de las oscilaciones hematológicas \eqref{eq:eritro_retardo} \\
\hline
\end{tabular}
\end{table}

\subsection{Ejercicios}

\paragraph{Ejercicio 1 (mecánico).} En condiciones de EPO saturada, la diferenciación satura, $d(E)\approx d_{\max}$, y la producción de \eqref{eq:eritro_1d} se vuelve una constante $P_0 = H^* d_{\max}$. Resuelve el equilibrio de $\dot C = P_0 - \gamma_C C$, clasifícalo por linealización y calcula el tiempo de relajación. Con $\gamma_C = 10^{-7}\,\mathrm{s}^{-1}$, exprésalo en días.
\aside{Pista: es una EDO lineal de primer orden; el equilibrio es $C^*=P_0/\gamma_C$ y $g'(C^*)=-\gamma_C$. Recuerda el \cref{sec:estabilidad_lineal}: $\tau=1/|g'|$.}
% SOL: C* = P0/gamma_C, estable pues g'(C*)=-gamma_C<0; tau = 1/gamma_C = 1e7 s = 116 dias.

\paragraph{Ejercicio 2 (análisis).} Anemia hemolítica: el paper la modela con $\gamma_C \to 4.8\,\gamma_C$ (vida media de RBC de $120$ a $25$ días). Con producción constante $P_0$ (Ejercicio 1), demuestra que el nuevo equilibrio es $C^*_{\text{nuevo}} = C^*_{\text{sano}}/4.8$ y que el tiempo de relajación se acorta en el mismo factor. Interpreta biológicamente: ¿por qué el paciente alcanza su (menor) equilibrio \emph{más rápido}?
\aside{Pista: tanto $C^*=P_0/\gamma_C$ como $\tau=1/\gamma_C$ son inversamente proporcionales a $\gamma_C$. Un flujo de muerte más veloz reequilibra antes.}
% SOL: C* y tau ambos escalan como 1/gamma_C -> se dividen por 4.8. La vida media corta acelera el recambio, luego el nuevo estado se alcanza en ~25 dias en vez de ~116.

\paragraph{Ejercicio 3 (análisis; bifurcación).} Toma la ecuación del progenitor saturada \eqref{eq:eritro_transcritica}, $\dot H = (a-d)H - (a/H_{\max})H^2$. (i) Halla sus dos puntos fijos. (ii) Linealiza y determina la estabilidad de cada uno en función del signo de $a-d$. (iii) Identifica el valor del parámetro $\rho=d/a$ en que ocurre la bifurcación y su tipo, comparando con la forma normal del \cref{sec:transcritica}. (iv) Interpreta el caso $\rho>1$.
\aside{Pista: es idéntica a la logística/transcrítica con $r=a-d$. En $H=0$, $g'=a-d$; en $H^*=H_{\max}(1-d/a)$, $g'=-(a-d)$.}
% SOL: FP en H=0 y H*=Hmax(1-d/a). Para a>d (rho<1): H=0 inestable, H* estable. Bifurcacion transcritica en rho=1 (a=d): intercambio de estabilidad. rho>1: solo H=0 estable = aplasia total (medula no se autosostiene).

\paragraph{Ejercicio 4 (análisis; 2D).} Para el sistema 2D \eqref{eq:eritro_2d}, ya calculamos $\operatorname{tr}J = -\gamma_C-\gamma_E H^*$ y $\det J>0$. (i) Argumenta con el diagrama traza--determinante del \cref{sec:lineales2d} por qué el equilibrio no puede ser nunca silla ni fuente. (ii) Escribe la condición $(\operatorname{tr}J)^2 = 4\det J$ que separa nodo de foco. (iii) ¿Qué observación clínica tras una donación de sangre distinguiría, en principio, un régimen de foco (sobreimpulso) de uno de nodo (recuperación monótona) en la trayectoria de hemoglobina?
\aside{Pista: silla exige $\det J<0$; fuente exige $\operatorname{tr}J>0$. Ambos están excluidos. Un foco produce una hemoglobina que sobrepasa su valor basal antes de asentarse; un nodo se aproxima sin rebasar.}
% SOL: (i) det>0 excluye silla; tr<0 excluye fuente -> siempre nodo o foco estable. (ii) frontera nodo/foco: (gamma_C+gamma_E H*)^2 = 4 det J. (iii) Foco: la Hb rebasa el basal (overshoot amortiguado) antes de estabilizar; nodo: monotona.

\paragraph{Ejercicio 5 (interpretación / paciente).} La anemia aplásica se modela con $H_{\max}\to H_{\max}/7$, lo que baja $H^*$. Usando la relación de equilibrio de la EPO \eqref{eq:eritro_EC}, $E^* = (\sigma_E/\gamma_E H^*)\,e^{-C^*/D}$, argumenta cualitativamente por qué la EPO del paciente aplásico está \emph{paradójicamente elevada} pese a que su producción de glóbulos rojos es baja. ¿Qué papel juega el que la EPO se aclare por unión a $H$?
\aside{Pista: $E^*\propto 1/H^*$. Si hay menos progenitores para captar EPO, la EPO se acumula. El aclaramiento dependiente de $H$ es la clave; con aclaramiento constante $\gamma_E E$ el efecto desaparecería.}
% SOL: Menor H* -> menor consumo de EPO (aclaramiento gamma_E H* E) y ademas menor C* -> mayor e^{-C*/D}; ambos suben E*. La EPO elevada refleja masa eritroide reducida, no fallo renal. Coincide con Dor-Alon: EPO como indicador de actividad medular.

\paragraph{Ejercicio 6 (numérico, Python).} Integra la reducción 2D \eqref{eq:eritro_2d} (parámetros adimensionalizados) con Euler explícito, partiendo del equilibrio, e introduce una hemorragia aguda que baje $C$ un $30\%$ en $t=t_0$. Grafica $C(t)$ y comprueba si la recuperación es monótona o con sobreimpulso.
\begin{minted}[fontsize=\small, framesep=4pt, frame=lines]{python}
import numpy as np

# parametros adimensionales (Hstar absorbido; escalas tipicas)
gC, gE, Hs, sE, D, dmax, Kd = 0.30, 0.45, 1.0, 3.6, 2.0, 1.0, 1.0

def f(C, E):
    dC = Hs*dmax*E/(Kd+E) - gC*C
    dE = sE*np.exp(-C/D)   - gE*Hs*E
    return dC, dE

# equilibrio aproximado por relajacion libre
C, E = 1.7, 1.6
dt = 0.01
for _ in range(20000):
    dC, dE = f(C, E)
    C, E = C + dt*dC, E + dt*dE
Ceq = C
C = 0.7*Ceq            # hemorragia: -30 % en t0
traj = []
for k in range(60000):
    dC, dE = f(C, E)
    C, E = C + dt*dC, E + dt*dE
    traj.append(C)
traj = np.array(traj)
print(f"C_eq={Ceq:.3f}  min={traj.min():.3f}  "
      f"max={traj.max():.3f}  overshoot={traj.max()>Ceq*1.001}")
\end{minted}
\coderef{eritropoyesis.py}{https://github.com/uncrayon/notas-biomate/blob/main/code/ch20-eritropoyesis/eritropoyesis.py}
\aside{Pista: si \texttt{max} supera a \texttt{C\_eq} tras el mínimo, hubo sobreimpulso (foco); si $C(t)$ sube monótonamente hacia \texttt{C\_eq}, fue un nodo. Prueba a variar \texttt{gE} para cruzar la frontera $(\operatorname{tr}J)^2=4\det J$ del Ejercicio 4.}
% SOL: Con estos valores tr=-(gC+gE*Hs)=-0.75, det>0; discriminante decide. Al bajar gE el sistema pasa de nodo a foco y aparece overshoot en C(t).


% ---- Demostración de estilo (mantiene el contenido existente) ----
% % ============================================================
% Contenido de demostración (idéntico en las tres variantes)
% Modelo Lotka--Volterra: matemáticas, figura, código, citas.
% ============================================================

\section{Modelos depredador--presa}

El modelo de Lotka--Volterra describe la interacción entre una población de
presas $x(t)$ y una de depredadores $y(t)$~\cite{volterra1926,murray2002}:
\begin{align}
  \frac{dx}{dt} &= \alpha\, x - \beta\, x y, \label{eq:lv-x}\\[2pt]
  \frac{dy}{dt} &= -\gamma\, y + \delta\, x y, \label{eq:lv-y}
\end{align}
donde $\alpha$ es la tasa de crecimiento de las presas, $\gamma$ la tasa de
mortandad de los depredadores, y $\beta$, $\delta$ cuantifican el encuentro
entre ambas especies.

\begin{definition}[Equilibrio de coexistencia]
El sistema~\eqref{eq:lv-x}--\eqref{eq:lv-y} admite, además del
equilibrio trivial $(0,0)$, el punto de coexistencia
$(x^{*}, y^{*}) = \left(\tfrac{\gamma}{\delta},\, \tfrac{\alpha}{\beta}\right)$.
\end{definition}

\begin{theorem}[Estabilidad del equilibrio de coexistencia]
\label{thm:lv}
Para $\alpha, \beta, \gamma, \delta > 0$, la matriz jacobiana del sistema
evaluada en $(x^{*}, y^{*})$ tiene autovalores puramente imaginarios
$\lambda_{\pm} = \pm i \sqrt{\alpha \gamma}$; el equilibrio es un centro
(linealmente neutro) y las trayectorias cercanas son órbitas cerradas.
\end{theorem}

\begin{proof}
El jacobiano es
$J(x,y) = \begin{pmatrix} \alpha - \beta y & -\beta x \\ \delta y & -\gamma + \delta x \end{pmatrix}$,
de modo que $J(x^{*}, y^{*}) = \begin{pmatrix} 0 & -\beta\gamma/\delta \\ \delta\alpha/\beta & 0 \end{pmatrix}$,
con $\operatorname{tr} J = 0$ y $\det J = \alpha\gamma > 0$.
La \cref{fig:phase} muestra el retrato de fases para
$\alpha = \beta = \gamma = \delta = 1$.
\end{proof}

\begin{figure}[htbp]
  \centering
  \begin{tikzpicture}
    \begin{axis}[
        view={0}{90}, domain=0:2, y domain=0:2,
        xmin=0, xmax=2, ymin=0, ymax=2,
        samples=13, width=0.72\linewidth,
        xlabel={$x$ (presas)}, ylabel={$y$ (depredadores)},
        axis lines=middle, grid=both, grid style={gray!15},
        xlabel near ticks, ylabel near ticks,
      ]
      \addplot3[gray!70, -latex, quiver={
          u={x - x*y}, v={-y + x*y}, scale arrows=0.055}] (x,y,0);
      \draw[red!70!black, thick, dashed] (axis cs:1,0) -- (axis cs:1,2)
        node[below right, font=\footnotesize]{$x^{*}=1$};
      \draw[blue!70!black, thick, dashed] (axis cs:0,1) -- (axis cs:2,1)
        node[above right, font=\footnotesize]{$y^{*}=1$};
      \fill (axis cs:1,1) circle (2.2pt) node[above right, font=\footnotesize]{$(x^{*},y^{*})$};
    \end{axis}
  \end{tikzpicture}
  \caption{Retrato de fases del modelo de Lotka--Volterra con
  $\alpha=\beta=\gamma=\delta=1$. Las nulas roja y azul son los
  \emph{nullclines}; el equilibrio $(1,1)$ es un centro.}
  \label{fig:phase}
\end{figure}

\aside{Los autovalores imaginarios puros no prueban estabilidad: el
sistema conservativo tiene un primer integral $V(x,y)=\delta x - \gamma
\ln x + \beta y - \alpha \ln y$, y sus curvas de nivel son las órbitas
cerradas~\cite{strogatz2015}.}

\subsection{Integración numérica en Python}

La \cref{eq:lv-x}--\eqref{eq:lv-y} se integra con Euler explícito:

\begin{minted}[fontsize=\small, framesep=4pt, frame=lines]{python}
import numpy as np

alpha, beta, gamma, delta = 1.0, 1.0, 1.0, 1.0

def lv(u, t):
    """Campo vectorial de Lotka-Volterra."""
    x, y = u
    return np.array([alpha*x - beta*x*y,
                     -gamma*y + delta*x*y])

dt, T = 0.01, 20.0        # paso temporal y horizonte
u = np.array([1.2, 0.6])  # condición inicial
for t in np.arange(0, T, dt):
    u = u + dt * lv(u, t)
print(f"estado final: x = {u[0]:.3f}, y = {u[1]:.3f}")
\end{minted}

\subsection{Simulación con agentes en NetLogo}

La misma dinámica puede explorarse como un modelo basado en
agentes~\cite{netlogo}, donde cada presa y depredador es una tortuga:

\begin{minted}[fontsize=\small, framesep=4pt, frame=lines]{lexer/netlogo.py:NetLogoLexer}
globals [presas-iniciales depredadores-iniciales]
turtles-own [energia]

to setup
  clear-all
  create-turtles presas-iniciales [
    setxy random-xcor random-ycor
    set color green  set energia 10
  ]
  create-turtles depredadores-iniciales [
    setxy random-xcor random-ycor
    set color red
  ]
  reset-ticks
end

to go
  ask turtles [
    ifelse color = red
      [ let presa min-one-of turtles with [color = green]
                        [distance-to self]
        if presa != nobody [ face presa  fd 1.2 ] ]
      [ fd 1  set energia energia - 1
        if energia < 2 [ die ] ]
  ]
  tick
end
\end{minted}

\subsection{Dinámica discreta: del mapa logístico al caos}

Al discretizar el crecimiento logístico con tiempo
discreto~\cite{may1976} aparece el mapa
\begin{equation}
  N_{t+1} = r\, N_t \left(1 - \frac{N_t}{K}\right),
  \label{eq:logistic-map}
\end{equation}
que, bajo el cambio $z_t = N_t/K$, se reduce al mapa logístico
$z_{t+1} = \mu\, z_t (1 - z_t)$ con $\mu = 1 + r$. Para $\mu > 3{,}57$
la dinámica es caótica: un modelo determinista de una sola variable
produce trayectorias impredecibles, un fenómeno que May argumentó que
debe tomarse en serio en la gestión de poblaciones reales.

% ---------------------------------------------------------
\nocite{edelstein2007,tongqft}
\printbibliography[heading=bibintoc, title={Referencias}]
   % Lotka--Volterra: mover a Parte IV (sistemas 2D) cuando redacte esa parte

\printbibliography[heading=bibintoc, title={Referencias}]
\end{document}
